\chapter{The geometric cobar complex}
\label{chap:cobar-construction}

% Local macros (provide-command idempotent; needed for cross-chapter
% references to thm:koszul-reflection and chap:climax-platonic).
\providecommand{\Bbarch}{\bar{B}^{\mathrm{ch}}}
\providecommand{\Omegach}{\Omega^{\mathrm{ch}}}
\providecommand{\barBch}{\bar{B}^{\mathrm{ch}}}
\providecommand{\Kosz}{\mathsf{Kosz}}
\providecommand{\co}{\mathrm{co}}

The bar functor destroys the multiplication: it encodes~$\cA$ as the
symmetric factorization coalgebra $\barB(\cA)$, where the collision
data survive but the product one must recover is no longer visible.
The geometric cobar complex
$\Omega^{\mathrm{ch}}(\cC)=\bigoplus_{n\geq 0}\Omega^{\mathrm{ch}}_n(\cC)$
is the reconstruction machine for this loss on a conilpotent chiral
coalgebra~$\cC$, meaning a factorization coalgebra whose iterated
reduced coproduct vanishes on every local section. It realizes the
inversion $\Omega(\barB(\cA)) \xrightarrow{\sim} \cA$ on the
Koszul/completed locus (Theorem~B). The universal twisting morphism
is the datum behind the adjunction counit; the modular MC element
$\Theta_\cA \in \MC(\gAmod)$ records the compatible obstruction tower,
not a separate Koszul-dual object. Definitions~\ref{def:factorization-coalgebra},
\ref{def:conilpotent-cobar}, and
\ref{def:geom-cobar-intrinsic} give the formal versions used here.
Bar and cobar are represented by the same twisting-morphism calculus
$\tau \in \MC\bigl(\operatorname{Conv}(\cC, \cA)\bigr)$
(Definition~\ref{def:twisting-morphism}): bar records collision
residues via the logarithmic propagator
$\eta_{ij} = d\log(z_i - z_j)$; cobar uses the Verdier-transposed
boundary currents. The exchange between residues and currents,
$j_*$ and $j_!$, is the Verdier exchange on
$\operatorname{Conv}(\cC, \cA)$.

Concretely, given a conilpotent chiral
coalgebra~$\cC$ on~$X$ valued in holonomic $\cD$-modules, the
\emph{geometric cobar complex} is
$\Omega^{\mathrm{ch}}_n(\cC)
:= \mathbb{D}_{\overline{C}_{n+1}(X)}
\bigl(j_*\, j^*(\cC^{\vee,\boxtimes(n+1)})
\otimes \Omega^n(\log D)\bigr)$
(Definition~\ref{def:geom-cobar-intrinsic}), where
$\cC^\vee = \mathbb{D}_X(\cC)$ is the $\cD$-module Verdier dual
and $j \colon C_{n+1}(X) \hookrightarrow \overline{C}_{n+1}(X)$
is the open inclusion. In the uncurved conilpotent lane the cobar
differential is the Verdier transport of the bar differential through
the Hom functor:
$d_{\mathrm{cobar}} := \mathbb{D}(d_{\mathrm{bar}})$.
When the bar object is curved, the same Verdier transport produces a
CDG cobar object; its square is the inner curvature operator, not zero
on the nose.
The exchange $\mathbb{D} \circ j_* \cong j_! \circ \mathbb{D}$
(Lemma~\ref{lem:verdier-extension-exchange}) identifies the
cobar kernel as the distributional partner of the logarithmic
propagator: the Schwartz kernel
$K \in \cD'(C_n(X) \times C_m(X))$ of
Theorem~\ref{thm:schwartz-kernel-cobar} is the $j_!$-side of the
same MC element~$\tau$ whose $j_*$-side is the bar propagator.
This is the geometric content of the representability theorem
$\operatorname{Tw}(\cC, \cA) = \MC(\operatorname{Conv}(\cC, \cA))$
(Theorem~\ref{thm:bar-cobar-adjunction}).

At genus~$g \geq 1$ the fiberwise differential satisfies
$d_{\mathrm{fib}}^2 = \kappa(\cA) \cdot \omega_g$, reflecting the curved
$A_\infty$ structure $\mu_1^2 = [\mu_0, -]$
(Chapter~\ref{chap:bar-cobar-adjunction}). Thus the raw fiberwise
operator is not a differential unless the curvature acts trivially.
Strict nilpotence belongs to the uncurved genus-$0$ complex, to the
flat total comparison differential, to a Maurer--Cartan twist killing
the curvature, or to the coderived totalization. The cobar functor
applied to the bar coalgebra viewed as an FCom-algebra
(Theorem~\ref{thm:bar-modular-operad}) encodes the shadow
obstruction tower $\Theta_\cA^{\leq r}$
(Proposition~\ref{prop:cobar-modular-shadow}).
Strict inversion may require completion or coderived contexts
at higher genus.

The chapter sits inside the unified Koszul reflection of
Theorem~\ref{thm:koszul-reflection}: bar and cobar are the two
coordinates $K = \barBch$ and $K^{-1} = \Omegach$ of one involutive
adjoint equivalence $K^{2}\simeq\mathrm{id}$, with the cobar
construction the explicit chain-level expression of $K^{-1}$. The
cobar-bar inversion proved here in
Theorem~\ref{thm:bar-cobar-inversion-qi} is the chain-level
realisation of clause~\ref{KR-ii} on the Koszul locus, refined by
distributional kernels to a complete and constructive form. Off the
Koszul locus, clause~\ref{KR-iii} replaces chain-level inversion by
coderived involutivity in
$D^{\co}(\barBch\text{-}\mathrm{CoFact})$, and the
weight-completed extension to class $\mathsf{M}$ is given by
Chapter~\ref{chap:mc5-class-m-chain-level-platonic}.

\subsection{Motivation: reversing the prism}

\begin{remark}[Three functors on the bar coalgebra]
\label{rem:cobar-three-functors}
Before entering the construction, we fix a critical distinction that
pervades the chapter. Three functors act on the bar coalgebra
$\barB(\cA)$, producing three different objects; confusing any two
produces false statements.
\begin{enumerate}[label=\textup{(\roman*)}]
\item \emph{Cobar / bar-cobar inversion}
 (Theorem~B, Corollary~\ref{cor:bar-cobar-inverse}):
 \[\Omegach(\barB(\cA)) \;\xrightarrow{\sim}\; \cA\qquad
 \text{(recovers $\cA$ itself: bar-cobar inversion)}.\]
\item \emph{Verdier duality on $\operatorname{Ran}(X)$}
 (Theorem~A, Convention~\ref{conv:bar-coalgebra-identity}):
 \[\mathbb{D}_{\operatorname{Ran}}(\barB(\cA)) \;\simeq\;
 (\cA)^!_\infty\qquad
 \text{(gives the homotopy Koszul dual factorization algebra)}.\]
\item \emph{Chiral derived centre / chiral Hochschild cochains}
 (Theorem~H):
 \[Z^{\mathrm{der}}_{\mathrm{ch}}(\cA)
 \;=\; C^\bullet_{\mathrm{ch}}(\cA, \cA)\qquad
 \text{(the cochain-level closed-sector observables)}.\]
\end{enumerate}
\emph{Cobar is the inverse, not the dual.} Throughout this
chapter $\Omegach$ denotes the bar-cobar inverse functor of
(i); it recovers the algebra from its bar coalgebra. The
homotopy Koszul dual factorization algebra $\cA^!_\infty$ is obtained by
Verdier duality on $\operatorname{Ran}(X)$, and on the Koszul locus
Theorem~A identifies $\cA^!_\infty \simeq \cA^!$. This is a
\emph{different} operation on $\barB(\cA)$. A statement such as
``the cobar of the bar gives the Koszul dual''
is false; the correct statement is
``the cobar of the bar gives $\cA$, and Verdier duality of the bar
gives the homotopy Koszul dual factorization algebra.''
\end{remark}

\begin{remark}[Which bar complex does the cobar invert?]
\label{rem:cobar-which-bar}
The manuscript distinguishes three bar complexes on a chiral
algebra~$\cA$:
\begin{enumerate}[label=\textup{(\roman*)}]
\item $\barB^{\mathrm{FG}}(\cA) = \mathrm{Lie}^c(s^{-1}\bar{\cA})$, the
 \emph{Francis--Gaitsgory} bar using the cofree coLie
 coalgebra on the augmentation coideal (the Harrison complex);
\item $\barB^{\Sigma}(\cA) = \mathrm{Sym}^c(s^{-1}\bar{\cA})$, the
 \emph{symmetric / factorization} bar whose underlying
 coalgebra is cofree cocommutative coassociative with the
 coshuffle coproduct ($2^n$ terms; the factorization coalgebra of
 Definition~\ref{def:factorization-coalgebra});
\item $\barB^{\mathrm{ord}}(\cA) = T^c(s^{-1}\bar{\cA})$, the
 \emph{ordered} bar using the cofree coassociative
 (non-cocommutative) coalgebra with the deconcatenation
 coproduct ($n+1$ terms).
\end{enumerate}
The cobar functor of this chapter is the left adjoint to
$\barB^{\Sigma}$ (the factorization/symmetric bar): given a
factorization coalgebra~$\cC$ on~$X$, $\Omegach(\cC)$ is the free
chiral algebra on the desuspension $s^{-1}\bar{\cC}$ with
differential extending the reduced comultiplication
(Theorem~\ref{thm:cobar-free}). Bar-cobar inversion
(Theorem~B, Corollary~\ref{cor:bar-cobar-inverse}) is stated for
this symmetric/factorization bar: $\Omegach(\barB^{\Sigma}(\cA))
\xrightarrow{\sim} \cA$. The ordered bar $\barB^{\mathrm{ord}}(\cA)$
carries more information (it tracks point order) and its cobar is a
related but distinct construction; we work exclusively with the
symmetric cobar except where explicitly noted.
\end{remark}

\begin{remark}[Inverse prism principle]
The bar construction extracts residues at collision divisors; the cobar integrates over configuration spaces. Their duality is the residue-integration pairing on logarithmic forms.

\emph{QFT comparison.}
A perturbative field theory may supply an analytic kernel model in
which logarithmic bar forms become compactified configuration-space
integrands and FM cobar currents become source terms for Green
kernels. This is extra comparison data, not a classification theorem:
the bar coalgebra, the cobar algebra, the Koszul dual, and QFT
propagator distributions remain distinct objects.

\emph{Residue-current dictionary.}
\begin{center}
\begin{tabular}{c|c|c}
& \textbf{Bar} & \textbf{Cobar} \\ \hline
Space & Compactified $\overline{C}_n(X)$ & Open $C_n(X)$ \\
Forms & Logarithmic (residues) & Distributional (delta functions) \\
Operation & Extract (analyze) & Insert (synthesize) \\
Boundary & Normal crossing divisors & Diagonal singularities \\
QFT comparison & Compactified configuration integrands & Green-kernel source currents \\
\end{tabular}
\end{center}
\end{remark}

\subsection{Distribution theory prerequisites}

\begin{definition}[Test function space]\label{def:test-functions}
For the open configuration space $C_n(X)$, define the test function space:
\[\mathcal{D}(C_n(X)) = C_c^\infty(C_n(X), \mathbb{C})\]
consisting of smooth, compactly supported functions. This is equipped with the
inductive limit topology from exhaustion by compact sets.
\end{definition}

\begin{definition}[Distribution space]\label{def:distributions}
\index{distribution!on configuration space}
The space $\mathcal{D}'(C_n(X))$ of \emph{distributions} on $C_n(X)$ is the
continuous dual:
\[\mathcal{D}'(C_n(X)) = \mathcal{D}(C_n(X))^*\]
equipped with the weak-$*$ topology. A distribution $T \in \mathcal{D}'(C_n(X))$
is a continuous linear functional:
\[\langle T, \phi \rangle \in \mathbb{C} \quad \text{for all } \phi \in \mathcal{D}(C_n(X))\]
\end{definition}

\begin{example}[Fundamental distributions]\label{ex:fundamental-distributions}
\emph{Dirac delta.} For $p \in C_n(X)$:
\[\langle \delta_p, \phi \rangle = \phi(p)\]

\emph{Principal value.} On a local chart of $X^n$ or
$\overline{C}_n(X)$ with collision divisor $D_{ij}$ lying over
$z_i=z_j$:
\[\langle \text{PV}\left(\frac{1}{z_i - z_j}\right), \phi \rangle =
\lim_{\epsilon \to 0} \int_{|z_i - z_j| > \epsilon} \frac{\phi(z_1, \ldots, z_n)}{z_i - z_j}
dz_1 \cdots dz_n\]

\emph{Hadamard finite part.} For higher-order poles:
\[\text{FP}\left(\frac{1}{(z_i - z_j)^k}\right) =
\lim_{\epsilon \to 0} \left[\int_{|z_i - z_j| > \epsilon} \frac{\phi}{(z_i - z_j)^k} -
\frac{\text{(divergent terms)}}{\epsilon^{k-1}}\right]\]
\end{example}

\begin{theorem}[Schwartz kernel theorem for cobar {\cite{Hormander}}; \ClaimStatusProvedElsewhere]\label{thm:schwartz-kernel-cobar}
\index{Schwartz kernel!cobar}
Every continuous linear operator:
\[K: \mathcal{D}(C_n(X)) \to \mathcal{D}'(C_m(X))\]
is represented by a distribution kernel:
\[K \in \mathcal{D}'(C_n(X) \times C_m(X))\]
such that:
\[(K\phi)(z_1, \ldots, z_m) = \int_{C_n(X)} K(z_1, \ldots, z_m; w_1, \ldots, w_n) \phi(w_1, \ldots, w_n)\]
\end{theorem}

\begin{proof}
This is a special case of the Schwartz kernel theorem for distributions on smooth manifolds (H\"ormander~\cite{Hormander}, Theorem~5.2.1). The configuration spaces $C_n(X)$ and $C_m(X)$ are smooth manifolds, and continuity of $K$ in the distributional topology guarantees the existence of the integral kernel $K \in \mathcal{D}'(C_n(X) \times C_m(X))$.
\end{proof}

\subsection{Intrinsic cobar via Verdier duality}\label{subsec:intrinsic-cobar}

\begin{definition}[Conilpotent and completed coalgebras]\label{def:conilpotent-cobar}
\index{conilpotent!coalgebra}
\index{completed coalgebra!cobar}
Let $\cC$ be a coaugmented coalgebra with coaugmentation
$\eta\colon\omega_X\to\cC$, counit
$\varepsilon\colon\cC\to\omega_X$, coaugmentation coideal
$\bar\cC:=\operatorname{coker}(\eta)\cong\ker(\varepsilon)$,
and reduced comultiplication
$\bar\Delta\colon\bar\cC\to\bar\cC\otimes\bar\cC$.
It is \emph{conilpotent} if for every local section
$c\in\bar\cC$ there is $N=N(c)$ such that
\[
\bar\Delta^{(N)}(c)=0 .
\]
Equivalently,
\[
\bar\cC=\bigcup_{N\geq 0}F_N\bar\cC,\qquad
F_N\bar\cC:=\ker(\bar\Delta^{(N)}),
\]
the increasing coradical filtration exhausts the reduced coalgebra.
Then the cobar differential on the generator $s^{-1}c$ is a finite
sum, and its iterates terminate on each element of the free chiral
algebra.

The completed replacement is a descending coaugmentation-adic
filtration
\[
\cC=F^0\cC\supset F^1\cC=\bar\cC\supset F^2\cC\supset\cdots,
\qquad
\cC\simeq\varprojlim_N \cC/F^{N+1}\cC,
\]
with separated, exhaustive, finite-type quotients and continuous
coproduct. In that lane $\Omega^{\mathrm{ch}}(\cC)$ means the
continuous completed cobar object. It is an ordinary dg chiral algebra
only in the square-zero uncurved or flat-total model; with curvature it
is a CDG object satisfying
$d_{\Omega,h}^{\,2}=[h_\Omega,-]$.
See Definition~\ref{def:conilpotent-complete} and
Definition~\ref{def:strong-completion-tower} for the global filtered
form used in the inversion theorem.
\end{definition}

\begin{definition}[Factorization coalgebra]\label{def:factorization-coalgebra}
\index{factorization coalgebra|textbf}
A \emph{factorization coalgebra} on a smooth
curve~$X$ is a $\cD_X$-module~$\cC$ equipped with:
\begin{enumerate}[label=\textup{(\roman*)}]
\item a comultiplication
 $\Delta \colon \cC \to
 j_* j^*(\cC \boxtimes \cC)$, where
 $j \colon X^2 \setminus \Delta_X \hookrightarrow X^2$;
\item a counit
 $\varepsilon \colon \cC \to \omega_X$;
\item the \emph{factorization} (colocality) property:
 $\Delta$ factors through $j_* j^*$, so the
 comultiplication is comeromorphic with copoles only
 along the diagonal;
\item coassociativity: the Borcherds co-identity
 (the dual of
 equation~\eqref{eq:borcherds-axiom}) holds for the
 iterated comultiplication on codimension-$2$ strata.
\end{enumerate}
Under the finite-type holonomic hypotheses used in this chapter,
Verdier duality sends chiral algebras to factorization coalgebras.
If $\cA$ is an augmented chiral algebra, then
$\cC=\barB(\cA)$ is such a factorization coalgebra. Conversely, a
coaugmented conilpotent factorization coalgebra has a cobar
dg chiral algebra $\Omega(\cC)$; it recovers a prescribed algebra only
when $\cC$ lies in the bar image and the Koszul/completed hypotheses of
Corollary~\ref{cor:bar-cobar-inverse} hold.
\end{definition}

\begin{definition}[Intrinsic geometric cobar complex]\label{def:geom-cobar-intrinsic}
\index{cobar construction!geometric|textbf}
Let $\mathcal{C}$ be a conilpotent chiral coalgebra on a smooth algebraic
curve $X$, valued in regular holonomic finite-type $\mathcal{D}$-modules,
with comultiplication
$\Delta\colon \mathcal{C} \to \mathcal{C} \boxtimes \mathcal{C}$ and
coaugmentation $\eta\colon \omega_X \to \mathcal{C}$. Write
$\mathcal{C}^\vee := \mathbb{D}_X(\mathcal{C})$ for the $\mathcal{D}$-module
Verdier dual on $X$, which is a chiral algebra with multiplication dual to
$\Delta$ whenever the relevant tensor products are finite-type; in the
completed case all duals, tensor products, and Homs below are interpreted
continuously.

The \emph{geometric cobar complex} is the graded object:
\[
\Omega^{\mathrm{ch}}_n(\mathcal{C})
:= \mathbb{D}_{\overline{C}_{n+1}(X)}\!\left(
 j_*\, j^*\bigl((\mathcal{C}^\vee)^{\boxtimes(n+1)}\bigr)
 \otimes \Omega^n_{\overline{C}_{n+1}(X)}(\log D)
\right)
\]
where $j\colon C_{n+1}(X) \hookrightarrow \overline{C}_{n+1}(X)$ is the
inclusion, $D = \partial\overline{C}_{n+1}(X)$ is the boundary divisor, and
$\mathbb{D}_{\overline{C}_{n+1}(X)}$ is the Verdier duality functor on the
compactified configuration space.

In the uncurved case the \emph{cobar differential} is:
\[
d_{\mathrm{cobar}} := \mathbb{D}(d_{\mathrm{bar}})
\]
the Verdier dual of the bar differential
(Definition~\ref{def:bar-differential-complete}). In the curved CDG
case this formula denotes the degree-$1$ derivation underlying the
curved cobar object; its square is governed by
Corollary~\ref{cor:cobar-nilpotence-verdier}, not by strict
nilpotence.
\end{definition}

\begin{remark}\label{rem:cobar-intrinsic-meaning}
The definition says: \emph{to build the cobar of a coalgebra $\mathcal{C}$,
first form the bar of its dual algebra $\mathcal{C}^\vee$, then apply
Verdier duality on the configuration space.} Concretely, the bar uses
$j_*$ (maximal extension across diagonals) and logarithmic forms; the cobar
uses $j_!$ (extension by zero) and distributional sections. This is
precisely the exchange effected by Verdier duality:
$\mathbb{D} \circ j_* \cong j_! \circ \mathbb{D}$
(Lemma~\ref{lem:verdier-extension-exchange}).
\end{remark}

\begin{lemma}[Holonomicity of the bar complex; \ClaimStatusProvedHere]%
\label{lem:bar-holonomicity}
\index{bar complex!holonomicity}
Let $\cA$ be a chiral algebra on a smooth projective curve~$X$,
finitely generated over~$\mathcal{D}_X$ with regular singularities.
Then for each $n \geq 0$:
\begin{enumerate}[label=\textup{(\alph*)}]
\item The $n$-th bar component $\bar{B}^{\mathrm{ch}}_n(\cA)$ is a
 regular holonomic $\mathcal{D}$-module on the Fulton--MacPherson
 compactification~$\overline{C}_n(X)$.
\item The bar differential
 $d_{\mathrm{bar}} \colon \bar{B}^{\mathrm{ch}}_n(\cA) \to
 \bar{B}^{\mathrm{ch}}_{n-1}(\cA)$
 is a morphism of holonomic $\mathcal{D}$-modules.
\end{enumerate}
\end{lemma}

\begin{proof}
\emph{Step~1: External tensor product.}
Since $\cA$ is a holonomic $\mathcal{D}_X$-module by hypothesis,
the external tensor product $\cA^{\boxtimes n}$ is a holonomic
$\mathcal{D}_{X^n}$-module
\cite[Theorem~3.2.3]{HTT08}.

\emph{Step~2: Pullback to the FM compactification.}
The blowdown morphism $\pi \colon \overline{C}_n(X) \to X^n$ is
proper and birational, so $\pi^!\cA^{\boxtimes n}$ is a holonomic
$\mathcal{D}_{\overline{C}_n(X)}$-module
\cite[Theorem~3.2.3]{HTT08}.

\emph{Step~3: Logarithmic extension.}
The bar sheaf on the open configuration space $C_n(X)$ is
$j^*\pi^!\cA^{\boxtimes n} \otimes \Omega^*_{\log}(D)$,
where $j \colon C_n(X) \hookrightarrow \overline{C}_n(X)$ is the
inclusion and $D = \overline{C}_n(X) \setminus C_n(X)$ is a normal
crossings divisor (Theorem~\ref{thm:FM}).
The logarithmic differential forms $\Omega^*_{\log}(D)$ define a
regular integrable connection on the NCD complement, so
$j_*$ of the resulting regular holonomic module on~$C_n(X)$ is
regular holonomic on~$\overline{C}_n(X)$ by Deligne's regularity
theorem \cite[Theorem~5.2.17]{HTT08}.
This identifies $\bar{B}^{\mathrm{ch}}_n(\cA)$ as a regular
holonomic $\mathcal{D}_{\overline{C}_n(X)}$-module.

\emph{Step~4: Bar differential.}
The bar differential $d_{\mathrm{bar}} = d_{\mathrm{internal}}
+ d_{\mathrm{residue}} + d_{\mathrm{form}}$
(Definition~\ref{def:bar-differential-complete}).
The internal differential $d_{\mathrm{internal}}$ is induced by the
chiral algebra structure, which is a $\mathcal{D}_X$-module morphism
by definition.
The form differential $d_{\mathrm{form}}$ is the de~Rham differential
on $\Omega^*_{\log}(D)$, which is a $\mathcal{D}$-module morphism.
The residue differential $d_{\mathrm{residue}} = \sum_D \operatorname{Res}_D$
is the composition of restriction to a boundary divisor~$D$ followed
by the trace map, both of which are morphisms of
$\mathcal{D}$-modules \cite[Chapter~4]{KS90}.
Hence $d_{\mathrm{bar}}$ is a morphism of holonomic
$\mathcal{D}$-modules.
\end{proof}

\begin{lemma}[Verdier duality exchanges extensions; \ClaimStatusProvedHere]%
\label{lem:verdier-extension-exchange}
Let $Y$ be a smooth variety, $U \xhookrightarrow{j} Y$ an open inclusion
with complement a normal crossings divisor, and $\mathcal{M}$ a holonomic
$\mathcal{D}_U$-module. Then there is a canonical isomorphism in the
derived category of holonomic $\mathcal{D}_Y$-modules:
\[
\mathbb{D}_Y(j_*\mathcal{M}) \;\cong\; j_!\,\mathbb{D}_U(\mathcal{M}).
\]
\end{lemma}

\begin{proof}
This is a standard consequence of the theory of holonomic
$\mathcal{D}$-modules; we recall the argument. By
adjunction, $j_*$ is right adjoint to $j^*$ and $j_!$ is left adjoint to
$j^!$. For $j$ an open inclusion we have $j^! = j^*$, so:
\[
j_! \dashv j^* \dashv j_*.
\]
Verdier duality exchanges left and right adjoints: for any holonomic
$\mathcal{M}$ on $U$,
\[
\mathbb{D}_Y(j_*\mathcal{M})
= \mathbb{D}_Y(j_* \mathbb{D}_U \mathbb{D}_U \mathcal{M})
\cong j_!(\mathbb{D}_U \mathcal{M})
\]
where the last step uses the identity
$\mathbb{D}_Y \circ j_* \circ \mathbb{D}_U \cong j_!$, which holds because
$j_*$ preserves holonomicity (the complement is a normal crossings divisor)
and Verdier duality is involutive on the holonomic category
\cite[Theorem~2.6.1]{HTT08}.
\end{proof}

\begin{theorem}[Distributional model of the cobar; \ClaimStatusProvedHere]%
\label{thm:cobar-distributional-model}
\index{cobar construction!distributional model}
Let $\mathcal{C}$ be a regular holonomic finite-type factorization
coalgebra that is either conilpotent
(Definition~\ref{def:conilpotent-complete}) or complete, separated, and
exhaustive for the coaugmentation-adic filtration, with finite-type
associated graded and continuous coproduct. In the completed case assume
the transition maps satisfy the Mittag--Leffler condition, so inverse
limits commute with the finite bar-degree constructions used below.
Since each component of the bar object of $\mathcal{C}^{\vee}$ is
holonomic by Lemma~\ref{lem:bar-holonomicity}, the intrinsic cobar
(Definition~\ref{def:geom-cobar-intrinsic}) admits a canonical
realization by Fulton--MacPherson boundary currents:
\[
\Omega^{\mathrm{ch}}_n(\mathcal{C})
\;\cong\;
\Gamma\!\left(C_{n+1}(X),\,
 \mathrm{Hom}_{\mathcal{D}}\!\left(\pi^*\mathcal{C}^{\otimes(n+1)},\,
 \mathcal{D}_{C_{n+1}(X)}\right)
 \otimes \Omega^*_{C_{n+1}(X),\mathrm{dist}}
\right)
\]
Here the right side is shorthand for the $j_!$-supported current model
on $\overline{C}_{n+1}(X)$ restricted to the open stratum. The symbols
$\delta(z_i-z_j)$ below denote boundary currents on the
Fulton--MacPherson divisor, not ordinary distributions on the open
manifold $C_{n+1}(X)$.

The three components of the cobar differential correspond to
Verdier duals of the bar differential components:
\begin{enumerate}[label=\textup{(\roman*)}]
\item $d_{\mathrm{extend}} = \mathbb{D}(d_{\mathrm{res}})$:
 Verdier-dual of the residue differential. Residues at boundary
 divisors (extracting OPE poles) dualize to delta-function insertions
 at diagonals (enforcing coincidence constraints).
\item $d_{\mathrm{comult}} = \mathbb{D}(d_{\mathrm{fact}})$:
 Verdier-dual of the factorization differential. Configuration-splitting
 in the bar dualizes to comultiplication in the cobar.
\item $d_{\mathrm{internal}} = \mathbb{D}(d_{\mathrm{int}})$:
 The internal differential is self-dual.
\end{enumerate}
\end{theorem}

\begin{proof}
\emph{Step 1: Identification of the underlying graded object.}
By Lemma~\ref{lem:verdier-extension-exchange},
\[
\mathbb{D}_{\overline{C}_{n+1}}\bigl(
 j_* j^* (\mathcal{C}^\vee)^{\boxtimes(n+1)}
 \otimes \Omega^n_{\log}
\bigr)
\;\cong\;
j_!\, j^*\, \mathbb{D}\bigl(
 (\mathcal{C}^\vee)^{\boxtimes(n+1)}
\bigr)
\otimes \mathbb{D}(\Omega^n_{\log}).
\]
Since $\mathbb{D}_X(\mathcal{C}^\vee) = \mathcal{C}$ (Verdier duality is
involutive on holonomic $\mathcal{D}$-modules) and
$\mathbb{D}(\Omega^n_{\log})$ is the sheaf of distributional $(d-n)$-currents
on $\overline{C}_{n+1}(X)$ (where $d = \dim_{\mathbb{R}} C_{n+1}(X)$), we obtain
distributional sections of $\mathcal{C}^{\boxtimes(n+1)}$ with
$j_!$-support; that is, sections on $C_{n+1}(X)$ extended by zero across the
boundary. Passing to global sections gives the distributional model.

\emph{Step 2: Differential components.}
The bar differential $d_{\mathrm{bar}} = d_{\mathrm{int}} + d_{\mathrm{res}}
+ d_{\mathrm{config}}$ is a morphism of holonomic $\mathcal{D}$-modules.
Since $\mathbb{D}$ is an exact functor on the holonomic category, each
component dualizes individually.

For (i): the residue differential $d_{\mathrm{res}}$ acts by
$\mathrm{Res}_{z_i = z_j}[\omega \cdot \eta_{ij}]$ (extracting the residue
along the boundary divisor $D_{ij}$). Under the Verdier pairing
$\langle \eta_{ij}, \delta(z_i - z_j)\rangle = 1$
(Theorem~\ref{thm:bar-cobar-verdier}, Step~3), the dual operation is
$K \mapsto \iota_{ij,!}\operatorname{Tr}_{D_{ij}}(K)$, written
mnemonically as $\delta(z_i-z_j)\otimes K|_{D_{ij}}$: inserting the
boundary current dual to the residue. This is exactly the extension
differential $d_{\mathrm{extend}}$.

For (ii): the factorization differential splits configurations by restricting
to strata where subsets of points collide. Its dual reassembles
configurations by applying the comultiplication $\Delta$ of $\mathcal{C}$.

For (iii): the internal differential $d_{\mathrm{int}}$ acts
coefficient-wise by the differential of $\mathcal{A}$ (resp.\ $\mathcal{C}$)
and does not interact with configuration space geometry; it is self-dual
under $\mathbb{D}_X(\mathcal{C}^\vee) = \mathcal{C}$.
\end{proof}

\begin{corollary}[Uncurved cobar nilpotence and curved square via Verdier duality;
\ClaimStatusProvedHere]\label{cor:cobar-nilpotence-verdier}
Let $\mathcal{C}$ be a conilpotent chiral coalgebra whose Verdier-dual
bar object is uncurved. Then the cobar differential satisfies
\[
d_{\mathrm{cobar}}^2 = 0.
\]
If $\mathcal{C}$ carries CDG curvature, the curved cobar is a triple
\[
\bigl(\Omega^{\mathrm{ch}}_{\mathrm{curv}}(\mathcal{C}),
  d_{\Omega,h}, h_\Omega\bigr)
\]
with $h_\Omega\in \Omega^{\mathrm{ch}}_{\mathrm{curv}}(\mathcal{C})^2$
and
\begin{equation}\label{eq:curved-cobar-square}
d_{\Omega,h}^{\,2}(x)
= [h_\Omega,x]
:= h_\Omega\cdot x - x\cdot h_\Omega,
\qquad
d_{\Omega,h}(h_\Omega)=0.
\end{equation}
For a Maurer--Cartan element $\alpha$ the twisted operator
$d_{\Omega,\alpha}=d_{\Omega,h}+[\alpha,-]$ satisfies
\begin{equation}\label{eq:curved-cobar-twist-square}
d_{\Omega,\alpha}^{\,2}
=
\left[
h_\Omega+d_{\Omega,h}\alpha+\frac{1}{2}[\alpha,\alpha],
-
\right].
\end{equation}
Thus a curved cobar object becomes strictly square-zero only after a
twist with
$h_\Omega+d_{\Omega,h}\alpha+\frac{1}{2}[\alpha,\alpha]=0$, or after
passing to the flat comparison/coderived totalization where curvature
is absorbed into the total structure.
\end{corollary}

\begin{proof}
In the uncurved lane the bar differential satisfies
$d_{\mathrm{bar}}^2 = 0$
(Theorem~\ref{thm:genus-induction-strict}). Since
$d_{\mathrm{cobar}} = \mathbb{D}(d_{\mathrm{bar}})$ and Verdier
duality is exact on holonomic $\mathcal{D}$-modules,
\[
d_{\mathrm{cobar}}^2
= \mathbb{D}(d_{\mathrm{bar}}^2)
= \mathbb{D}(0)
=0.
\]

In the curved CDG lane the bar side carries curvature $h_{\bar B}$ and
satisfies $d_{\bar B,h}^{\,2}=[h_{\bar B},-]$. Verdier duality
transports $h_{\bar B}$ to a curvature element
$h_\Omega=\mathbb{D}(h_{\bar B})$ on the cobar side and transports the
inner commutator identity to
Equation~\eqref{eq:curved-cobar-square}. The Bianchi identity
$d_{\Omega,h}(h_\Omega)=0$ is the CDG relation dual to
$d_{\bar B,h}(h_{\bar B})=0$.

For the twist, expand
$d_{\Omega,\alpha}=d_{\Omega,h}+[\alpha,-]$ and use the graded Jacobi
identity. Since $h_\Omega$ has even degree,
\[
(d_{\Omega,h}+[\alpha,-])^2
=[h_\Omega+d_{\Omega,h}\alpha+\tfrac12[\alpha,\alpha],-],
\]
which is Equation~\eqref{eq:curved-cobar-twist-square}. The right-hand
side vanishes for a curvature-killing Maurer--Cartan solution. Without
such a solution, ordinary cohomology of the raw curved object is not
the invariant; Positselski's coderived totalization records the CDG
object with its curvature.
\end{proof}

\begin{remark}[Resolution of delta-function products]\label{rem:delta-product-resolution}
A classical objection to the distributional cobar is that products of delta
distributions may be ill-defined: the product
$\delta(z_1 - z_2) \cdot \delta(z_2 - z_3)$ requires a renormalization
choice in the analytic category. The intrinsic
definition (Definition~\ref{def:geom-cobar-intrinsic}) avoids this
product: the cobar differential is the Verdier transpose of the bar
differential in the holonomic $\mathcal{D}$-module category, and
compositions are pull--push compositions through the normal-crossing
strata of $\overline{C}_n(X)$. When a displayed distributional product
fails the wavefront transversality criterion, it is only a local
notation for this Verdier composition; it is not an additional analytic
definition of the cobar differential.

The distributional representation
(Theorem~\ref{thm:cobar-distributional-model}) is a \emph{consequence}: it
describes how the intrinsic $\mathcal{D}$-module object looks when evaluated
on test sections, but the uncurved nilpotence and the curved inner-square
identity are established at the $\mathcal{D}$-module level, where no
analytic subtleties arise.
\end{remark}


\subsection{Explicit distributional model}\label{subsec:distributional-model}

\begin{definition}[Geometric cobar complex]\label{def:geom-cobar}
The bar complex uses logarithmic forms on \emph{compactified}
configuration spaces; the cobar uses the Verdier-dual current model on
the open stratum together with controlled boundary currents on the
Fulton--MacPherson divisors. For the Heisenberg algebra
(\S\ref{sec:frame-koszul-dual}), the cobar of
$\mathrm{coLie}^{\mathrm{ch}}(V^*)$ produced
$\mathrm{Sym}^{\mathrm{ch}}(V^*)$; the following definition gives the
general construction.

For a conilpotent chiral coalgebra $\mathcal{C}$ on $X$ with coaugmentation
$\eta: \omega_X \to \mathcal{C}$ and comultiplication $\Delta: \mathcal{C} \to
\mathcal{C} \boxtimes \mathcal{C}$, the \emph{geometric cobar complex} is:
\[
\Omega^{\text{ch}}_{p,q}(\mathcal{C}) = \Gamma\left(C_{p+1}(X), \text{Hom}_{\mathcal{D}}(\pi^*\mathcal{C}^{\otimes(p+1)}, \mathcal{D}_{C_{p+1}(X)}) \otimes \Omega^q_{C_{p+1}(X),\text{dist}}\right)
\]
where:
\begin{itemize}
\item $C_{p+1}(X)$ is the open stratum of the
Fulton--MacPherson compactification
\item $\pi: C_{p+1}(X) \to X^{p+1}$ is the projection
\item $\Omega^q_{C_{p+1}(X),\text{dist}}$ are distributional $q$-forms: currents with
prescribed boundary values along the collision divisors
$D_{ij}\subset\overline{C}_{p+1}(X)$
\item $\text{Hom}_{\mathcal{D}}$ denotes $\mathcal{D}$-module homomorphisms
\end{itemize}

Equivalently, using the Schwartz kernel theorem (Theorem~\ref{thm:schwartz-kernel-cobar}):
\[\Omega^{\text{ch}}_n(\mathcal{C}) = \text{Dist}\left(C_n(X), \mathcal{C}^{\boxtimes n}\right)
\otimes \Omega^*_{C_n(X)}\]
consisting of distributional sections of $\mathcal{C}^{\boxtimes n}$
on the FM current model, restricted to the open stratum.
\end{definition}

\begin{remark}[Distributions]\label{rem:why-distributions}
The cobar differential inserts boundary currents
$\delta_{D_{ij}}$, written locally as $\delta(z_i-z_j)$, so the
distributional notation lives on the FM compactification rather than on
the open configuration space alone. Geometrically, these currents are
Verdier-dual to logarithmic forms on $\overline{C}_n(X)$. Under a QFT
comparison datum, the same boundary current maps to the source term in
a Green-kernel equation:
$(\Box - m^2)G(z,w)=\delta^{(2)}(z-w)$.
\end{remark}

\begin{example}[Simplest cobar element]\label{ex:simplest-cobar}
For $n=2$ with trivial coalgebra $\mathcal{C} = \omega_X$, the basic
cobar boundary current is
\[K_2=\delta_{D_{12}}\otimes (dz\wedge d\bar z),\]
where $D_{12}\subset \overline{C}_2(X)$ is the collision divisor. It
acts on a test density $\phi$ on $\overline{C}_2(X)$ whose boundary
trace is defined by
\[
\langle K_2,\phi\rangle
=\int_X \operatorname{Tr}_{D_{12}}(\phi)(z)\,dz\wedge d\bar z.
\]
Thus $K_2$ is not an ordinary distribution on the open manifold
$C_2(X)$; it is the FM boundary-current representative of the
Verdier-dual residue.

\emph{QFT comparison.} In the free scalar comparison this is the
diagonal source current for the Green kernel, not the Green kernel
itself.
\end{example}

\begin{theorem}[Cobar differential; \ClaimStatusProvedHere]\label{thm:cobar-diff-geom}
The cobar differential is a degree +1 operator:
\[d_{\text{cobar}}: \Omega^{\text{ch}}_{p,q}(\mathcal{C}) \to
\Omega^{\text{ch}}_{p-1,q+1}(\mathcal{C}) \oplus \Omega^{\text{ch}}_{p,q}(\mathcal{C})
\oplus \Omega^{\text{ch}}_{p+1,q}(\mathcal{C})\]

It decomposes into three components:
\[
d_{\text{cobar}} = d_{\text{comult}} + d_{\text{internal}} + d_{\text{extend}}
\]
where each component has precise meaning:

\emph{Component 1: Comultiplication differential.}
\[d_{\text{comult}}: \Omega^{\text{ch}}_{p,q}(\mathcal{C}) \to
\Omega^{\text{ch}}_{p-1,q}(\mathcal{C})\]
Uses the comultiplication $\Delta: \mathcal{C} \to \mathcal{C} \boxtimes \mathcal{C}$
to split configurations. For $K \in \Omega^{\text{ch}}_n(\mathcal{C})$ represented as:
\[K = \int_{C_n(X)} k(z_1, \ldots, z_n) \otimes c_1(z_1) \otimes \cdots \otimes c_n(z_n)\]

We have:
\[(d_{\text{comult}}K)(c_0, \ldots, c_{n-2}) = \sum_{i=0}^{n-2} (-1)^{\epsilon_i}
K(c_0, \ldots, \Delta(c_i), \ldots, c_{n-2})\]
where $\epsilon_i = |c_0| + \cdots + |c_{i-1}|$ is the Koszul sign.

\emph{Component 2: Internal differential.}
\[d_{\text{internal}}: \Omega^{\text{ch}}_{p,q}(\mathcal{C}) \to
\Omega^{\text{ch}}_{p,q}(\mathcal{C})\]
Applies the internal differential of $\mathcal{C}$ coefficient-wise:
\[(d_{\text{internal}}K)(c_0, \ldots, c_n) = \sum_{i=0}^n (-1)^{\epsilon_i}
K(c_0, \ldots, d_{\mathcal{C}}(c_i), \ldots, c_n)\]

\emph{Component 3: Extension differential.}
\[d_{\text{extend}}: \Omega^{\text{ch}}_{p,q}(\mathcal{C}) \to
\Omega^{\text{ch}}_{p+1,q}(\mathcal{C})\]
This extends distributions across collision divisors and is the
Verdier adjoint of taking residues in the bar complex.

For a distribution $K$ on $C_n(X)$ with singularities along $\Delta_{ij} =
\{z_i = z_j\}$:
\[
(d_{\text{extend}}K)(z_0, \ldots, z_n)
=\sum_{i<j}\iota_{ij,!}\operatorname{Tr}_{D_{ij}}(K),
\]
written in local analytic notation as
$\sum_{i<j}\delta(z_i-z_j)\otimes K|_{D_{ij}}$.

\end{theorem}

\begin{proof}
We construct each component explicitly with all signs and conventions.

\emph{Step 1: Comultiplication component.}

For $K \in \Omega^{\text{ch}}_n(\mathcal{C})$, write:
\[K = \sum_{\sigma \in \mathfrak{S}_n} K_\sigma \otimes c_{\sigma(1)} \otimes \cdots
\otimes c_{\sigma(n)}\]
where $K_\sigma \in \mathcal{D}'(C_n(X))$ and $c_i \in \mathcal{C}$.

The comultiplication differential acts by:
\[(d_{\text{comult}}K)(c_1, \ldots, c_{n-1}) = \sum_{i=1}^{n-1} \sum_{\Delta(c_i) =
\sum c_i' \otimes c_i''} (-1)^{\epsilon_i} K(c_1, \ldots, c_{i-1}, c_i', c_i'',
c_{i+1}, \ldots, c_{n-1})\]

\emph{Sign convention.} $\epsilon_i = |c_1| + \cdots + |c_{i-1}|$ accounts for
moving $c_i$ past previous elements.

\emph{Geometric picture.} In local coordinates $(z_1, \ldots, z_n)$ on $C_n(X)$:
\[
(d_{\text{comult}}K)(z_1, \ldots, z_{n-1}) = \int_X K(z_1, \ldots, z_i, w, z_{i+1},
\ldots, z_{n-1}) \otimes \Delta_w
\]
where $\Delta_w$ is the coproduct evaluated at point $w \in X$, and we sum over
all insertion positions $i$.

\emph{Step 2: Internal component.}

\[(d_{\text{internal}}K)(c_1, \ldots, c_n) = \sum_{i=1}^n (-1)^{|c_1| + \cdots +
|c_{i-1}|} K(c_1, \ldots, d_{\mathcal{C}}(c_i), \ldots, c_n)\]

\emph{Step 3: Extension component.}

The extension differential:
\[
d_{\text{extend}}\colon
\mathrm{Dist}_{\mathrm{FM}}(C_n(X))\to
\mathrm{Dist}_{\mathrm{FM}}(C_{n+1}(X))
\]
extends FM currents by inserting boundary currents at collision loci.

\emph{Local coordinate formula.} Near the boundary divisor
$D_{ij}\subset\overline{C}_n(X)$ lying over $z_i=z_j$, introduce
coordinates:
\[\epsilon = z_i - z_j, \quad \zeta = \frac{z_i + z_j}{2}, \quad z_k \text{ for } k
\neq i,j\]

A distribution $K$ singular along $\Delta_{ij}$ has Laurent expansion:
\[K(\epsilon, \zeta, \{z_k\}) = \sum_{m=-\infty}^{M} \frac{K_m(\zeta, \{z_k\})}{\epsilon^m}
+ \text{(regular terms)}\]

The extension across $D_{ij}$ is:
\[
(d_{\text{extend}}K)(z_1, \ldots, z_n, w) = \sum_{i<j} \delta(z_i - z_j) \otimes
\text{Res}_{\epsilon=0}[K] \otimes \delta(w - \zeta)
\]

\emph{Explicit formula using regularization.}
\[\langle d_{\text{extend}}K, \phi \rangle = \lim_{\epsilon_0 \to 0} \int_{|z_i - z_j|
< \epsilon_0} K \cdot \phi - \text{(regularization counterterms)}\]

The regularization removes divergences, leaving a finite distributional value.

\emph{Example.} For $K = \frac{1}{(z_1 - z_2)^2}$, let
$\epsilon = z_1 - z_2$. The $d\log$ kernel $\eta_{12} = d\epsilon/\epsilon$
of the bar complex absorbs one power of $\epsilon$ (pole-order
absorption), so the Verdier dual of a double pole is a
\emph{single} distribution of order~$2$ supported on
$\Delta_{12} = \{z_1 = z_2\}$, not a tensor product of distributions:
\[
d_{\text{extend}}\!\left[\frac{1}{(z_1 - z_2)^2}\right]
= -\,\delta'(z_1 - z_2)
\]
where $\delta'$ is the derivative of the delta distribution, paired with
test functions by $\langle \delta', \phi \rangle = -\phi'(0)$. In
general, a pole of order~$k$ in $\epsilon$ dualizes to
$(-1)^{k-1}(k-1)!^{-1}\,\delta^{(k-1)}(\epsilon)$, a derivative of
order~$k-1$: one less than the pole order, because the $d\log$ kernel
on the bar side has already absorbed one pole.
\end{proof}

\begin{theorem}[Uncurved distributional verification of
\texorpdfstring{$d_{\text{cobar}}^2 = 0$}{d\_cobar\textasciicircum 2 = 0};
\ClaimStatusProvedHere]\label{thm:cobar-d-squared-zero}
Assume the uncurved conilpotent regime of
Corollary~\ref{cor:cobar-nilpotence-verdier}. The rigorous proof of
$d_{\mathrm{cobar}}^2 = 0$ is the Verdier-duality argument there; the
present verification is a distributional consistency check. It shows
that the distributional model, where the displayed products are defined
or regularized compatibly with Verdier duality, gives the same
square-zero answer. At Term~3 below, the product
$\delta(z_i - z_j)^2$ is formally undefined; the Arnold cancellation is
a computational shadow of the Verdier proof, not an independent
definition of the product.
The uncurved cobar differential satisfies $d_{\text{cobar}}^2 = 0$.
This requires verifying the diagonal and off-diagonal cancellation
families mirroring the bar complex from Section~\ref{sec:bar-nilpotency}:

\[d_{\text{cobar}}^2 = (d_{\text{comult}} + d_{\text{internal}} + d_{\text{extend}})^2
= \sum_{i,j} d_i \circ d_j = 0\]

\emph{The cancellation families.}
\begin{enumerate}
\item $d_{\text{comult}}^2 = 0$ (coassociativity)
\item $d_{\text{internal}}^2 = 0$ (differential property)
\item $d_{\text{extend}}^2 = 0$ (Verdier dual of the residue-square
identity, equivalently the Arnold boundary relation)
\item $d_{\text{comult}} \circ d_{\text{internal}} + d_{\text{internal}} \circ
d_{\text{comult}} = 0$ (chain map property)
\item $d_{\text{comult}} \circ d_{\text{extend}} + d_{\text{extend}} \circ
d_{\text{comult}} = 0$ (compatibility)
\item $d_{\text{internal}} \circ d_{\text{extend}} + d_{\text{extend}} \circ
d_{\text{internal}} = 0$ (compatibility)
\end{enumerate}
\end{theorem}

\begin{proof}
We verify each term systematically, providing the geometric and algebraic reasoning.

\emph{Term 1: $d_{\text{comult}}^2 = 0$}

This follows from coassociativity of the comultiplication $\Delta$. By definition:
\[(\Delta \otimes \text{id}) \circ \Delta = (\text{id} \otimes \Delta) \circ \Delta\]

Applied twice:
\begin{align*}
d_{\text{comult}}^2(K)(c_1, \ldots, c_{n-2}) &= \sum_{i < j} (-1)^{\epsilon_i +
\epsilon_j} K(\ldots, \Delta(c_i), \ldots, \Delta(c_j), \ldots) \\
&= \sum_{i < j} (-1)^{\epsilon_i + \epsilon_j} K(\ldots, (\Delta \otimes \text{id})
\Delta(c_i), \ldots)
\end{align*}

By coassociativity, terms with different orderings cancel pairwise. This completes term~1.

\emph{Term 2: $d_{\text{internal}}^2 = 0$}

In the uncurved lane the internal differential acts on each tensor
factor independently:
\[
d_{\text{internal}}(K)(c_1, \ldots, c_n) = \sum_i K(c_1, \ldots, d_{\mathcal{C}}(c_i), \ldots, c_n).
\]
Applying $d_{\text{internal}}$ twice, each~$c_i$ contributes
$d_{\mathcal{C}}^2(c_i) = 0$ (since $(\mathcal{C}, d_{\mathcal{C}})$
is a cochain complex), and cross-terms
$d_{\mathcal{C}}(c_i) \otimes d_{\mathcal{C}}(c_j)$ for $i \neq j$
appear with equal signs in both applications, hence cancel:
\[d_{\text{internal}}^2(K) = \sum_i K(\ldots, d_{\mathcal{C}}^2(c_i), \ldots) + \sum_{i < j}(\cdots) - \sum_{i < j}(\cdots) = 0.\]
This completes term~2.

\emph{Term 3: $d_{\text{extend}}^2 = 0$}

The extension differential inserts delta functions. Applied twice:
\begin{align*}
d_{\text{extend}}^2(K) &= d_{\text{extend}}\left(\sum_{i<j} \delta(z_i - z_j) \otimes
K|_{\Delta_{ij}}\right) \\
&= \sum_{i<j} \sum_{k<\ell} \delta(z_i - z_j) \otimes \delta(z_k - z_\ell) \otimes
K|_{\Delta_{ij} \cap \Delta_{k\ell}}
\end{align*}

The displayed product is not a definition. If read as a product of
ordinary distributions, terms such as $\delta(z_i-z_j)^2$ are generally
undefined. In the intrinsic model one computes instead
\[
d_{\mathrm{extend}}^2
=\mathbb{D}(d_{\mathrm{res}})^2
=\mathbb{D}(d_{\mathrm{res}}^2).
\]
The bar residue differential satisfies $d_{\mathrm{res}}^2=0$ because
the oriented codimension-two boundary of
\(\overline{C}_n(X)\) is the Arnold--Orlik--Solomon combination of
three collision faces. Therefore
$d_{\mathrm{extend}}^2=0$ without forming any product
$\delta\cdot\delta$. If one chooses a regulator
$\delta_\epsilon$ and expands the mnemonic expression, the divergent
terms occur in precisely this Arnold combination; the regularized
calculation is a consistency check on the Verdier proof, not a
replacement for it.

This completes term~3.

\emph{Term 4: $d_{\text{comult}} \circ d_{\text{internal}} + d_{\text{internal}}
\circ d_{\text{comult}} = 0$}

This states that $\Delta: \mathcal{C} \to \mathcal{C} \boxtimes \mathcal{C}$ is a
chain map (compatible with the differential). By hypothesis:
\[\Delta \circ d_{\mathcal{C}} = (d_{\mathcal{C}} \otimes \text{id} + \text{id}
\otimes d_{\mathcal{C}}) \circ \Delta\]

Applied to cobar elements:
\begin{align*}
(d_{\text{comult}} \circ d_{\text{internal}})(K) &= d_{\text{comult}}\left(\sum_i
K(\ldots, d_{\mathcal{C}}(c_i), \ldots)\right) \\
&= \sum_{i,j} K(\ldots, \Delta(d_{\mathcal{C}}(c_i)), \ldots)
\end{align*}

By the chain map property:
\[\Delta(d_{\mathcal{C}}(c_i)) = (d_{\mathcal{C}} \otimes \text{id} + \text{id}
\otimes d_{\mathcal{C}})(\Delta(c_i))\]

Substituting and using Koszul signs, this precisely cancels $(d_{\text{internal}}
\circ d_{\text{comult}})(K)$. This completes term~4.

\emph{Term 5: $d_{\text{comult}} \circ d_{\text{extend}} + d_{\text{extend}}
\circ d_{\text{comult}} = 0$}

\emph{Calculation.}
\begin{align*}
(d_{\text{comult}} \circ d_{\text{extend}})(K) &= d_{\text{comult}}\left(\sum_{i<j}
\delta(z_i - z_j) \otimes K|_{\Delta_{ij}}\right) \\
&= \sum_{i<j} \sum_k \delta(z_i - z_j) \otimes \Delta_k(K|_{\Delta_{ij}})
\end{align*}

where $\Delta_k$ applies the coproduct at position $k$.

Similarly:
\begin{align*}
(d_{\text{extend}} \circ d_{\text{comult}})(K) &= d_{\text{extend}}\left(\sum_k
\Delta_k(K)\right) \\
&= \sum_k \sum_{i<j} \delta(z_i - z_j) \otimes (\Delta_k(K))|_{\Delta_{ij}}
\end{align*}

By the Leibniz rule for distributions:
\[\delta(z_i - z_j) \otimes \Delta_k(K) = \Delta_k(\delta(z_i - z_j) \otimes K)
\quad \text{if } k \notin \{i, j\}\]

For $k \in \{i,j\}$, the coproduct \emph{splits the collision point}, and the
contributions from the two orderings cancel by coassociativity.

All terms cancel pairwise, completing term~5.

\emph{Term 6: $d_{\text{internal}} \circ d_{\text{extend}} + d_{\text{extend}}
\circ d_{\text{internal}} = 0$}

\emph{Calculation.}
\begin{align*}
(d_{\text{internal}} \circ d_{\text{extend}})(K)(c_1, \ldots, c_n) &=
d_{\text{internal}}\left(\sum_{i<j} \delta(z_i - z_j) \otimes K|_{\Delta_{ij}}\right) \\
&= \sum_{i<j} \sum_k (-1)^{\epsilon_k} \delta(z_i - z_j) \otimes
K|_{\Delta_{ij}}(c_1, \ldots, d_{\mathcal{C}}(c_k), \ldots)
\end{align*}

Similarly:
\begin{align*}
(d_{\text{extend}} \circ d_{\text{internal}})(K) &= d_{\text{extend}}\left(\sum_k
(-1)^{\epsilon_k} K(\ldots, d_{\mathcal{C}}(c_k), \ldots)\right) \\
&= \sum_k \sum_{i<j} (-1)^{\epsilon_k} \delta(z_i - z_j) \otimes
(K(\ldots, d_{\mathcal{C}}(c_k), \ldots))|_{\Delta_{ij}}
\end{align*}

The differential $d_{\mathcal{C}}$ acts coefficient-wise,
while $\delta(z_i - z_j)$ acts geometrically. They commute as operators:
\[[\delta(z_i - z_j), d_{\mathcal{C}}(c_k)] = 0\]

The two compositions differ by a Koszul sign: commuting the degree-$1$ operator $d_{\mathcal{C}}$ past the geometric factor $\delta(z_i - z_j)$ introduces $(-1)^1$, so $d_{\mathrm{internal}} \circ d_{\mathrm{extend}} = -d_{\mathrm{extend}} \circ d_{\mathrm{internal}}$, and the anticommutator vanishes.

This completes the verification of $d^2 = 0$ in the uncurved lane. All
diagonal and off-diagonal families vanish:
\begin{center}
\begin{tabular}{c|ccc}
& $d_{\text{comult}}$ & $d_{\text{internal}}$ & $d_{\text{extend}}$ \\ \hline
$d_{\text{comult}}$ & coassoc. & chain map & Leibniz \\
$d_{\text{internal}}$ & chain map & $d^2=0$ & commute \\
$d_{\text{extend}}$ & Leibniz & commute & Arnold \\
\end{tabular}
\end{center}

Therefore, in the uncurved conilpotent regime:
\[d_{\text{cobar}}^2 = 0\]

This establishes the cobar construction as a chain complex, and more
precisely as a differential graded algebra with the $A_\infty$
structure, only in the uncurved regime. In the curved CDG regime the
same calculation leaves the inner curvature term
$[h_\Omega,-]$ of Equation~\eqref{eq:curved-cobar-square}.
\end{proof}

\begin{remark}[Duality with the uncurved bar
\texorpdfstring{$d^2=0$}{d\textasciicircum 2=0} proof]\label{rem:bar-cobar-d2-duality}
The structure of this proof mirrors the uncurved bar $d^2=0$ proof from
Section~\ref{sec:bar-nilpotency}:

\begin{center}
\begin{tabular}{l|l}
\textbf{Bar (Section~\ref{sec:bar-nilpotency})} & \textbf{Cobar (this section)} \\ \hline
Residues at divisors & Delta functions at diagonals \\
Compactified space $\overline{C}_n(X)$ & Open space $C_n(X)$ \\
Logarithmic forms & Distributional currents \\
Stratification by collisions & Singular support on diagonals \\
Arnold--Orlik--Solomon relations & Arnold relations for distributions \\
Extract (analyze) & Insert (synthesize) \\
\end{tabular}
\end{center}

The proofs are dual under Verdier duality, as is the bar-cobar
adjunction itself. In the curved lane the same duality transports the
bar curvature into the cobar curvature and gives
$d_{\Omega,h}^{\,2}=[h_\Omega,-]$.
\end{remark}

\subsection{Sign conventions for cobar operations}


\begin{convention}[Cobar sign system]\label{conv:cobar-signs}
The cobar complex inherits signs from three sources. Throughout the
monograph the signs appendix
(Appendix~\ref{app:signs}, \S\ref{app:sign-conventions}) is
authoritative; this convention records the chapter-local shorthand.

\emph{Desuspension sign (from $s^{-1}$).}
The cobar complex uses desuspension: $|s^{-1}v| = |v| - 1$ (lowers
cohomological degree by one). In particular, a degree-$h$ generator
$v \in \bar{\mathcal{C}}$ becomes a degree-$(h{-}1)$ generator
$s^{-1}v \in \Omegach(\mathcal{C})$, and the cobar differential
$d_\Omega$ has degree $+1$. The cofree-algebra structure on
$\Omegach(\mathcal{C})$ is on the desuspension $s^{-1}\bar{\mathcal{C}}$,
never on $\bar{\mathcal{C}}$ itself. This is the same convention as
the bar complex (both bar and cobar use desuspension in this
monograph; see the signs appendix).

\emph{Koszul signs (from grading).}
The standard Koszul sign rule (Appendix~\ref{app:signs}, \S\ref{app:sign-conventions}).

\emph{Symmetry signs (from permutations).}
The symmetric group $\mathfrak{S}_n$ acts on $C_n(X)$ and $\mathcal{C}^{\boxtimes n}$.
For $\sigma \in \mathfrak{S}_n$ and elements $c_1, \ldots, c_n$:
\[\sigma(c_1 \otimes \cdots \otimes c_n) = (-1)^{\epsilon(\sigma, c)}
c_{\sigma(1)} \otimes \cdots \otimes c_{\sigma(n)}\]
where $\epsilon(\sigma, c)$ is the Koszul sign for moving graded elements according
to $\sigma$.

\emph{Boundary-current signs (from orientation lines).}
The symbol $\delta_{D_{ij}}$ denotes the orientation line of the
codimension-one FM boundary divisor $D_{ij}$, transported by Verdier
duality from the logarithmic residue map. Interchanging two boundary
traces contributes the Koszul sign of the two degree-one residue
operators, together with the sign of the induced permutation of the
coefficient factors. No product of ordinary delta distributions is
part of the definition.
\end{convention}

\begin{lemma}[Sign consistency for the uncurved cobar differential; \ClaimStatusProvedHere]\label{lem:cobar-sign-consistency}
The sign conventions above ensure that, in the uncurved cobar
differential, double applications of the elementary operations carry
the signs needed for the Verdier square-zero identity.
\end{lemma}

\begin{proof}
The extension operator is the Verdier transpose of the residue
operator:
\[
d_{\mathrm{extend}}=\mathbb D(d_{\mathrm{res}}).
\]
The bar residue operator is a signed sum over oriented boundary
divisors of $\overline C_n(X)$. On a codimension-two stratum there are
two ways to arrive by successive residues. If the collisions are
disjoint, the two oriented paths differ by the interchange of two
degree-one residue maps. If the collisions are nested, the three
oriented paths around the triple-collision stratum occur in the
Arnold--Orlik--Solomon combination
\[
\eta_{ij}\wedge\eta_{jk}
+\eta_{jk}\wedge\eta_{ki}
+\eta_{ki}\wedge\eta_{ij}=0 .
\]
Thus $d_{\mathrm{res}}^2=0$ in the uncurved bar complex with the signs
fixed by the orientation lines of the FM boundary strata. Applying
Verdier duality gives
\[
d_{\mathrm{extend}}^2
=\mathbb D(d_{\mathrm{res}}^2)=0.
\]
The displayed current notation in Examples below records these
orientation signs; it never defines a product
$\delta_{D_{ij}}\delta_{D_{kl}}$ in the ordinary distribution algebra.
\end{proof}

\begin{example}[Three-point sign shadow of Verdier nilpotence]\label{ex:three-point-signs-cobar}
Consider the uncurved cobar complex for $n=3$ with
$\mathcal{C} = \omega_X$ (trivial). Elements are:
\[K_3(z_1, z_2, z_3) = \sum_{\text{perms}} k_\sigma(z_1, z_2, z_3) \cdot
\text{sgn}(\sigma)\]

Apply $d_{\text{extend}}$:
\begin{align*}
d_{\text{extend}}(K_3) &= \delta(z_1 - z_2) \otimes K_3|_{z_1=z_2} \\
&\quad + \delta(z_2 - z_3) \otimes K_3|_{z_2=z_3} \\
&\quad + \delta(z_1 - z_3) \otimes K_3|_{z_1=z_3}
\end{align*}

Apply again:
\begin{align*}
d_{\text{extend}}^2(K_3) &= \delta(z_1 - z_2) \wedge \delta(z_2 - z_3) \otimes
K_3|_{z_1=z_2=z_3} \\
&\quad + \delta(z_2 - z_3) \wedge \delta(z_1 - z_3) \otimes K_3|_{z_1=z_2=z_3} \\
&\quad + \delta(z_1 - z_3) \wedge \delta(z_1 - z_2) \otimes K_3|_{z_1=z_2=z_3}
\end{align*}

After replacing the displayed products by Verdier pull--push
composition along the three codimension-two boundary strata, the
transported Arnold relation is:
\[\delta(z_1 - z_2) \wedge \delta(z_2 - z_3) + \delta(z_2 - z_3) \wedge \delta(z_1 - z_3) + \delta(z_1 - z_3) \wedge \delta(z_1 - z_2) = 0\]
so the three Verdier-transposed terms in
$d_{\text{extend}}^2(K_3)$ sum to zero:
\[\text{term}_1 + \text{term}_2 + \text{term}_3 = 0\]

Hence $d_{\text{extend}}^2(K_3) = 0$ in the uncurved three-point
Verdier model. The display records the sign shadow of
Corollary~\ref{cor:cobar-nilpotence-verdier}; it is not an independent
definition of a product of singular distributions.
\end{example}

\subsection{Low-degree explicit computations}

\begin{example}[Divided-power cobar differential through degree 5]
\label{ex:cobar-linear-complete}

Let
$\mathcal{C}=\operatorname{Sym}^{c}_{\mathrm{ch}}(V)$ be the
cofree cocommutative chiral coalgebra on
$V=\operatorname{span}\{v\}$ with $|v|=h$, written in the ordinary
monomial normalization. Its coproduct is
\[
\Delta(v^n)=\sum_{k=0}^{n}\binom{n}{k}v^k\otimes v^{n-k}.
\]
The reduced coproduct omits the two counit summands $k=0,n$.

\emph{Cobar complex.}
\[
\Omega^{\text{ch}}(\mathcal{C})
=\mathrm{Free}_{\text{ch}}\bigl(s^{-1}\bar{\mathcal C}\bigr),
\qquad
\bar{\mathcal C}=\bigoplus_{n\geq1}\mathbb C\,v^n .
\]

\emph{Generators.} $s^{-1}v, s^{-1}v^2, s^{-1}v^3, s^{-1}v^4, s^{-1}v^5, \ldots$
in degrees $h-1, 2h-1, 3h-1, 4h-1, 5h-1, \ldots$ respectively.

\emph{Differential formulas.}

\emph{Degree 1 (h-1):}
\[d(s^{-1}v) = 0\]
(Primitive element, no coproduct.)

\emph{Degree 2 (2h-1):}
The cobar differential uses the \emph{reduced} coproduct $\bar{\Delta}$, which excludes the counit terms:
\begin{align*}
d(s^{-1}v^2) &= -\bar{\Delta}(v^2) = -\sum_{k=1}^{1} \binom{2}{k} (s^{-1}v^k) \cdot (s^{-1}v^{2-k}) \\
&= -\binom{2}{1}(s^{-1}v) \cdot (s^{-1}v) \\
&= -2(s^{-1}v)^2
\end{align*}

\emph{Degree 3 (3h-1):}
\begin{align*}
d(s^{-1}v^3) &= -\bar{\Delta}(v^3) = -\sum_{k=1}^{2} \binom{3}{k} (s^{-1}v^k) \cdot (s^{-1}v^{3-k}) \\
&= -3(s^{-1}v) \cdot (s^{-1}v^2) - 3(s^{-1}v^2) \cdot (s^{-1}v)
\end{align*}

In a commutative setting:
\[d(s^{-1}v^3) = -6(s^{-1}v) \cdot (s^{-1}v^2)\]

\emph{Degree 4 (4h-1):}
\[d(s^{-1}v^4) = -4(s^{-1}v) \cdot (s^{-1}v^3) - 6(s^{-1}v^2) \cdot (s^{-1}v^2)\]

\emph{Degree 5 (5h-1):}
\[d(s^{-1}v^5) = -5(s^{-1}v) \cdot (s^{-1}v^4) - 10(s^{-1}v^2) \cdot (s^{-1}v^3)\]

\emph{General pattern.} For generator $s^{-1}v^n$:
\[d(s^{-1}v^n) = -\sum_{k=1}^{n-1} \binom{n}{k} (s^{-1}v^k) \cdot (s^{-1}v^{n-k})\]

\emph{Geometric interpretation.} These formulas encode how one
divided-power insertion with charge $v^n$ splits into two insertion
points with charges $v^k$ and $v^{n-k}$, weighted by the chosen
monomial normalization. This example is a coefficient and sign check
for the cobar differential. It is not, by itself, a bar-cobar
inversion theorem: the cohomology depends on the target algebra whose
bar coalgebra is being resolved. If the coalgebra is replaced by the
primitive truncation $\omega_X\oplus V$ with
$\bar\Delta(V)=0$, then
$\Omega^{\mathrm{ch}}(\omega_X\oplus V)=
\mathrm{Free}_{\mathrm{ch}}(s^{-1}V)$ with zero cobar differential.
\end{example}

\begin{example}[Cobar of exterior coalgebra: free fermions]\label{ex:cobar-fermion-complete}

Let $\mathcal{C} = \Lambda^*_{\text{ch}}(V)$ be the chiral exterior coalgebra on
$V = \text{span}\{\psi\}$ with $|\psi| = \frac{1}{2}$ (fermionic). The comultiplication:
\[\Delta(\psi) = \psi \otimes 1 + 1 \otimes \psi, \quad \Delta(\psi^2) = 0\]
(since $\psi^2 = 0$ by anticommutativity).

\emph{Cobar complex.}
\[\Omega^{\text{ch}}(\Lambda^*_{\text{ch}}(V)) = \text{Free}_{\text{ch}}(s^{-1}\psi)\]

\emph{Generator.} $s^{-1}\psi$ in degree $-\frac{1}{2}$.

\emph{Differential.}
The reduced comultiplication $\bar{\Delta}$ removes the $1 \otimes \psi + \psi \otimes 1$
term. For the reduced coproduct:
\[\bar{\Delta}(\psi) = 0\]

Therefore:
\[d(s^{-1}\psi) = 0\]

\emph{Cohomology.}
\[H^*(\Omega^{\text{ch}}(\Lambda^*_{\text{ch}}(V))) = \text{Free}_{\text{ch}}(s^{-1}\psi)\]

The desuspension $s^{-1}$ converts the fermionic generator $\psi$ (with
anticommuting multiplication) into a bosonic generator $s^{-1}\psi$ (with commuting
multiplication in the free algebra).

\emph{Interpretation.} This is the cobar of the exterior coalgebra, not
the assertion that $\Omega\bar{B}(\mathcal F)$ is the
$\beta\gamma$ system. When this exterior coalgebra is the
Koszul-dual coalgebra $\mathcal F^i$, finite-type linear duality gives
the Koszul-dual algebra $(\mathcal F^i)^\vee$; the cobar construction
itself is the inverse operation appearing in
$\Omega(\bar B(\mathcal F))\simeq\mathcal F$. Desuspension converts the
fermionic generator $\psi$ into the bosonic generator
$\phi=s^{-1}\psi$ inside this cobar algebra.
\end{example}

\begin{example}[Free fermion to beta-gamma via bar-cobar]
\label{ex:fermion-betagamma-bar-cobar}

\emph{The Koszul duality chain.}
\[\text{Free fermion algebra } \mathcal{F}
\xrightarrow{\text{bar}}
\bar{B}(\mathcal{F})
\xrightarrow{H^*}
\mathcal F^i
\xrightarrow{(-)^\vee}
\beta\gamma \text{ system}\]
The bar complex resolves the Koszul-dual coalgebra~$\mathcal F^i$.
Linear duality of this coalgebra yields the Koszul dual
\emph{algebra}~$\mathcal{F}^!=(\mathcal F^i)^\vee$. The cobar
construction $\Omega(\bar{B}(\mathcal{F}))$ instead recovers
$\mathcal{F}$ itself by bar-cobar inversion.

\emph{Explicit construction.}

\begin{theorem}[Fermion-boson Koszul duality; \ClaimStatusProvedHere]\label{thm:fermion-boson-koszul}

\smallskip\noindent
\textup{[Regime: quadratic
\textup{(}Convention~\textup{\ref{conv:regime-tags}}\textup{)}.]}
The $\beta\gamma$ system is the \emph{Koszul dual} of free fermions.
This is the chiral analog of classical
$\mathrm{Sym}(V) \leftrightarrow \Lambda(V^*)$ Koszul duality
(see Part~\ref{part:characteristic-datum} for the complete bar-cobar computation).

The Koszul duality correspondence exchanges:
\begin{enumerate}
\item \emph{Relations}: anticommuting fields $\{\psi, \psi\} = 0$ $\leftrightarrow$ symplectic bosons $[\beta,\gamma] = 1$; exterior relation $\psi \boxtimes \psi = 0$ $\leftrightarrow$ symplectic pairing $\langle \beta, \gamma \rangle = \frac{1}{z_1-z_2}$.

\item \emph{Algebra structure}: exterior algebra $\Lambda^*(\psi)$ $\leftrightarrow$ polynomial-type algebra; bar complex $\bar{B}(\mathcal{F}) = \Lambda^*(\psi, \partial\psi, \ldots)$ $\leftrightarrow$ $\bar{B}(\beta\gamma)$ with symplectic differential.

\item \emph{Statistics}: fermionic $\leftrightarrow$ bosonic.
\end{enumerate}
\end{theorem}

\begin{proof}
The free fermion algebra $\mathcal{F} = \Lambda^{\mathrm{ch}}(V)$ is generated by $V$ with the exterior (anticommutative) chiral product. By the classical Koszul duality $\Lambda(V)^! = \mathrm{Sym}(V^*)$ (Loday--Vallette~\cite{LV12}, Theorem~3.2.1), the chiral Koszul dual is $\mathcal{F}^! = \mathrm{Sym}^{\mathrm{ch}}(V^*)$, the commutative chiral algebra on $V^*$, which is the $\beta\gamma$ system with $\beta \in V^*$, $\gamma \in V$. The complete bar-cobar computation verifying the Koszul property (acyclicity of $B(\mathcal{F}) \otimes_\tau \mathcal{F}^!$) is carried out in Part~\ref{part:characteristic-datum}.
\end{proof}

\emph{Propagator}: The bosonic $\beta\gamma$ system has propagator:
\[\langle \beta(z) \gamma(w) \rangle = \frac{1}{z - w}\]

This matches the fermion two-point function via Koszul duality (not to be confused with bosonization, which is a distinct equivalence $\mathcal{F} \cong V_{\mathbb{Z}}$).
\end{example}

\begin{example}[Cobar \texorpdfstring{$A_\infty$}{A-infinity} operations: explicit formulas through \texorpdfstring{$n_5$}{n5}]
\label{ex:cobar-ainfty-n5}

The cobar construction carries a canonical $A_\infty$ structure. We compute the
first five operations explicitly.

\emph{Operation $n_1$: The differential}
\[n_1 = d_{\text{cobar}}: \Omega^n(\mathcal{C}) \to \Omega^{n+1}(\mathcal{C})\]
(Already computed above.)

\emph{Operation $n_2$: Convolution product}
\[n_2: \Omega^p(\mathcal{C}) \otimes \Omega^q(\mathcal{C}) \to \Omega^{p+q}(\mathcal{C})\]
(In cohomological grading, $m_k$ has degree $2-k$; for $k=2$, degree $0$: $p+q \mapsto p+q$.)

\emph{Formula.} For integration kernels $K_1, K_2$:
\[(n_2(K_1, K_2))(z_1, \ldots, z_{p+q-1}) = \int_X K_1(z_1, \ldots, z_p; w) \cdot
K_2(w, z_{p+1}, \ldots, z_{p+q-1}) \, dw\]

\emph{Geometric interpretation.} Glue two configuration spaces at a common point
$w$, then integrate over $w$.

\emph{Sign.} $(-1)^{|K_1| \cdot |K_2|}$ from Koszul rule.

\emph{Example.} For $K_1 = \frac{1}{z_1 - w}$, $K_2 = \frac{1}{w - z_2}$, using partial fractions:
\begin{align*}
n_2(K_1, K_2)(z_1, z_2) &= \oint \frac{1}{z_1 - w} \cdot \frac{1}{w - z_2} \, dw \\
&= \frac{1}{z_1 - z_2}\oint \left[\frac{1}{w - z_2} - \frac{1}{w - z_1}\right] dw \\
&= \frac{2\pi i}{z_1 - z_2} \quad \text{(residue at } w = z_2\text{)}
\end{align*}
This reproduces the Cauchy kernel $1/(z_1-z_2)$ up to the
normalization $2\pi i$.

\emph{Operation $n_3$: Triple kernel composition}
\[n_3: \Omega^{p_1}(\mathcal{C}) \otimes \Omega^{p_2}(\mathcal{C}) \otimes
\Omega^{p_3}(\mathcal{C}) \to \Omega^{p_1+p_2+p_3-2}(\mathcal{C})\]

\emph{Formula.}
\[(n_3(K_1, K_2, K_3))(z_1, \ldots, z_N) = \int_{X \times X} K_1(\ldots; w_1) \cdot
K_2(w_1, \ldots; w_2) \cdot K_3(w_2, \ldots) \, dw_1 dw_2\]

\emph{Geometric interpretation.} Glue three configuration spaces in a chain,
then integrate over the two gluing points.

\emph{Operation $n_4$: Fourfold operation}
\[n_4: \bigotimes_{i=1}^4 \Omega^{p_i}(\mathcal{C}) \to \Omega^{\sum p_i - 3}(\mathcal{C})\]

\emph{Formula.} Similar, but integrate over three intermediate points $w_1, w_2, w_3$.

\emph{Operation $n_5$: Fivefold operation}
\[n_5: \bigotimes_{i=1}^5 \Omega^{p_i}(\mathcal{C}) \to \Omega^{\sum p_i - 4}(\mathcal{C})\]

\emph{General pattern.}
\[n_k: \bigotimes_{i=1}^k \Omega^{p_i}(\mathcal{C}) \to \Omega^{\sum p_i - (k-2)}(\mathcal{C})\]

\emph{Geometric realization.} Integrate over the moduli space $\overline{M}_{0,k+1}$
of stable curves:
\[n_k(K_1, \ldots, K_k) = \int_{\overline{M}_{0,k+1}} K_1 \wedge \cdots \wedge K_k
\wedge \omega_{0,k+1}\]

\emph{QFT comparison.} If a BV/factorization comparison datum supplies
interaction vertices and a propagator kernel, the genus-$0$ operation
$n_k$ maps to a tree-level $k$-ary kernel composition. Loop corrections
come from the higher-genus modular bar/cobar object, not from the
genus-$0$ operation $n_k$ itself.

\emph{$A_\infty$ relations.} These operations satisfy:
\[\sum_{i+j=n+1} \sum_{k} (-1)^{\epsilon} n_i(\text{id}^{\otimes k} \otimes n_j
\otimes \text{id}^{\otimes (n-k-j)}) = 0\]

This encodes associativity up to homotopy, with $n_3$ measuring the failure of
$n_2$ to be associative, $n_4$ measuring the failure of $n_3$ to be coherent, etc.
\end{example}

\subsection{Conditional comparison with propagator kernels}

\begin{theorem}[Conditional cobar--propagator comparison;
\ClaimStatusConditional]\label{thm:cobar-propagator-comparison}
\label{conj:cobar-physical}
Let $\mathcal{C}$ be a coaugmented conilpotent chiral coalgebra and
let $\Omega^{\mathrm{ch}}(\mathcal{C})$ be its cobar object: a dg
chiral algebra in the uncurved genus-$0$ regime, and the curved CDG
object
$(\Omega^{\mathrm{ch}}_{\mathrm{curv}}(\mathcal{C}),d_{\Omega,h},h_\Omega)$
when curvature is present. Assume a perturbative QFT comparison datum
consisting of:
\begin{enumerate}[label=\textup{(\roman*)}]
\item a local BV/factorization theory on~$X$ with field bundle~$E$,
 kinetic operator $P\colon \Gamma(E)\to\Gamma(E^\vee\otimes\omega_X)$,
 BRST operator $Q_{\mathrm{BV}}$, and propagator distribution
 $G\in\mathcal{D}'(X^2,E\boxtimes E^\vee)$ satisfying
 \[
 P_zG(z,w)=\delta_\Delta(z,w)-\Pi_{\ker P}(z,w),
 \]
 where $\Pi_{\ker P}=0$ in the acyclic local model and is the harmonic
 projector on a compact curve;
\item a continuous kernel map
 \[
 \Phi_n\colon
 \operatorname{Dist}_{\mathrm{FM}}(C_n(X),\mathcal{C}^{\boxtimes n})
 \longrightarrow
 \mathcal{D}'(X^n,E^{\boxtimes n})
 \]
 compatible with wavefront restrictions, boundary traces, and the
 coalgebra coproduct;
\item compatibility with differentials:
 \[
 \Phi(d_{\Omega}K)
 =
 (Q_{\mathrm{BV}}+P)\Phi(K)
 \]
 on the source-current summand, and compatibility with products on the
 interaction vertices;
\item in the curved case, either $h_\Omega=0$, or a
 Maurer--Cartan element $\alpha$ with
 $h_\Omega+d_{\Omega,h}\alpha+\frac12[\alpha,\alpha]=0$, or a chosen
 coderived/flat totalization in which the curvature is absorbed.
\end{enumerate}
Then $\Phi$ sends the FM boundary current
$\delta_{D_{ij}}$ in the cobar differential to the diagonal source
$\delta_\Delta$ in the Green-kernel equation, and sends genus-$0$
cobar $A_\infty$ compositions to tree-level kernel compositions.
It does not identify an arbitrary cobar element with a propagator
kernel, and it does not identify
$\Omega^{\mathrm{ch}}(\mathcal{C})$ with the Koszul dual algebra.
If $\mathcal{C}\simeq \bar{B}^{\mathrm{ch}}(\mathcal{A})$, then
$\Omega^{\mathrm{ch}}(\mathcal{C})\simeq\mathcal{A}$ under
Corollary~\ref{cor:bar-cobar-inverse}; the Koszul dual is instead
the Verdier-dual factorization algebra
$\mathbb{D}_{\operatorname{Ran}}\bar{B}^{\mathrm{ch}}(\mathcal{A})$,
identified with $\mathcal{A}^!$ only on the Koszul locus.
For $g\geq1$, the comparison must use the genus-$g$ corrected
bar/cobar object: period corrections and the curvature term
$h_\Omega$ are additional data, not hidden genus-$0$ higher products.
\end{theorem}

\begin{proof}
The first assertion is formal from the definition of
$d_{\mathrm{extend}}$: it is the Verdier transpose of the residue map
and is represented by the FM boundary current $\delta_{D_{ij}}$
(Theorem~\ref{thm:cobar-diff-geom}). Condition~(ii) is precisely the
analytic datum needed to push this current to the ordinary diagonal
distribution in the QFT kernel space, and condition~(iii) says that
the pushed differential is the BV kinetic differential. Thus
$P_zG=\delta_\Delta-\Pi_{\ker P}$ is the analytic image of the
boundary-current identity, not the definition of the cobar
differential.

The tree statement uses only the genus-$0$ operadic bar-cobar
resolution: the operations are indexed by the boundary strata of
$\overline{M}_{0,k+1}$ and obey the $A_\infty$ identities by
Theorem~\ref{thm:cobar-ainfty}. Stable graphs with loops lie in the
higher-genus modular bar/cobar theory. In the curved lane,
Corollary~\ref{cor:cobar-nilpotence-verdier} gives
$d_{\Omega,h}^{\,2}=[h_\Omega,-]$, so ordinary cobar cohomology is
available only under condition~(iv). The last assertion is exactly the
separation of the three functors in
Remark~\ref{rem:cobar-three-functors}: cobar inverts the bar coalgebra,
whereas Koszul duality is Verdier duality of the bar object.
\end{proof}

\begin{remark}[Comparison data for QFT kernels]
The comparison theorem becomes an identification of a concrete QFT only
after four independent obligations are discharged:
\begin{enumerate}[label=\textup{(\alph*)}]
\item construct the BV/factorization model and its propagator
distribution with a fixed gauge, measure, and harmonic projector;
\item prove the wavefront-set and boundary-trace compatibility of
$\Phi$ for all collision strata used by the cobar differential;
\item prove the Ward/BRST identity
$\Phi(d_\Omega K)=(Q_{\mathrm{BV}}+P)\Phi(K)$ and the matching of
interaction vertices;
\item for $g\geq1$, prove compatibility with period corrections,
sewing, and the curved square
$d_{\Omega,h}^{\,2}=[h_\Omega,-]$ or exhibit a curvature-killing
Maurer--Cartan twist.
\end{enumerate}
Without these data, the algebraic statement is only the Verdier
bar-cobar construction, not a statement about BRST state spaces or
scattering amplitudes.
\end{remark}

\begin{example}[Free scalar parametrix and the cobar boundary current]\label{ex:free-scalar-cobar}
Let $P=-\Delta$ be the scalar kinetic operator in a local flat chart,
and let $G$ be a chosen parametrix:
\[
P_zG(z,w)=\delta^{(2)}(z-w)
\]
locally, or
$P_zG(z,w)=\delta_\Delta(z,w)-\Pi_{\ker P}(z,w)$ on a compact curve.
With the standard two-dimensional normalization,
$G(z,w)$ has logarithmic singular part
$-\frac{1}{2\pi}\log|z-w|$ modulo a smooth kernel.

The cobar comparison sends the FM current
$\delta_{D_{12}}$ to the source distribution $\delta_\Delta$.
It does not say that the Green kernel itself is a cobar cocycle; the
kernel is the chosen analytic inverse for~$P$, while
$\delta_{D_{12}}$ is the algebraic boundary current in
$d_{\mathrm{extend}}$.

For the Gaussian theory the two-point function is $G(z_1,z_2)$, and
the four-point function is Wick's sum
\[
\langle\phi_1\phi_2\phi_3\phi_4\rangle
=G_{12}G_{34}+G_{13}G_{24}+G_{14}G_{23}.
\]
Interaction vertices, when present, are additional BV/factorization
data matched to the genus-$0$ cobar $A_\infty$ operations by the map
$\Phi$ of Theorem~\ref{thm:cobar-propagator-comparison}.
\end{example}

\begin{remark}[Vertex insertions and cobar products]\label{rem:vertex-operators-cobar}
A CFT vertex operator is not automatically a cobar element. A
comparison datum may assign to a vertex operator $V_\alpha$ a class
$K_\alpha\in H^0(\Omega^{\mathrm{ch}}(\mathcal{C}))$ only after a
BRST-closed insertion map has been constructed. If that map is
compatible with OPE residues and the cobar product, then the singular
part of
\[
V_\alpha(z)V_\beta(w)
\sim
\sum_\gamma
\frac{C_{\alpha\beta}^{\gamma}}
{(z-w)^{h_\alpha+h_\beta-h_\gamma}}V_\gamma(w)
\]
is transported to
\[
n_2(K_\alpha,K_\beta)
=\sum_\gamma C_{\alpha\beta}^{\gamma}K_\gamma.
\]
The equality of the constants is a Ward-identity and normalization
check inside the chosen comparison datum, not a formal consequence of
the cobar construction alone.
\end{remark}

\subsection{Verdier duality on bar and cobar}


\begin{theorem}[Bar-cobar Verdier; \ClaimStatusProvedHere]\label{thm:bar-cobar-verdier}
\label{thm:verdier-bar-cobar}
\index{Verdier duality!bar-cobar|textbf}
\textup{[Regime: quadratic, genus-$0$
\textup{(}Convention~\textup{\ref{conv:regime-tags})}.]}

Let $\mathcal{A}$ be on the Koszul locus and write $\mathcal{A}^!$ for its
Koszul dual.

Write
\[
(\mathcal{A})^!_\infty
:=
\mathbb{D}_{\operatorname{Ran}}(\bar{B}^{\text{ch}}(\mathcal{A}))
\]
for the homotopy Koszul dual factorization algebra. There is a perfect pairing:
\[\langle \cdot, \cdot \rangle: \bar{B}^{\text{ch}}_n(\mathcal{A}) \otimes
(\mathcal{A})^!_{\infty,n} \to \mathbb{C}\]

given by:
\[\langle \omega_{\text{bar}}, K_{\text{Verdier}} \rangle = \int_{\overline{C}_n(X)}
\omega_{\text{bar}} \wedge \iota^* K_{\text{Verdier}}\]

where:
\begin{itemize}
\item $\omega_{\text{bar}} \in \Gamma(\overline{C}_n(X), \mathcal{A}^{\boxtimes n}
\otimes \Omega^*_{\log})$ is a bar element (logarithmic form on compactified space)
\item $K_{\text{Verdier}} \in (\mathcal{A})^!_{\infty,n}$ is represented on $C_n(X)$ by
distributional sections with values in $(\mathcal{A}^!)^{\boxtimes n}$; this is the
post-Verdier factorization-algebra side of Theorem~A, not a pre-Verdier coalgebra
\item $\iota: C_n(X) \hookrightarrow \overline{C}_n(X)$ is the inclusion of the
open configuration space
\item The integration is well-defined because logarithmic forms pair with distributions
\end{itemize}

\emph{Properties of the pairing.}
\begin{enumerate}
\item \emph{Perfect pairing:} Non-degenerate in both arguments by the
canonical holonomic Verdier pairing
$\mathcal M\otimes\mathbb D\mathcal M\to\omega_{\overline C_n(X)}$;
finite-type hypotheses identify the global pairing with the finite
pushforward of this local duality pairing.
\item \emph{Differential compatibility:} $\langle d_{\text{bar}}\omega, K \rangle
= -\langle \omega, d_{\text{cobar}}K \rangle$ (graded Leibniz rule)
\item \emph{Residue-distribution duality:} $\langle \text{Res}_{D}[\omega],
\delta_D \rangle = 1$ for any divisor $D$
\item \emph{Verdier duality:} The perfect pairing realizes
\[
\mathbb{D}_{\operatorname{Ran}}(\bar{B}^{\text{ch}}(\mathcal{A}))
\;\simeq\; (\mathcal{A})^!_\infty
\]
as factorization algebras on $\operatorname{Ran}(X)$; on the Koszul locus,
Theorem~A identifies $(\mathcal{A})^!_\infty \simeq \mathcal{A}^!$
\end{enumerate}
\end{theorem}

\begin{proof}
\emph{Step 1: Well-definedness of the pairing}

Logarithmic forms on $\overline{C}_n(X)$ restrict to distributional
forms on $C_n(X)$. Explicitly, near a divisor $D = \{z_i = z_j\}$ with local
coordinate $\epsilon = z_i - z_j$:

Logarithmic form: $\omega = \frac{d\epsilon}{\epsilon} \wedge (\text{smooth forms})$

Restriction to $C_n(X)$: $\iota^*\omega$ has a pole at $\epsilon = 0$, hence is
a distribution on $C_n(X) = \overline{C}_n(X) \setminus D$.

The pairing integrates this distribution against the corresponding
distributional representative on the Verdier-dual side:
\[\langle \omega, K \rangle = \int_{\overline{C}_n(X)} \omega \wedge K\]

This is well-defined by the theory of currents (de Rham's theorem on distributions).

\emph{Step 2: Differential compatibility}

We verify:
\[\langle d_{\text{bar}}\omega, K \rangle = -\langle \omega, d_{\text{cobar}}K \rangle\]

At the level of $\mathcal{D}$-modules on $\overline{C}_n(X)$, the bar differential and its Verdier-transported differential on the post-Verdier factorization-algebra side are flat connections (the bar complex is a complex of holonomic $\mathcal{D}$-modules by Lemma~\ref{lem:bar-holonomicity}). Verdier duality is an exact contravariant involution on $D^b_{\mathrm{hol}}(\mathcal{D}_{\overline{C}_n(X)})$ that sends a flat connection $\nabla$ to the adjoint connection $\nabla^\dagger$ (see~\cite{KS90}, Chapter~4).

We make the adjoint identity explicit. For a holonomic $\mathcal{D}$-module $\mathcal{M}$ on a smooth variety $Y$ of dimension~$d$, the Verdier dual is $\mathbb{D}\mathcal{M} = R\mathcal{H}\!\mathit{om}_{\mathcal{D}_Y}(\mathcal{M}, \mathcal{D}_Y \otimes \Omega_Y^{-1})[d]$. When $\mathcal{M}$ carries a flat connection $\nabla \colon \mathcal{M} \to \mathcal{M} \otimes \Omega^1_Y$, sections $m^* \in \mathbb{D}\mathcal{M}$ inherit the adjoint connection $\nabla^\dagger$ defined by
\[\langle \nabla^\dagger m^*,\, m \rangle \;=\; -\,\langle m^*,\, \nabla m \rangle\]
for all local sections $m \in \mathcal{M}$, where $\langle -,- \rangle$ is the canonical pairing $\mathbb{D}\mathcal{M} \otimes \mathcal{M} \to \omega_Y$. This is the unique connection on $\mathbb{D}\mathcal{M}$ making the pairing flat, i.e., satisfying $d\langle m^*, m\rangle = \langle \nabla^\dagger m^*, m\rangle + \langle m^*, \nabla m\rangle = 0$.

Applying this to $Y = \overline{C}_n(X)$, $\mathcal{M} = \bar{B}^{\mathrm{ch}}_n(\mathcal{A})$ with $\nabla = d_{\mathrm{bar}}$, and the post-Verdier factorization-algebra side identified with $\mathbb{D}\mathcal{M}$ by the pairing constructed in Steps~3--4, one obtains $d_{\mathrm{cobar}} = d_{\mathrm{bar}}^\dagger$ as the Verdier-adjoint connection. The differential compatibility
\[\langle d_{\text{bar}}\omega, K \rangle = -\langle \omega, d_{\text{cobar}}K \rangle\]
therefore holds by the flatness of the canonical pairing on holonomic $\mathcal{D}$-modules, without requiring a component-by-component Stokes argument.

\emph{Step 3: Residue-current pairing}

Let $D=D_{ij}$ and choose a local normal coordinate
$\epsilon=z_i-z_j$. The residue morphism is the Poincar\'e--Leray map
\[
\operatorname{Res}_{D}\colon
\Omega^1_{\overline C_n(X)}(\log D)\longrightarrow \mathcal O_D,
\qquad
\operatorname{Res}_{D}\!\left(\frac{d\epsilon}{\epsilon}\right)=1.
\]
Its Verdier transpose is the boundary-current inclusion
\[
\iota_{D,!}\omega_D[-1]\longrightarrow
\mathbb D_{\overline C_n(X)}
\Omega^1_{\overline C_n(X)}(\log D).
\]
We denote the image of the unit boundary density by $\delta_D$. With
the standard normalization of the residue map,
\[
\langle\operatorname{Res}_{D}(\eta_{ij}),\delta_D\rangle
=\int_D 1=1
\]
on a normalized local fiber. Thus the notation
$\langle\eta_{ij},\delta(z_i-z_j)\rangle=1$ means
``residue paired with its Verdier-transposed boundary current''. It is
not the product of the singular form $d\epsilon/\epsilon$ with an
ordinary Dirac distribution on the open stratum.

\emph{Step 4: Verdier duality realization}

The pairing establishes an isomorphism of factorization algebras on
$\operatorname{Ran}(X)$:
\[
\mathbb{D}_{\operatorname{Ran}}(\bar{B}^{\text{ch}}(\mathcal{A}))
\xrightarrow{\sim} (\mathcal{A})^!_\infty
\]

where $(\mathcal{A})^!_\infty$ denotes the homotopy Koszul dual factorization
algebra of Theorem~A; on the Koszul locus this is the same post-Verdier
factorization algebra that Theorem~A identifies with $\mathcal{A}^!$.

\emph{Non-degeneracy.}
At bar degree~$n = 1$, Step~3 establishes the fundamental pairing
$\langle \eta_{ij}, \delta(z_i - z_j) \rangle = 1$. For $n > 1$,
non-degeneracy follows by induction using the coalgebra structure: the
comultiplication on $\bar{B}^n$ decomposes an $n$-point integrand into
products of lower-point integrands via the K\"unneth formula on the
FM boundary strata $\partial\overline{C}_n(X) \cong
\coprod_{|S|=k} \overline{C}_k(X) \times \overline{C}_{n-k+1}(X)$.
The induced pairing at degree~$n$ is the tensor product of
pairings at degrees~$k$ and~$n{-}k{+}1$, which are non-degenerate
by the induction hypothesis. Therefore the map is an isomorphism
at each bar degree.

\emph{Geometric meaning.}
\begin{itemize}
\item Bar = logarithmic-form realization of the pre-Verdier bar coalgebra
\item Verdier-dual side = distributional realization of the post-Verdier factorization algebra
\item Pairing = Poincar\'e--Verdier duality between these two realizations
\end{itemize}

This completes the proof.
\end{proof}

\begin{corollary}[Bar-cobar mutual inverses; \ClaimStatusProvedHere]\label{cor:bar-cobar-inverse}
\index{bar-cobar inversion|textbf}
\index{quasi-isomorphism!bar-cobar}
Let $\mathcal{A}$ be an augmented chiral algebra on a smooth curve $X$,
valued in holonomic $\mathcal{D}$-modules
(Convention~\ref{subsec:ambient-category}), and let $\mathcal{C}$ be a
coaugmented chiral coalgebra.

\emph{Algebra side.} If
\begin{enumerate}[label=(A\arabic*)]
\item $\mathcal{A}$ is \emph{chiral Koszul}: the bar complex
 $\bar{B}^{\mathrm{ch}}(\mathcal{A})$ has cohomology concentrated in a
 single degree (Definition~\ref{def:chiral-koszul-pair});
\item \emph{Conilpotence or completion}: either
 $\bar{B}^{\mathrm{ch}}(\mathcal{A})$ is conilpotent, or $\mathcal{A}$
 is complete with respect to its augmentation ideal
 (\S\ref{sec:i-adic-completion}),
\end{enumerate}
then the counit of the adjunction is a quasi-isomorphism:
\[\Omega^{\mathrm{ch}}(\bar{B}^{\mathrm{ch}}(\mathcal{A})) \xrightarrow{\sim} \mathcal{A}.\]

\emph{Coalgebra side.} If, dually,
\begin{enumerate}[label=(C\arabic*)]
\item $\mathcal{C}$ is \emph{coKoszul}: the cobar complex
 $\Omega^{\mathrm{ch}}(\mathcal{C})$ has cohomology concentrated in a
 single degree;
\item $\mathcal{C}$ is conilpotent
 (Definition~\ref{def:conilpotent-cobar}),
\end{enumerate}
then the unit of the adjunction is a quasi-isomorphism:
\[\mathcal{C} \xrightarrow{\sim} \bar{B}^{\mathrm{ch}}(\Omega^{\mathrm{ch}}(\mathcal{C})).\]

The maps are the unit and counit of the bar-cobar adjunction. The
Verdier pairing of Theorem~\ref{thm:bar-cobar-verdier} identifies their
geometric distributional realizations; the quasi-isomorphism assertion
comes from the Koszul collapse plus the conilpotent or completed
convergence hypotheses. When $(\mathcal{A},\mathcal{C})$ form a chiral
Koszul pair with
$\mathcal{C}\simeq\bar{B}^{\mathrm{ch}}(\mathcal{A})$, the two sides
exhibit bar and cobar as quasi-inverse on that Koszul/completed
subcategory, not on all chiral algebras and coalgebras.
\end{corollary}

\begin{proof}
\emph{Algebra side.}
The projection
\[
\pi\colon \bar{B}^{\mathrm{ch}}(\mathcal{A})\longrightarrow
\mathcal{A}
\]
to bar degree~$1$, followed by desuspension and inclusion of
$\bar{\mathcal A}$, is the universal chiral twisting morphism
(Proposition~\ref{prop:universal-twisting-adjunction} and
Definition~\ref{def:chiral-twisting-datum}). The corresponding
twisted tensor products are the Koszul complexes
$K^L_\pi(\mathcal A,\bar B^{\mathrm{ch}}(\mathcal A))$ and
$K^R_\pi(\bar B^{\mathrm{ch}}(\mathcal A),\mathcal A)$. Hypothesis~(A1)
is exactly their acyclicity after passing to the associated quadratic
graded object; hypothesis~(A2) is the convergence input that lets this
associated-graded acyclicity lift to the filtered chiral complexes.
In the completed lane the assertion is checked on the finite quotients
$\mathcal A_{\le N}$, where Lemma~\ref{lem:degree-cutoff} makes every
bar and cobar sum finite, and then passed to the inverse limit by the
complete filtered comparison lemma.

Under the twisting-morphism adjunction, $\pi$ corresponds to the
counit
\[
\varepsilon_{\mathcal A}\colon
\Omega^{\mathrm{ch}}\bar{B}^{\mathrm{ch}}(\mathcal{A})
\longrightarrow \mathcal{A}.
\]
The fundamental theorem of twisting morphisms
\cite[Theorem~2.3.1]{LV12}, applied in the chiral filtered ambient,
identifies acyclicity of the two twisted tensor products with
$\varepsilon_{\mathcal A}$ being a quasi-isomorphism. This proves the
algebra-side claim.

\emph{Coalgebra side.}
The universal inclusion
$\iota\colon\mathcal C\to\Omega^{\mathrm{ch}}(\mathcal C)$ is the
cobar universal twisting morphism. By the same adjunction it
corresponds to the unit
\[
\eta_{\mathcal C}\colon
\mathcal C\longrightarrow
\bar B^{\mathrm{ch}}\Omega^{\mathrm{ch}}(\mathcal C).
\]
Conilpotence gives an exhaustive finite coradical filtration, so the
coalgebra-side twisted tensor products are defined without hidden
infinite sums. The coKoszul hypothesis is the dual acyclicity
condition. The same fundamental twisting theorem therefore makes
$\eta_{\mathcal C}$ a quasi-isomorphism.
\end{proof}

\begin{proposition}[Explicit cobar-bar augmentation; \ClaimStatusProvedHere]
\label{prop:cobar-bar-augmentation}
\index{augmentation map!cobar-bar|textbf}
\index{quasi-isomorphism!explicit chain map}
Let $\mathcal{A}$ be an augmented chiral algebra satisfying the
algebra-side hypotheses (A1)--(A2) of
Corollary~\ref{cor:bar-cobar-inverse}.
The \emph{augmentation map}
$\varepsilon\colon
\Omega^{\mathrm{ch}}(\bar{B}^{\mathrm{ch}}(\mathcal{A}))
\to \mathcal{A}$
is constructed as follows.

\emph{On generators.}
By Theorem~\ref{thm:cobar-free},
$\Omega^{\mathrm{ch}}(\bar{B}^{\mathrm{ch}}(\mathcal{A}))
= \mathrm{Free}_{\mathrm{ch}}(s^{-1}\overline{B})$
where $\overline{B} = \bigoplus_{n \geq 1}
\bar{B}^{\mathrm{ch}}_n(\mathcal{A})$ is the coaugmentation
coideal of the bar coalgebra.
Let
\[
\tau_{\mathcal A}\colon s^{-1}\overline B\longrightarrow \mathcal A
\]
be the generator-level representative of the canonical chiral
twisting morphism
$\pi_{\mathcal A}\colon\bar B^{\mathrm{ch}}(\mathcal A)\to\mathcal A$
under the free-algebra adjunction. Define
\begin{equation}\label{eq:augmentation-on-generators}
\varepsilon|_{s^{-1}\overline B}:=\tau_{\mathcal A}.
\end{equation}
Equivalently, $\tau_{\mathcal A}$ is the bar-degree-one projection
followed by the fixed desuspension convention of
Appendix~\ref{app:signs}; it is zero on
$s^{-1}\bar B^{\mathrm{ch}}_n(\mathcal A)$ for $n\geq2$.

\emph{Extension.}
By the universal property of the free chiral algebra
(Theorem~\ref{thm:cobar-free}), the restriction
\eqref{eq:augmentation-on-generators} extends uniquely to a
chiral algebra morphism
$\varepsilon\colon
\mathrm{Free}_{\mathrm{ch}}(s^{-1}\overline{B}) \to \mathcal{A}$.

\emph{Low-degree formulas.}
On generators:
\begin{align}
\varepsilon(s^{-1}b_{\bar a}) &= \bar{a}
 &&\text{(inclusion of the augmentation ideal),}
 \label{eq:aug-degree-1} \\
\varepsilon(s^{-1}b_{\bar a_1,\bar a_2}) &= 0
 &&\text{(projection kills bar degree $\geq 2$).}
 \label{eq:aug-degree-2-gen}
\end{align}
Here $b_{\bar a}\in\bar B^{\mathrm{ch}}_1(\mathcal A)$ is the
bar-degree-one generator corresponding to
$\bar a\in\bar{\mathcal A}$, and
$b_{\bar a_1,\bar a_2}\in\bar B^{\mathrm{ch}}_2(\mathcal A)$ is the
bar-degree-two generator represented by the two insertions
$\bar a_1,\bar a_2$ in the symmetric factorization bar.
On products of generators (by the algebra map property):
\begin{equation}\label{eq:aug-product}
\varepsilon\bigl(s^{-1}b_{\bar a_1}
 \cdot s^{-1}b_{\bar a_2}\bigr)
= \bar{a}_{1\,(0)}\bar{a}_2
\qquad\text{(the zeroth chiral product).}
\end{equation}
The element $s^{-1}b_{\bar a_1,\bar a_2}$ is a \emph{single
generator} from bar degree~$2$; the element
$s^{-1}b_{\bar a_1}\cdot s^{-1}b_{\bar a_2}$ is a
\emph{product} of two bar-degree-one generators, which
$\varepsilon$ evaluates via the chiral multiplication
$\bar{a}_{1\,(0)}\bar{a}_2$.
\end{proposition}

\begin{proof}
\emph{Step~1: Chain map property.}
We verify
\[
d_{\mathcal A}\varepsilon=\varepsilon d_\Omega .
\]
Since $\varepsilon$ is an algebra morphism and $d_\Omega$ is a derivation
(Theorem~\ref{thm:cobar-free}), it suffices to check on generators
$s^{-1} b$ for $b \in \bar{B}^{\mathrm{ch}}_n(\mathcal{A})$.

The cobar differential on generators decomposes as
(Theorem~\ref{thm:cobar-diff-geom}):
\[
d_\Omega(s^{-1} b)
\;=\; \pm\, s^{-1}(d_{\mathrm{bar}}\, b)
\;+\; \sum \pm\, s^{-1} b' \cdot s^{-1} b''
\]
where $d_{\mathrm{bar}}$ is the bar differential and
$\bar{\Delta}(b) = \sum b' \otimes b''$ is the reduced coproduct
(signs determined by the Koszul rule for desuspension).

\emph{Bar degree $n=1$:}
$b = b_{\bar a} \in \bar{B}^{\mathrm{ch}}_1(\mathcal A)$.
The reduced coproduct vanishes and the internal part of the bar
differential corresponds, under the same desuspension convention, to
$d_{\mathcal A}\bar a$. Hence
\[
\varepsilon d_\Omega(s^{-1}b_{\bar a})
=d_{\mathcal A}\bar a
=d_{\mathcal A}\varepsilon(s^{-1}b_{\bar a}).
\]
For a strict algebra with zero internal differential this specializes
to the displayed zero computation.

\emph{Bar degree $n=2$:}
$b = b_{\bar a_1,\bar a_2} \in
\bar{B}^{\mathrm{ch}}_2(\mathcal A)$.
The bar differential has a bar-degree-one collision residue term
corresponding to $\bar a_{1\,(0)}\bar a_2$. The reduced coshuffle
coproduct has singleton--singleton summands
$b'\otimes b''$ corresponding to the two bar-degree-one generators.
After applying $\varepsilon$, the collision contribution and the
product contribution are exactly the two terms of the Maurer--Cartan
equation
\[
d_{\mathcal A}\pi_{\mathcal A}
\;+\;\pi_{\mathcal A}d_{\bar B}
\;+\;\pi_{\mathcal A}\smile\pi_{\mathcal A}=0
\]
for the universal twisting morphism
\textup{(}\cite[\S2.2, Proposition~2.2.6]{LV12}\textup{)}. Hence
$\varepsilon d_\Omega(s^{-1}b)=0=d_{\mathcal A}\varepsilon(s^{-1}b)$
in bar degree~$2$.

\emph{Bar degree $n \geq 3$:}
$d_{\mathrm{bar}}\, b$ has an internal part in bar degree~$n$ and
collision parts in bar degree~$n{-}1\geq2$, all killed by
$\varepsilon$. Each summand $s^{-1}b'\cdot s^{-1}b''$ in the reduced
coproduct has at least one factor of bar degree~$\geq2$, except for
the adjacent binary faces already treated in bar degree~$2$ after
applying the MC equation locally. Thus
$\varepsilon d_\Omega(s^{-1}b)=0=d_{\mathcal A}\varepsilon(s^{-1}b)$.

\emph{Step~2: Quasi-isomorphism.}
The map $\varepsilon$ is the adjunction counit corresponding to the
universal twisting morphism~$\pi$. Under the hypotheses of
Corollary~\ref{cor:bar-cobar-inverse}, the twisted tensor products for
$\pi$ are acyclic and the conilpotent or completed filtration
converges. The fundamental theorem of twisting morphisms therefore
implies that this explicit counit is a quasi-isomorphism. This is the
same map denoted $\psi_{\mathcal A}$ in
Theorem~\ref{thm:bar-cobar-adjunction}.
\end{proof}

\begin{remark}[The chiral bar-cobar adjunction via twisting morphisms]
\label{rem:chiral-adjunction-via-tw}
\index{bar-cobar adjunction!via twisting morphisms!chiral}
\index{twisting morphism!chiral universal}
The restriction
$\varepsilon|_{s^{-1}\overline{B}}\colon
s^{-1}\overline B\to\cA$
is the generator-level representative of the chiral universal twisting
morphism
$\pi\colon\barB^{\mathrm{ch}}(\cA) \to \cA$
(Proposition~\ref{prop:universal-twisting-adjunction}(i),
Definition~\ref{def:chiral-twisting-datum}).
By the two universal properties
(Proposition~\ref{prop:universal-twisting-adjunction}(i)--(ii)),
the \emph{chiral bar-cobar adjunction} is:
\begin{equation}\label{eq:chiral-tw-adjunction}
\mathrm{Hom}_{\mathrm{dg\,CoAlg}}(\mathcal{C},\,
 \barB^{\mathrm{ch}}(\cA))
\;\cong\;
\mathrm{Tw}^{\mathrm{ch}}(\mathcal{C},\, \cA)
\;\cong\;
\mathrm{Hom}_{\mathrm{dg\,Alg}}(\Omega^{\mathrm{ch}}(\mathcal{C}),\,
 \cA).
\end{equation}
On the Koszul locus,
$\barB^{\mathrm{ch}}(\cA) \otimes_\pi \cA \simeq \mathbb{k}$
(Theorem~\ref{thm:fundamental-twisting-morphisms}): unit and counit
are quasi-isomorphisms.
\end{remark}

\begin{example}[Augmentation for \texorpdfstring{$\widehat{\mathfrak{sl}}_2$}{sl2-hat}]
\label{ex:augmentation-sl2}
For the affine Kac--Moody algebra
$\mathcal{A} = V_k(\mathfrak{sl}_2)$ with basis
$\{J^+, J^0, J^-\}$ and OPE
$J^a(z) J^b(w) \sim f^{ab}_c J^c(w)/(z-w)
+ k\,(J^a, J^b)_{\mathrm{tr}}/(z-w)^2$,
the augmentation acts as:
\begin{align*}
\varepsilon(s^{-1}b_{J^a}) &= J^a
 &&\text{(recovery of generators),} \\
\varepsilon(s^{-1}b_{J^a} \cdot s^{-1}b_{J^b})
&= J^a_{(0)}J^b = f^{ab}_c J^c
 &&\text{(recovery of structure constants).}
\end{align*}
Here $b_{J^a}\in\bar B^{\mathrm{ch}}_1(\mathcal A)$ denotes the
bar-degree-one generator corresponding to~$J^a$. Thus
$\varepsilon(s^{-1}b_{J^+} \cdot s^{-1}b_{J^-}) = J^0$
(= $[J^+, J^-] = H$),
confirming that the augmentation recovers the
$\mathfrak{sl}_2$ Lie bracket from the cobar-bar resolution.
\end{example}

\begin{example}[Explicit pairing: two-point function]\label{ex:pairing-two-point}

Consider $n=2$. The bar element is:
\[\omega_{\text{bar}} = a_1(z_1) \otimes a_2(z_2) \otimes \frac{dz_1 - dz_2}{z_1 - z_2}\]

The cobar element is:
\[K_{\text{cobar}} = c_1(z_1) \otimes c_2(z_2) \otimes \delta(z_1 - z_2)\]

The pairing:
\begin{align*}
\langle \omega_{\text{bar}}, K_{\text{cobar}} \rangle &= \int_{\overline{C}_2(X)}
(a_1 \otimes a_2) \cdot (c_1 \otimes c_2) \wedge \frac{dz_1 - dz_2}{z_1 - z_2}
\wedge \delta(z_1 - z_2) \\
&= \int_X (a_1 \otimes a_2)(z, z) \cdot (c_1 \otimes c_2)(z, z) \wedge dz \wedge d\bar{z}
\end{align*}

By the residue-distribution identity:
\[\int \frac{dz_1 - dz_2}{z_1 - z_2} \wedge \delta(z_1 - z_2) = 1\]

Therefore:
\[\langle \omega_{\text{bar}}, K_{\text{cobar}} \rangle = \int_X \langle a_1, c_1
\rangle \cdot \langle a_2, c_2 \rangle \, dz \wedge d\bar{z}\]

Under a CFT comparison datum this pairing is the two-point kernel
pairing; algebraically it is the degree-two Verdier pairing.
\end{example}

\subsection{Configuration-integral comparison}

\begin{theorem}[Kontsevich formality (1997) {\cite{Kon99}}; \ClaimStatusProvedElsewhere]\label{thm:kontsevich-formality}
For any smooth manifold $M$, there exists an $L_\infty$
quasi-isomorphism:
\[\mathcal{U}: T_{\text{poly}}(M) \xrightarrow{\sim} D_{\text{poly}}(M)\]
from polyvector fields to polydifferential operators, given by configuration space integrals:
\[\mathcal{U}_n(\gamma_1, \ldots, \gamma_n) = \sum_{\Gamma \in G_n} w_\Gamma
\int_{\overline{C}_{n,m}(\mathbb{H})} \omega_\Gamma\]
where $G_n$ = admissible graphs, $w_\Gamma$ = combinatorial weights, and
$\omega_\Gamma$ involves propagators $d\log(z_i - z_j)$ and angle forms $d\theta_i$.
\end{theorem}

\begin{remark}[Configuration-integral comparison]\label{rem:kontsevich-chiral}
Kontsevich formality and the chiral bar construction both use
compactified configuration spaces and logarithmic angle forms. The
comparison stops there unless the source and target operads are fixed:
formality compares polyvector fields with polydifferential operators,
whereas the present chapter compares the chiral bar coalgebra with its
cobar inverse and, separately, the Verdier-dual Koszul algebra.
\end{remark}

\begin{remark}[Costello--Gwilliam factorization algebras]\label{rem:CG-factorization-detailed}
Our construction extends~\cite{CG17}: Vol.~1 develops factorization homology on manifolds, Vol.~2 treats quantum corrections and curved $A_\infty$ in QFT. We specialize to complex curves (essential for chiral structure) and use configuration geometry rather than the BV formalism.
\end{remark}

\begin{remark}[Cobar summary]
The intrinsic cobar differential is obtained by Verdier transport of
the bar differential of the dual algebra. Thus, in the uncurved lane,
the three-component differential, the $d^2 = 0$ verification, and the
Arnold relations mirror their bar counterparts; in the curved lane the
same transport gives the inner-curvature square
$d_{\Omega,h}^{\,2}=[h_\Omega,-]$.
The functor $\Omega^{\mathrm{ch}}$ remains the bar-cobar inverse, not
Verdier duality on $\barBch(\cA)$.
\end{remark}

\begin{proposition}[Finite shadow extraction from the cobar side]
\label{prop:cobar-modular-shadow}
\ClaimStatusConditional
Let $\cA$ lie in the finite-window or completed pro-conilpotent lane in
which Theorem~\ref{thm:bar-modular-operad} gives the bar coalgebra an
$\mathrm{FCom}$-algebra structure and the counit
$\Omega(\barB(\cA))\to\cA$ is defined. For each finite arity and genus
cutoff $r$, the induced cobar-side structure maps determine the
projection $\Theta_\cA^{\le r}$ of the bar-intrinsic MC element after
passing through the counit and the finite quotient defining that
cutoff. This is a finite extraction statement; it is not an assertion
that an unconverged infinite graph sum acts on the raw direct-sum
cobar object.
\end{proposition}

\begin{proof}
The bar-intrinsic element
$\Theta_\cA:=D_\cA-d_0$ lives in the completed modular convolution
algebra of Theorem~\ref{thm:mc2-bar-intrinsic}. The
$\mathrm{FCom}$-structure of Theorem~\ref{thm:bar-modular-operad}
expresses each finite cutoff as a signed sum over oriented modular
graphs:
\[
\sum_{\Gamma\in\mathcal G_{\le r}}
\frac{1}{|\operatorname{Aut}(\Gamma)|}\,\ell_\Gamma .
\]
The orientation signs, automorphism denominators, and stability
conditions are part of the modular-operadic structure map. Applying
the cobar counit transports this finite graph operator to the same
finite quotient of the convolution algebra. Corollary~\ref{cor:shadow-extraction}
identifies that quotient with the projection
$\Theta_\cA^{\le r}$. Completion is used only to pass from the
compatible family of finite extractions to the full tower.
\end{proof}

\subsection{Čech-Alexander complex realization}

\begin{remark}[Lurie comparison]\label{rem:v1-lurie-comparison-bar-cobar}
The comparison with Lurie's bar-cobar formalism has four layers.
\begin{enumerate}[label=\textup{(\roman*)}]
\item In Higher Algebra, the bar construction of an augmented algebra is the simplicial bar object, equivalently the \v{C}ech nerve of the augmentation
\textup{(}Lurie~\cite[Example~5.2.2.3, Proposition~5.2.2.5]{HA}\textup{)}.
Our reduced bar uses the same augmentation pattern, with
$\bar{\cA} = \ker(\varepsilon)$ and the Fulton--MacPherson
compactifications $\overline{C}_{n+1}(X)$ replacing the ordinary simplices.

\item Lurie identifies the bar construction with the Koszul-dual
coalgebra in a stable monoidal setting
\textup{(}Lurie~\cite[Proposition~5.2.5.1]{HA}\textup{)}.
This matches Theorem~A only after the distinction of
Remark~\ref{rem:cobar-three-functors}: Verdier duality on
$\operatorname{Ran}(X)$ produces the homotopy Koszul-dual
factorization algebra, whereas the cobar functor produces the inverse
$\Omega(\barB(\cA)) \simeq \cA$ on the Koszul locus.

\item The Ran-space descent used here is the chiral analogue of
Lurie's description of factorizable cosheaves on
$\operatorname{Ran}(M)$
\textup{(}Lurie~\cite[Theorems~5.5.4.10 and~5.5.4.14]{HA}\textup{)}.
For algebraic curves with $\cD$-modules, the direct comparison is
Francis--Gaitsgory's chiral Koszul duality on $\operatorname{Ran}(X)$
\textup{(}\cite[Remark~2.5.7, Theorem~5.1.1]{FG12}\textup{)}.
This is the piece of the Gaitsgory--Rozenblyum programme closest to
the chiral cobar construction: the prestack and correspondence background is
nearer to Volume~I than to Volume~II, while Volume~II enters later
through formal geometry and derived $\cD$-modules.

\item The coderived continuation off the Koszul locus is not supplied
by Higher Algebra alone. Lurie's stable $\infty$-categories provide
the derived ambient category, but the curved and completed
continuation used in
Theorem~\ref{thm:higher-genus-inversion}\textup{(}b\textup{)}
is Positselski's coderived enlargement
\textup{(}\cite{Positselski11}\textup{)}.
\end{enumerate}

\noindent
The truthful dictionary is therefore: Higher Algebra explains the
augmentation nerve and the operadic bar-cobar package; Francis--Gaitsgory
supplies the chiral Ran-space realization; Positselski controls the
completed coderived continuation. We do not claim here a separate
proved-here theorem identifying $\Omega^{\mathrm{ch}}(\cC)$ with a
coordinate \v{C}ech complex; the rigorous comparison surface in this
chapter is Verdier duality plus factorization descent.
\end{remark}

\subsection{Integration kernels and cobar operations}

\begin{definition}[Cobar integration kernel]\label{def:cobar-kernel}
Elements of the cobar complex can be represented by integration kernels:
\[
K_{p+1}(z_0, \ldots, z_p; w_0, \ldots, w_p) \in \Gamma\left(C_{p+1}(X) \times C_{p+1}(X), \text{Hom}(\mathcal{C}^{\otimes(p+1)}, \mathbb{C}) \otimes \Omega^*\right)
\]
acting on sections of $\mathcal{C}$ by:
\[
(\Phi_K \cdot c)(z_0, \ldots, z_p) = \int_{C_{p+1}(X)} K_{p+1}(z_0, \ldots, z_p; w_0, \ldots, w_p) \wedge c(w_0) \otimes \cdots \otimes c(w_p)
\]
\end{definition}

\begin{remark}[No universal scalar cobar kernel]\label{ex:fundamental-cobar}
For the trivial chiral coalgebra $\mathcal{C}=\omega_X$, the
arity-two cobar current is the FM boundary-current class dual to the
residue on $D_{12}$. There is no canonical rational scalar kernel on
$C_2(X)\times C_2(X)$ that reconstructs a nontrivial chiral
multiplication from $\omega_X$ alone. Any displayed Green or Cauchy
kernel requires extra data: a coordinate, a density, and usually a
parametrix. The intrinsic cobar object supplies the unit/counit
boundary current; a QFT or analytic comparison datum supplies
propagator kernels.
\end{remark}

\begin{theorem}[Cobar as free chiral algebra; \ClaimStatusProvedHere]\label{thm:cobar-free}
As a graded chiral algebra (forgetting the differential), the cobar
construction $\Omega^{\mathrm{ch}}(\mathcal{C})$ is the free chiral
algebra generated by $s^{-1}\bar{\mathcal{C}}$, where
$\bar{\mathcal{C}} = \ker(\epsilon\colon \mathcal{C} \to \omega_X)$.
The differential $d_\Omega$ is the unique derivation extending the
reduced comultiplication
$\bar{\Delta}\colon \bar{\mathcal{C}} \to
\bar{\mathcal{C}} \otimes \bar{\mathcal{C}}$
(this is essential data: $\Omega^{\mathrm{ch}}(\mathcal{C})$ is a
\emph{DG} chiral algebra whose differential is part of the structure).
\end{theorem}

\begin{proof}
By construction (Definition~\ref{def:geom-cobar-precise}), the underlying graded chiral algebra of $\Omega^{\mathrm{ch}}(\mathcal{C})$ is $\mathrm{Free}_{\mathrm{ch}}(s^{-1}\bar{\mathcal{C}})$, the free chiral algebra generated by $s^{-1}\bar{\mathcal{C}}$ as a $\mathcal{D}_X$-module. The universal property of free algebras then gives: for any chiral algebra~$\mathcal{A}$ and graded $\mathcal{D}_X$-module morphism $f\colon s^{-1}\bar{\mathcal{C}} \to \mathcal{A}$, there exists a unique chiral algebra morphism $\tilde{f}\colon \mathrm{Free}_{\mathrm{ch}}(s^{-1}\bar{\mathcal{C}}) \to \mathcal{A}$ extending~$f$.

It remains to verify that $d_\Omega$ is a derivation. Since $\mathrm{Free}_{\mathrm{ch}}$ is left adjoint to the forgetful functor, a derivation on a free algebra is determined by its restriction to generators. The restriction $d_\Omega|_{s^{-1}\bar{\mathcal{C}}}$ equals the desuspended reduced comultiplication $s^{-1}\bar{\Delta}\colon s^{-1}\bar{\mathcal{C}} \to (s^{-1}\bar{\mathcal{C}})^{\otimes 2} \hookrightarrow \mathrm{Free}_{\mathrm{ch}}(s^{-1}\bar{\mathcal{C}})$, which is a well-defined map of $\mathcal{D}_X$-modules. The Leibniz extension to all of $\Omega^{\mathrm{ch}}(\mathcal{C})$ is unique, giving $d_\Omega$ as the unique derivation extending $\bar{\Delta}$.
\end{proof}

\subsection{Geometric bar-cobar composition}

\begin{theorem}[Chiral bar-cobar adjunction backbone, geometric form;
\ClaimStatusConditional]
\label{thm:geom-unit}%
\label{thm:bar-cobar-adjunction}%
\index{bar-cobar adjunction|textbf}%
\index{Theorem A!chiral bar-cobar adjunction backbone}%
\textup{[Regime: any augmented chiral algebra; no Koszul hypothesis required for the
adjunction. Quasi-isomorphism is the content of
Theorem~B, not of this
theorem; cf.\ Remark~\ref{rem:eta-adjunction-vs-qi-scope},
Remark~\ref{rem:fundamental-distinction}.]}

Let $\cA$ be an augmented chiral algebra on a smooth curve~$X$, valued in
holonomic $\cD_X$-modules
\textup{(}Convention~\ref{subsec:ambient-category}\textup{)}, and write
$\bar{\cA} = \ker(\varepsilon\colon \cA \to \omega_X)$.

\begin{enumerate}[label=\textup{(\arabic*)}]
\item \emph{Existence of the counit.} There is a canonical morphism of dg
chiral algebras
\[
\psi_{\cA} \;\colon\; \Omega^{\mathrm{ch}}(\bar{B}^{\mathrm{ch}}(\cA))
 \;\longrightarrow\; \cA,
\]
the \emph{counit of the bar-cobar adjunction
$\Omega^{\mathrm{ch}} \dashv \bar{B}^{\mathrm{ch}}$}
\textup{(}Remark~\ref{rem:adjunction-direction-convention}\textup{)},
characterised by the universal property
\textup{(}Proposition~\ref{prop:universal-twisting-adjunction}(ii)\textup{)}:
$\psi_{\cA}$ is the unique chiral-algebra extension of the chiral universal
twisting morphism
$\pi \colon \bar{B}^{\mathrm{ch}}(\cA) \to \cA$ along the desuspension
inclusion
$\iota \colon \bar{B}^{\mathrm{ch}}(\cA) \xrightarrow{s^{-1}}
s^{-1}\overline{\bar{B}^{\mathrm{ch}}(\cA)}
\hookrightarrow \Omega^{\mathrm{ch}}(\bar{B}^{\mathrm{ch}}(\cA))$.
Concretely, $\psi_{\cA}$ is the augmentation map of
Proposition~\ref{prop:cobar-bar-augmentation}.

\item \emph{Geometric Fulton--MacPherson realisation.} On the
configuration-space side, $\psi_{\cA}$ is realised on each tensor power by
the residue pairing
\begin{equation}\label{eq:psi-FM-realisation}
\psi_{\cA}\bigl(s^{-1}b_1 \cdots s^{-1}b_{n+1}\bigr)(z)
\;=\;
\int_{\overline{C}_{n+1}(X)}
 \mathrm{Res}_{\mathrm{ev}_0 = z}\,(b_1 \boxtimes \cdots \boxtimes b_{n+1})
 \;\wedge\; \omega_n,
\end{equation}
where $\mathrm{ev}_0 \colon \overline{C}_{n+1}(X) \to X$ evaluates at the
zeroth point, $\omega_n$ is the Poincar\'e--Verdier dual of the small
diagonal $X \hookrightarrow X^{n+1}$, and the residue is taken along the
divisor where the remaining $n$ points collide with~$z$. The sum
$\psi_{\cA} = \sum_{n\geq 0} \psi_{\cA,n}$ converges as a section of
$\cA$ when $\bar{B}^{\mathrm{ch}}(\cA)$ is conilpotent or $\cA$ is complete
with respect to the augmentation ideal
\textup{(}Hypothesis (H1)+(H2) of
Theorem~\ref{thm:bar-cobar-platonic}\textup{)}: each weight piece of $\cA$
receives contributions from only finitely many bar degrees.

\item \emph{Naturality.} For every morphism $f \colon \cA \to \cA'$ of
augmented chiral algebras, the diagram
\[
\begin{tikzcd}[ampersand replacement=\&]
\Omega^{\mathrm{ch}}(\bar{B}^{\mathrm{ch}}(\cA))
 \arrow[r, "\psi_{\cA}"]
 \arrow[d, "\Omega^{\mathrm{ch}}(\bar{B}^{\mathrm{ch}}(f))"'] \&
\cA \arrow[d, "f"] \\
\Omega^{\mathrm{ch}}(\bar{B}^{\mathrm{ch}}(\cA'))
 \arrow[r, "\psi_{\cA'}"'] \&
\cA'
\end{tikzcd}
\]
commutes; that is, $\psi$ is a natural transformation
$\Omega^{\mathrm{ch}} \circ \bar{B}^{\mathrm{ch}} \Rightarrow \mathrm{id}$
on $\mathrm{Alg}^{\mathrm{ch,aug}}_X$.

\item \emph{Adjunction Hom-bijection.} For every coaugmented conilpotent
chiral coalgebra $\cC$ on~$X$, the counit~$\psi_{\cA}$ together with the
unit
$\eta_{\cC} \colon \cC \to \bar{B}^{\mathrm{ch}}\Omega^{\mathrm{ch}}(\cC)$
\textup{(}coalgebra side, by the dual of
Lemma~\ref{lem:cobar-derivation-extension}\textup{)} implement the natural
bijection
\begin{equation}\label{eq:psi-eta-hom-bijection}
\Hom_{\mathrm{dg\,Alg}^{\mathrm{ch}}_X}\bigl(\Omega^{\mathrm{ch}}(\cC), \cA\bigr)
\;\cong\;
\mathrm{Tw}^{\mathrm{ch}}(\cC, \cA)
\;\cong\;
\Hom_{\mathrm{dg\,CoAlg}^{\mathrm{ch,conil}}_X}\bigl(\cC, \bar{B}^{\mathrm{ch}}(\cA)\bigr),
\end{equation}
which is the chiral specialisation of LV12 Theorem~2.2.9 and which is the
content of the bar-cobar \emph{adjunction}
$\Omega^{\mathrm{ch}} \dashv \bar{B}^{\mathrm{ch}}$ \textup{(}Theorem~A of
the programme\textup{)}.
The counit $\psi_{\cA}$ is a quasi-isomorphism under the
Koszul/collapse hypotheses of Theorem~B; off that strict lane the
conclusion of (1)--(4) still holds, but $\psi_{\cA}$ is only an
isomorphism in the coderived category under the off-Koszul filtered
pro-conilpotent hypotheses.
\end{enumerate}
\end{theorem}

\begin{proof}[Geometric proof]
\emph{Step 1: existence and chiral-algebra-map property.}
By Theorem~\ref{thm:cobar-free}, the underlying graded chiral algebra of
$\Omega^{\mathrm{ch}}(\bar{B}^{\mathrm{ch}}(\cA))$ is the free chiral algebra
$\mathrm{Free}_{\mathrm{ch}}(s^{-1}\overline{B})$ on the desuspended reduced
bar coalgebra. The generator-level map
$s^{-1}\overline B\to\cA$ attached to the canonical twisting morphism
$\pi_{\cA}\colon\bar B^{\mathrm{ch}}(\cA)\to\cA$,
together with the universal property of the free chiral algebra
\textup{(}Lemma~\ref{lem:cobar-derivation-extension}\,(i)\textup{)} produces
a unique chiral-algebra morphism extending that generator-level map. This is exactly the
augmentation~$\varepsilon$ of
Proposition~\ref{prop:cobar-bar-augmentation}, which we rename
$\psi_{\cA}$ here to align with the canonical adjunction-counit notation of
Remark~\ref{rem:adjunction-direction-convention}.

\emph{Step 2: chain-map property.}
By Lemma~\ref{lem:cobar-derivation-extension}\,(ii), to verify
$\psi_{\cA} \circ d_\Omega = d_{\cA} \circ \psi_{\cA}$ it suffices to check
the identity on generators. This is the content of Step~1 of the proof of
Proposition~\ref{prop:cobar-bar-augmentation}: bar degree~$1$ is the
internal-differential compatibility of the twisting morphism; bar
degree~$2$ uses the Maurer--Cartan equation for the chiral universal
twisting morphism
$\pi_{\cA}$, whose generator-level representative is
$\psi_{\cA}|_{s^{-1}\overline{B}}$
\textup{(}\cite[\S2.2, Proposition~2.2.6]{LV12}\textup{)}; bar degrees
$n \geq 3$ vanish because $\psi_{\cA}$ kills generators from bar
degree $\geq 2$.

\emph{Step 3: geometric Fulton--MacPherson formula.}
Each tensor power $\Omega^{\mathrm{ch}}_{(n+1)}(\bar{B}^{\mathrm{ch}}(\cA))$
is the chiral product on $\overline{C}_{n+1}(X)$ of $n+1$ desuspended bar
generators
\textup{(}Definition~\ref{def:geom-cobar-precise}\textup{)}; the
small-diagonal class $\omega_n$ is its Poincar\'e--Verdier dual under the
configuration-stratum residue pairing
\textup{(}Theorem~\ref{thm:bar-cobar-verdier},
Lemma~\ref{lem:bar-holonomicity}\textup{)}.
The formula~\eqref{eq:psi-FM-realisation} matches the algebraic $\psi_{\cA}$
of Step~1 termwise on each FM stratum because the residue of the bar
generator on the boundary stratum $\partial_{ij}\overline{C}_{n+1}(X)$
equals the chiral OPE of $\cA$ at the merged point; this is the chiral
incarnation of the residue-distribution duality
\textup{(}Theorem~\ref{thm:bar-cobar-verdier} Step~3\textup{)}, applied
inductively along the FM degeneration tower as in
Proposition~\ref{prop:cobar-bar-augmentation} Step~1.

\emph{Step 4: convergence.}
Conilpotency of $\bar{B}^{\mathrm{ch}}(\cA)$
\textup{(}automatic under (H1)+(H2) by
Theorem~\ref{thm:coalgebra-via-NAP}(4) and the conformal-weight grading;
cf.\ the geometric inputs in Chapter~\ref{chap:bar-cobar-adjunction}
preceding Theorem~\ref{thm:bar-cobar-inversion-qi}\textup{)} ensures that
each weight piece of $\Omega^{\mathrm{ch}}(\bar{B}^{\mathrm{ch}}(\cA))$ has
only finitely many bar-degree contributions, so the sum
$\psi_{\cA} = \sum_n \psi_{\cA,n}$ stabilises in each weight piece.
Equivalently, $I$-adic completeness of $\cA$
\textup{(}Hypothesis (H2)\textup{)} produces the same finite-stage
stabilisation by the $I$-adic-vs-conilpotent equivalence of
\S\ref{sec:i-adic-completion}.

\emph{Step 5: naturality.}
Both $\bar{B}^{\mathrm{ch}}$ and $\Omega^{\mathrm{ch}}$ are functorial on
$\mathrm{Alg}^{\mathrm{ch,aug}}_X$ and $\mathrm{CoAlg}^{\mathrm{ch,conil}}_X$
respectively
\textup{(}functoriality of the underlying tensor coalgebra/algebra by
Lemma~\ref{lem:cobar-derivation-extension}, naturality of the differential
by the same generator-restriction principle\textup{)}. The square in (3) is
the generator-restriction comparison: both composites agree on
$s^{-1}\bar{B}^{\mathrm{ch}}_1(\cA) \cong \bar{\cA}$ via the morphism~$f$,
hence on all of $\Omega^{\mathrm{ch}}(\bar{B}^{\mathrm{ch}}(\cA))$ by
Lemma~\ref{lem:cobar-derivation-extension}\,(ii).

\emph{Step 6: Hom-bijection.}
This is Proposition~\ref{prop:universal-twisting-adjunction}(iii) applied to
the chiral universal twisting morphism: the bijections
$\psi_\cA \mapsto \pi$ and $\eta_\cC \mapsto \pi$ exhibit
$\Omega^{\mathrm{ch}} \dashv \bar{B}^{\mathrm{ch}}$ via the universal
twisting bifunctor $\mathrm{Tw}^{\mathrm{ch}}(-,-)$, and the chiral
specialisation works verbatim because $\bar{B}^{\mathrm{ch}}$ is computed in
the symmetric monoidal $(\infty,1)$-category of factorisation $\cD$-modules
on $\Ran(X)$
\textup{(}Theorem~\ref{thm:bar-cobar-platonic}\textup{)} where the
convolution dGLA is well-defined in each weight piece by
Hypothesis~(H3).

\emph{Step 7: scope of qi.}
Whether $\psi_{\cA}$ is a quasi-isomorphism is governed by
Theorem~\ref{thm:bar-cobar-inversion-qi}: on the strict Koszul lane it is a
genus-graded qi; off the Koszul locus it is the cone of an
$\mathrm{IndCoh}(\Ran(X))$-coderived isomorphism whose promotion to a
chain-level qi requires the collapse input of clause~(4) of
Theorem~\ref{thm:bar-cobar-inversion-qi} or the flat case
$\kappa(\cA) = 0$. The present theorem makes \emph{no} qi claim for
$\psi_{\cA}$ off the Koszul locus; see
Remark~\ref{rem:eta-adjunction-vs-qi-scope} below.
\end{proof}

\begin{lemma}[Cobar derivation extension; \ClaimStatusProvedHere]
\label{lem:cobar-derivation-extension}
Let $\cC$ be an uncurved coaugmented chiral coalgebra on~$X$ with reduced
comultiplication
$\bar{\Delta} \colon \bar{\cC} \to \bar{\cC} \otimes \bar{\cC}$,
and write $V := s^{-1}\bar{\cC}$ for the desuspension. Let
$\Omega^{\mathrm{ch}}(\cC) = \mathrm{Free}_{\mathrm{ch}}(V)$ be the free
chiral algebra on~$V$.
\begin{enumerate}[label=\textup{(\roman*)}]
\item \emph{Free extension property.} For every chiral algebra~$\cA$ and every
$\cD_X$-module morphism $f \colon V \to \cA$, there is a unique
chiral-algebra morphism
$\widetilde{f} \colon \Omega^{\mathrm{ch}}(\cC) \to \cA$ extending~$f$.
\item \emph{Generator-determined chain maps.} If
$g, g' \colon \Omega^{\mathrm{ch}}(\cC) \to \cA$ are two chiral-algebra
morphisms whose restrictions to~$V$ agree and which both intertwine the
cobar differential $d_\Omega$ with the differential of~$\cA$, then $g = g'$.
\item \emph{Sign-tracked desuspension.} The cobar differential
$d_\Omega \colon \Omega^{\mathrm{ch}}(\cC) \to \Omega^{\mathrm{ch}}(\cC)$ is
the unique degree-$+1$ derivation whose restriction to generators
$V = s^{-1}\bar{\cC}$ is the composite
\[
s^{-1}\bar{\cC}
 \xrightarrow{\;s\;} \bar{\cC}
 \xrightarrow{\;-\bar{\Delta}\;}
 \bar{\cC} \otimes \bar{\cC}
 \xrightarrow{\;(s^{-1} \otimes s^{-1})\,\sigma\;}
 V \otimes V
 \;\hookrightarrow\;
 \mathrm{Free}_{\mathrm{ch}}(V),
\]
where $\sigma$ inserts the Koszul sign $(-1)^{|c'|}$ from the desuspension
of the leftmost factor in $\bar{\Delta}(c) = \sum c' \otimes c''$
\textup{(}cf.\ \cite[\S2.2.5]{LV12}\textup{)}. The same sign convention
makes $d_\Omega$ a derivation of degree~$+1$ for the chiral product on
$\mathrm{Free}_{\mathrm{ch}}(V)$, and yields $d_\Omega^2 = 0$ on
generators from coassociativity of $\bar{\Delta}$ in the uncurved case.
\end{enumerate}
\end{lemma}

\begin{proof}
\emph{Statement (i)} is the universal property of the free chiral algebra:
$\mathrm{Free}_{\mathrm{ch}} \dashv \mathrm{forget}$ on $\cD_X$-modules
\textup{(}Theorem~\ref{thm:cobar-free}\textup{)}.

\emph{Statement (ii)} follows from~(i) because the difference $g - g'$ is a
chiral-algebra-derivation that vanishes on generators, hence on the whole
free algebra by the Leibniz rule; the chain-map condition is a closed
condition on this derivation, also generator-determined.

\emph{Statement (iii).}
The desuspension $s^{-1}$ shifts cohomological degree by $-1$, so the
bigraded reduced coproduct $\bar{\Delta}$ at internal bidegree $(p,q)$
produces, after $(s^{-1}\otimes s^{-1}) \circ s$, a map of bidegree
$(p+1, q-1)$ on $V$, hence degree $+1$ as required for a cobar differential.
The Koszul sign $(-1)^{|c'|}$ in $\sigma$ records the transposition
$V \otimes V \cong s^{-1}(\bar{\cC} \otimes \bar{\cC})$ mediated by
$s \otimes s$ followed by the desuspension of the leftmost factor:
explicitly, $(s^{-1} \otimes s^{-1})(c' \otimes c'') = (-1)^{|c'|}\,
s^{-1}c' \otimes s^{-1}c''$.
The additional global sign $-1$ in front of $\bar{\Delta}$ is the
manuscript convention
$d_\Omega := -s^{-1} \circ \bar{\Delta} \circ s$; it is the
desuspended translation of the suspended convention in
\cite[\S2.2]{LV12} and makes the bar-cobar counit a chain map.
Coassociativity of $\bar{\Delta}$ then gives $d_\Omega^2 = 0$ on
generators in the uncurved case by the standard
\cite[Proposition~2.2.6]{LV12} computation:
the two iterated coproducts $(\bar{\Delta} \otimes \mathrm{id}) \bar{\Delta}$
and $(\mathrm{id} \otimes \bar{\Delta}) \bar{\Delta}$ produce equal
contributions to $d_\Omega^2$ with opposite signs, hence cancel. The
derivation property extends the uncurved identity $d_\Omega^2 = 0$ to all of
$\mathrm{Free}_{\mathrm{ch}}(V)$ by the Leibniz rule. This is the
generator-level reformulation of
Corollary~\ref{cor:cobar-nilpotence-verdier}.
\end{proof}

\begin{remark}[Scope of the chiral bar-cobar adjunction $\psi$ vs.\ qi]
\label{rem:eta-adjunction-vs-qi-scope}
\index{bar-cobar adjunction!scope}%
\index{Theorem A!scope vs Theorem B}%
The counit $\psi_{\cA}$ of Theorem~\ref{thm:bar-cobar-adjunction} \emph{exists}
without any Koszul hypothesis: it is determined by
Lemma~\ref{lem:cobar-derivation-extension}\,(i) from the augmentation alone,
and is the same map as the augmentation~$\varepsilon$ of
Proposition~\ref{prop:cobar-bar-augmentation}. This is the substance of
Theorem~A of the programme: the backbone adjunction
$\Omega^{\mathrm{ch}} \dashv \bar{B}^{\mathrm{ch}}$ is uniformly available
on the full category of augmented chiral algebras.

The further claim that $\psi_{\cA}$ is a \emph{quasi-isomorphism} is
\emph{Theorem~B} \textup{(}Theorem~\ref{thm:bar-cobar-inversion-qi}\textup{)};
it requires the Koszul hypothesis on the strict lane and a coderived
collapse input off it. Thus there are two distinct assertions:
existence of the adjunction and invertibility of the counit on the
Koszul or coderived locus.

The canonical natural transformation in this geometric direction is
the \emph{counit} $\psi$, and it points
$\Omega^{\mathrm{ch}}(\bar{B}^{\mathrm{ch}}(\cA)) \to \cA$, not the reverse.
The reverse arrow $\cA \to \Omega^{\mathrm{ch}}(\bar{B}^{\mathrm{ch}}(\cA))$
exists canonically only when $\psi_{\cA}$ is a qi, in which case it
is the homotopy inverse of $\psi_{\cA}$; this is exactly the Koszul
locus content of Theorem~\ref{thm:bar-cobar-inversion-qi}.
\end{remark}

\begin{corollary}[Chiral bar-cobar Hom-bijection backbone; \ClaimStatusConditional]
\label{cor:chiral-adjunction-hom-bijection}
For every augmented chiral algebra $\cA$ on~$X$ and every coaugmented
conilpotent chiral coalgebra~$\cC$ on~$X$, both valued in holonomic
$\cD_X$-modules, the counit and unit of
Theorem~\ref{thm:bar-cobar-adjunction} realise the natural bijection
\[
\Hom_{\mathrm{dg\,Alg}^{\mathrm{ch}}_X}\bigl(\Omega^{\mathrm{ch}}(\cC), \cA\bigr)
\;\cong\;
\mathrm{Tw}^{\mathrm{ch}}(\cC, \cA)
\;\cong\;
\Hom_{\mathrm{dg\,CoAlg}^{\mathrm{ch,conil}}_X}\bigl(\cC, \bar{B}^{\mathrm{ch}}(\cA)\bigr).
\]
This is the chiral specialisation of LV12 Theorem~2.2.9, exhibiting
$\Omega^{\mathrm{ch}} \dashv \bar{B}^{\mathrm{ch}}$ \textup{(}Theorem~A of
the programme\textup{)}, verified independently of any Koszulness or
quasi-isomorphism hypothesis on $\cA$ or~$\cC$.
\end{corollary}

\begin{proof}
The first bijection
$\Hom_{\mathrm{dg\,Alg}^{\mathrm{ch}}}(\Omega^{\mathrm{ch}}(\cC), \cA)
\cong \mathrm{Tw}^{\mathrm{ch}}(\cC, \cA)$
is the chiral instance of \cite[Theorem~2.2.7]{LV12}: a chiral-algebra map
$\Omega^{\mathrm{ch}}(\cC) \to \cA$ is determined by its restriction to the
generators $s^{-1}\bar{\cC}$
\textup{(}Lemma~\ref{lem:cobar-derivation-extension}\,(i)\textup{)}, and
the chain-map condition translates to the Maurer--Cartan equation for the
corresponding $\cD_X$-linear map $\cC \to \cA$, i.e.\ to membership in
$\mathrm{Tw}^{\mathrm{ch}}(\cC, \cA)$.

The second bijection
$\mathrm{Tw}^{\mathrm{ch}}(\cC, \cA)
\cong \Hom_{\mathrm{dg\,CoAlg}^{\mathrm{ch,conil}}}(\cC, \bar{B}^{\mathrm{ch}}(\cA))$
is dually the chiral instance of \cite[Theorem~2.2.4]{LV12}: by the
universal property of the cofree-conilpotent factorisation coalgebra
$T^c(s^{-1}\bar{\cA})$
\textup{(}Theorem~\ref{thm:coalgebra-via-NAP}\textup{)}, a coalgebra map
into $\bar{B}^{\mathrm{ch}}(\cA)$ is determined by its component projection
to $s^{-1}\bar{\cA}$, and the chain-map condition becomes the same
Maurer--Cartan equation.

Composition of the two bijections is the abstract bar-cobar adjunction. The
counit $\psi_{\cA}$ corresponds to
$\mathrm{id}_{\bar{B}^{\mathrm{ch}}(\cA)}$ under the second bijection; the
unit $\eta_{\cC}$ corresponds to $\mathrm{id}_{\Omega^{\mathrm{ch}}(\cC)}$
under the first. No Koszulness or qi hypothesis enters.
\end{proof}

\begin{remark}[Cobar inversion as the inverse half of the Koszul reflection]
\label{rem:cobar-as-koszul-reflection-inverse}
\index{cobar construction!inverse half of Koszul reflection}%
\index{Koszul reflection!cobar as $K^{-1}$}%
The counit $\psi\colon \Omegach(\barBch(\cA)) \to \cA$ of
Theorem~\ref{thm:bar-cobar-adjunction} is, under the Koszul reflection of
Theorem~\ref{thm:koszul-reflection}\,\ref{KR-i}, the counit of the
involutive symmetric-monoidal adjunction $K \dashv K^{-1}$ with
$K = \barBch$ and $K^{-1} = \Omegach$. On the Koszul locus $\Kosz(X)$
this counit is a chain-level qi
\textup{(}clause~\ref{KR-ii}: $K^{2}\simeq\mathrm{id}$\textup{)}; off it,
$\psi_{\cA}$ is a coderived isomorphism in
$D^{\co}(\barBch\text{-}\mathrm{CoFact})$
\textup{(}clause~\ref{KR-iii}\textup{)}. The cobar functor is
\emph{not} the Koszul dual of $\cA$ (Remark~\ref{rem:cobar-three-functors});
it is the inverse half of the reflection that sends $\cA$ to its bar
coalgebra and back. Three operations on $\barBch(\cA)$ remain
distinct, as in Remark~\ref{rem:cobar-three-functors}:
\begin{itemize}
\item $\Omegach(\barBch(\cA)) \simeq \cA$ (cobar = inverse, this section);
\item $\mathbb{D}_{\Ran}\,\barBch(\cA) \simeq \cA^!_\infty$
(Verdier duality on $\Ran$, gives Koszul dual factorization
\emph{algebra}; Theorem~A as recovered in
Theorem~\ref{thm:bar-cobar-isomorphism-main});
\item $C^\bullet_{\mathrm{ch}}(\cA,\cA) = Z^{\mathrm{der}}_{\mathrm{ch}}(\cA)$
(chiral Hochschild cochains, Theorem~H, closed-sector cochains).
\end{itemize}
The Koszul reflection unifies the first two as the involution
$K^{2} \simeq \mathrm{id}$ and supplies the chain-level qi on the
Koszul locus that the present cobar inversion realises in concrete
distributional form.
\end{remark}

\begin{definition}[Categorical logarithm and exponential]
\label{def:categorical-logarithm}
\index{categorical logarithm|textbf}
\index{categorical exponential}
The bar functor $\barBch$ and the cobar functor $\Omegach$ are
the \emph{categorical logarithm} and \emph{categorical
exponential} of chiral algebras on~$X$. Four properties
justify the terminology.
\begin{enumerate}[label=\textup{(\roman*)}]
\item \emph{Additivity} ($\log$ sends products to sums):
 \[
 \barBch(\cA\otimes\cA')
 \;\simeq\;
 \barBch(\cA)\;\oplus\;\barBch(\cA').
 \]
 This holds because the logarithmic forms
 $\eta_{ij}=d\log(z_i-z_j)$ on $\mathrm{Conf}_n(X)$ are
 additive under the product of chiral brackets:
 $\eta_{ij}^{\cA\otimes\cA'}
 =\eta_{ij}^\cA+\eta_{ij}^{\cA'}$.
\item \emph{Integral kernel}:
 the bar differential is the integral operator with kernel
 $\eta_{ij}=d\log(z_i-z_j)\in
 \Omega^1_{\log}(\mathrm{Conf}_2(X))$.
 The Fourier modes of $\eta_{ij}$ are the OPE residues
 $a_{(n)}b$ at all pole orders.
\item \emph{Genus-$0$ cocycle condition} ($d\log$ is closed on the affine
 screen):
 the Arnold relation
 \[
 \eta_{ij}\wedge\eta_{jk}
 +\eta_{jk}\wedge\eta_{ki}
 +\eta_{ki}\wedge\eta_{ij}
 \;=\;0
 \]
 implies $d_{\barB}^2=0$. The three terms correspond to the
 three boundary faces of $\overline{\mathrm{FM}}_3(X)$;
 the Arnold relation is their signed cancellation. For
 $X=\mathbb{P}^1$ one must also quotient by the projective linear
 relations, and for $g(X)\geq1$ the period relations of
 Proposition~\ref{prop:arnold-higher-genus} enter the same degree-$2$
 check.
\item \emph{Inversion} ($\exp\circ\log=\mathrm{id}$):
 $\Omegach(\barBch(\cA))\xrightarrow{\sim}\cA$ on the Koszul
 locus (Theorem~\ref{thm:bar-cobar-inversion-qi}). The cobar
 functor reconstructs the algebra from its bar coalgebra.
\end{enumerate}
\end{definition}

\section{Precise distribution spaces}

\begin{definition}[Fulton--MacPherson distribution space]
The space
$\operatorname{Dist}_{\mathrm{FM}}(C_n(X),\mathcal{C}^{\boxtimes n})$
is the continuous dual of compactly supported smooth test densities on
the FM compactification $\overline{C}_n(X)$, with the following
controlled boundary conditions:
\begin{itemize}
\item prescribed conormal singularities along the collision divisors
 $D_{ij}$;
\item finite-order growth along the normal-crossing boundary;
\item the Koszul-signed $\mathfrak{S}_n$ transformation law.
\end{itemize}
The older notation $\operatorname{Dist}(C_n(X),-)$ below refers to this
FM current space unless the ordinary distribution space on the open
manifold is explicitly named.
\end{definition}

\begin{theorem}[Topology; \ClaimStatusProvedHere]\label{thm:weak-topology}
The FM distribution space
$\operatorname{Dist}_{\mathrm{FM}}(C_n(X),\mathcal{C}^{\boxtimes n})$
is a complete locally convex topological vector space under the
weak-$*$ topology defined by
\[
\langle K,\phi\rangle
=\int_{\overline{C}_n(X)} K\cdot\phi
\]
for compactly supported smooth test densities~$\phi$ on
$\overline{C}_n(X)$ with boundary traces. The cobar differential
$d_\Omega$ is continuous in this topology.
\end{theorem}

\begin{proof}
\emph{Step 1: Topological structure.}
The test function space
$C_c^\infty(\overline{C}_n(X),(\mathcal{C}^{\boxtimes n})^*)$ with
boundary trace operators carries the standard nuclear LF topology. The
FM distribution space is its continuous dual, equipped with the
weak-$*$ topology: a net $K_\alpha\to K$ iff
$\langle K_\alpha,\phi\rangle\to\langle K,\phi\rangle$ for all test
densities~$\phi$. Completeness follows from the completeness of
$\mathbb{C}$ and Grothendieck completeness for duals of nuclear
LF-spaces.

\emph{Step 2: Prescribed singularities.}
The singularity condition is imposed along the FM divisors
$D_{ij}\subset\overline{C}_n(X)$, not along diagonals inside the open
stratum. Locally, if $\epsilon$ is a normal coordinate to $D_{ij}$, the
allowed conormal expansions have finite pole order
$K\sim\sum_k c_{ij}^{(k)}\epsilon^{-k}$, with $k$ bounded by the OPE
order. These finite-order conormal conditions cut out a closed subspace
of the full current space, hence preserve completeness.

\emph{Step 3: Continuity of $d_\Omega$.}
The cobar differential is the transpose of the bar residue map. For
each test density~$\phi$,
\[
\langle d_\Omega K,\phi\rangle
=\langle K,d_\Omega^*\phi\rangle,
\]
where $d_\Omega^*$ is the finite sum of boundary trace maps followed
by the coalgebra coproduct. Trace to a normal-crossing boundary stratum
and pullback along the FM structure maps are continuous maps of nuclear
LF test spaces. Hence their transpose $d_\Omega$ is weak-$*$ continuous.
\end{proof}

\begin{remark}[Regularization methods]\label{rem:regularization-methods}
When distributional products arising in the cobar differential require regularization, three standard methods are available:
\begin{enumerate}
\item \emph{Dimensional regularization}: analytic continuation in dimension $d = d_0 - \epsilon$;
\item \emph{Principal value prescription}: $\mathrm{P.V.}\int f(z)/z\,dz$;
\item \emph{Hadamard finite parts}: extracting the finite part of a divergent integral.
\end{enumerate}
For the cobar construction in the chiral setting, regularization is not
part of the definition: the intrinsic operation is the Verdier
transpose of the bar residue map. If one rewrites this operation as a
pointwise product of kernels and boundary currents, regularization may
be needed even at genus~$0$ unless the relevant wavefront sets are
transversal. The uncurved genus-$0$ identity is independent of such a
choice because the Verdier proof reduces it to the Arnold relation on
FM boundary strata. At higher genus the propagator acquires period
corrections; see the genus-$g$ analysis in
Theorem~\ref{thm:higher-genus-inversion}.
\end{remark}

\begin{example}[Cobar via integration kernels]\label{ex:cobar-kernels}
The cobar construction uses distributional integration kernels. For a chiral coalgebra $\mathcal{C}$
with coproduct $\Delta: \mathcal{C} \to \mathcal{C} \boxtimes \mathcal{C}$, elements of $\Omega^{\text{ch}}(\mathcal{C})$ are:

\[\sum_{n \geq 0} \int_{C_n(X)} K_n(z_1, \ldots, z_n) \cdot c_1(z_1) \cdots c_n(z_n) \, dz_1 \cdots dz_n\]

where:
\begin{itemize}
\item $K_n$ are FM boundary-current representatives, restricted to
 the open stratum
\item $c_i \in \mathcal{C}$ are coalgebra elements
\item any analytic regularization is auxiliary to the intrinsic
 Verdier definition
\end{itemize}

The cobar differential is written schematically as:
\[d_{\text{cobar}} = \sum_{i<j} \Delta_{ij} \cdot \delta(z_i - z_j)\]
meaning the finite sum of FM boundary trace maps followed by the
coalgebra coproduct.

This realizes the cobar complex as the distributional pairing partner
of the bar complex under configuration-space Verdier pairing, not as
the linear dual object that produces the Koszul dual algebra:
\[\langle \omega_{\text{bar}}, K_{\text{cobar}} \rangle = \int_{\overline{C}_n(X)} \omega_{\text{bar}} \wedge \iota^* K_{\text{cobar}}\]
where $\iota: C_n(X) \hookrightarrow \overline{C}_n(X)$ is the inclusion.

\emph{Conditional QFT comparison.}
\begin{itemize}
\item Bar elements = logarithmic compactified configuration kernels
\item Cobar elements = FM current kernels Verdier-dual to bar residues
\item A QFT comparison may send the pairing to a propagator/amplitude
 kernel pairing only after the BV, wavefront, and genus data of
 Theorem~\ref{thm:cobar-propagator-comparison} have been supplied
\end{itemize}
\end{example}

\subsection{Poincaré--Verdier duality realization}

\begin{theorem}[Bar-cobar Verdier pairing; \ClaimStatusProvedHere]\label{thm:poincare-verdier}
The bar and cobar constructions are related by the Verdier
intertwining of Theorem~\textup{\ref{thm:bar-cobar-isomorphism-main}}:
\[
\mathbb{D}_{\operatorname{Ran}}\bigl(\bar{B}^{\mathrm{ch}}(\cA)\bigr)
\;\simeq\; (\cA)^!_\infty,
\]
where $\mathbb{D}_{\operatorname{Ran}}$ denotes Verdier duality on
$\operatorname{Ran}(X)$ and $(\cA)^!_\infty$ is the homotopy Koszul
dual factorization algebra. On the Koszul locus,
Theorem~\textup{\ref{thm:bar-cobar-isomorphism-main}} identifies
$(\cA)^!_\infty \simeq \cA^!$.
The cobar construction $\Omega^{\mathrm{ch}}$ is the \emph{inverse},
not the dual: $\Omega^{\mathrm{ch}}(\bar{B}^{\mathrm{ch}}(\cA))
\simeq \cA$ on the Koszul locus
\textup{(}Theorem~B\textup{)}.
The Verdier pairing between bar forms and cobar distributions
is the configuration-space incarnation of this intertwining.
(The seven-face identification of $r(z)$ is the subject of Chapter~\ref{ch:holographic-datum-master}.)
\end{theorem}

\begin{proof}[Geometric realization]
The duality is realized through the perfect pairing:
\[
\langle \omega_{\text{bar}}, \omega_{\text{cobar}} \rangle = \int_{\overline{C}_n(X)} \omega_{\text{bar}} \wedge \iota^*\omega_{\text{cobar}}
\]
where $\iota: C_n(X) \hookrightarrow \overline{C}_n(X)$ is the inclusion.

The duality has the following structure:
\begin{itemize}
\item Logarithmic forms on $\overline{C}_n(X)$ (bar) are dual to distributions on $C_n(X)$ (cobar)
\item Residues at divisors (bar) are dual to principal value integrals (cobar)
\item Collision divisors (bar) correspond to extension loci (cobar)
\item The duality exchanges extraction (analysis) with reconstruction (synthesis)
\end{itemize}
\end{proof}

\subsection{Explicit cobar computations}

\begin{example}[Cobar of exterior coalgebra]\label{ex:cobar-exterior}
Let $\mathcal{E} = \Lambda^*_{\text{ch}}(V)$ be the chiral exterior coalgebra on generators $V$. Then:
\[
\Omega^{\text{ch}}(\mathcal{E}) \cong S_{\text{ch}}(s^{-1}V)
\]
the chiral symmetric algebra on the desuspension of $V$.

Geometrically, this duality is realized by:
\begin{itemize}
\item Fermionic fields $\psi \in V$ with antisymmetric OPE become bosonic fields $\phi \in s^{-1}V$ with symmetric OPE
\item The cobar differential vanishes since the reduced comultiplication $\bar{\Delta}(\psi) = 0$
\item Configuration space integrals enforce bosonic statistics through symmetric integration domains
\end{itemize}

This is the chiral analogue of the classical Koszul duality between exterior and symmetric algebras.
\end{example}

\begin{example}[Cobar of bar of free fermions]\label{ex:cobar-bar-fermion}
For the free fermion algebra $\mathcal{F}$, bar-cobar inversion gives:
\[
\Omega^{\text{ch}}(\bar{B}^{\text{ch}}(\mathcal{F})) \xrightarrow{\sim} \mathcal{F}
\]
recovering the free fermion algebra itself (not the $\beta\gamma$
system). The Koszul dual $\mathcal{F}^! \cong \beta\gamma$ is
obtained instead from the cohomological coalgebra
$\mathcal F^i=H^*(\bar{B}^{\text{ch}}(\mathcal{F}))$ by finite-type
linear duality:
$\mathcal{F}^!=(\mathcal F^i)^\vee$
(see Example~\ref{ex:fermion-betagamma-bar-cobar} and
Theorem~\ref{thm:fermion-boson-koszul}). This distinction is the
chiral analog of the classical fact that
$\Omega(\bar{B}({\Lambda}(V))) \simeq {\Lambda}(V)$
(bar-cobar inversion), while
${\Lambda}(V)^! = \mathrm{Sym}(V^*)$ (Koszul dual algebra via
linear duality).
\end{example}


\subsection{\texorpdfstring{Cobar $A_\infty$ structure}{Cobar A-infinity structure}}
\label{subsec:cobar-diff-ainfty}

\begin{theorem}[\texorpdfstring{$A_\infty$}{A-infinity} structure on cobar {\cite{LV12}}; \ClaimStatusProvedElsewhere]\label{thm:cobar-ainfty}
The cobar construction $\Omega^{\text{ch}}(\mathcal{C})$ carries a canonical $A_\infty$ structure with operations:
\[
m_k: \Omega^{\text{ch}}(\mathcal{C})^{\otimes k} \to \Omega^{\text{ch}}(\mathcal{C})[2-k]
\]
geometrically realized by:
\[
m_k(\alpha_1, \ldots, \alpha_k) = \int_{\partial \overline{M}_{0,k+1}} \alpha_1 \wedge \cdots \wedge \alpha_k \wedge \omega_{0,k+1}
\]
where $\overline{M}_{0,k+1}$ is the moduli space of stable curves with $k+1$ marked points.
\end{theorem}

\begin{proof}
By \cite[Theorem~9.2.8]{LV12}, $A_\infty$ structures arise from operadic bar-cobar resolutions.
The boundary stratification of the Deligne--Mumford compactification gives:
\[
\partial \overline{M}_{0,k+1} = \bigcup_{\substack{I \sqcup J = [k+1] \\ |I|,|J| \geq 2}} \overline{M}_{0,|I|+1} \times \overline{M}_{0,|J|+1}
\]
Each boundary stratum $\overline{M}_{0,|I|+1} \times \overline{M}_{0,|J|+1}$ contributes a composition $m_{|I|} \circ m_{|J|}$ to the boundary of $m_k$. The topological identity $\partial^2 = 0$ on the chain complex $C_*(\overline{M}_{0,k+1})$ then gives $\sum_{i+j=k+1} m_i \circ m_j = 0$, which is the $A_\infty$ relation. Signs are determined by the orientation of the boundary strata, matching the Koszul sign convention in \cite[\S10.3]{LV12}.
\end{proof}

\subsection{Geometric cobar for curved coalgebras}

\begin{definition}[Curved cobar]\label{def:curved-cobar}
Let $(\mathcal{C},d_\mathcal C,\bar\Delta,h_C)$ be a curved
coaugmented chiral coalgebra in the finite-type or completed
holonomic ambient, where $h_C$ is the curvature datum dual to the
curved bar square. The cobar construction is the CDG chiral algebra
\[
\Omega^{\mathrm{ch}}_{\mathrm{curv}}(\mathcal{C})
=
\bigl(\Omega^{\mathrm{ch}}(\mathcal{C}), d_{\Omega,h}, h_\Omega\bigr),
\]
not an ordinary cochain complex in general. The element $h_\Omega$ is
the Verdier-transposed curvature, equivalently the image of the bar
curvature under the pairing of Theorem~\ref{thm:bar-cobar-verdier}.
The defining CDG identities are
\[
d_{\Omega,h}^{\,2}=[h_\Omega,-],
\qquad
d_{\Omega,h}(h_\Omega)=0.
\]
Only after a separate QFT comparison datum has been supplied may
$h_\Omega$ be represented by an analytic kernel integral. That
representation is not part of the algebraic definition. The uncurved
cobar complex is the special case $h_\Omega=0$.
\end{definition}

\begin{theorem}[Curved Maurer--Cartan equation; \ClaimStatusProvedHere]\label{thm:curved-mc-cobar}
\textup{[Regime: curved-central
\textup{(}Convention~\textup{\ref{conv:regime-tags})}.]}

For $\alpha \in \Omega^{\text{ch}}(\mathcal{C})[-1]$, set
\[
h_{\Omega,\alpha}
:=
h_\Omega+d_{\Omega,h}\alpha+\frac{1}{2}[\alpha,\alpha].
\]
The twisted derivation
$d_{\Omega,\alpha}=d_{\Omega,h}+[\alpha,-]$ satisfies
\[
d_{\Omega,\alpha}^{\,2}=[h_{\Omega,\alpha},-].
\]
Thus $\alpha$ satisfies the curved Maurer--Cartan equation precisely
when $h_{\Omega,\alpha}=0$; only then does the twist produce a strict
square-zero cobar differential. Such curvature-killing twists
correspond geometrically to:
\begin{itemize}
\item Deformations of the chiral structure that do not preserve the grading
\item Quantum anomalies in the conformal field theory
\item Central extensions and their geometric representatives
\end{itemize}
\end{theorem}

\begin{proof}
By the general theory of curved $A_\infty$ and CDG algebras
(Positselski~\cite{Positselski11}, Loday--Vallette~\cite{LV12}
\S10.1), twisting a curved object by $\alpha$ changes the curvature to
\[
h_{\Omega,\alpha}
=h_\Omega+d_{\Omega,h}\alpha+\frac{1}{2}[\alpha,\alpha],
\]
and changes the derivation to
$d_{\Omega,\alpha}=d_{\Omega,h}+[\alpha,-]$. The CDG identity gives
$d_{\Omega,\alpha}^{\,2}=[h_{\Omega,\alpha},-]$. Hence a twist is an
ordinary differential exactly when the curved Maurer--Cartan equation
$h_{\Omega,\alpha}=0$ holds.

In the chiral setting, $h_\Omega$ is the cobar-side representative of
the same anomaly measured by the curved bar square. For affine
Kac--Moody algebras its scalar component is the curvature computed in
Proposition~\ref{prop:km-bar-curvature}. At genus~$1$, curvature-zero
solutions give central extensions by $\mathbb{C}$ as in
Section~\ref{sec:genus_1_central_extensions}; with nonzero curvature
one remains in the CDG/coderived category.
\end{proof}

\begin{proposition}[Affine modular curvature scalar; \ClaimStatusProvedHere]
\label{prop:km-bar-curvature}
\index{curvature!affine modular scalar}
For the universal affine Kac--Moody chiral algebra
$V_k(\mathfrak g)$ at level~$k$, the genus-$0$ PBW bar differential is
the uncurved chiral bar differential. In the curved-central modular
fiberwise lane, the scalar Hodge curvature of the genus-$g$ bar family is
\begin{equation}\label{eq:km-curvature-part1}
\kappa^{\mathrm{KM}}\!\bigl(V_k(\mathfrak g)\bigr)
=\frac{\dim(\mathfrak g)(k+h^\vee)}{2h^\vee}.
\end{equation}
Here $h^\vee$ is the dual Coxeter number and the trace form is
normalized with long roots of length~$2$. Thus the raw modular
fiberwise operator has scalar square
\[
d_{\mathrm{fib}}^{\,2}
=\bigl[\kappa^{\mathrm{KM}}(V_k(\mathfrak g))\,\omega_g,-\bigr]
\]
after projection to the scalar Hodge lane. In particular:
\begin{enumerate}[label=\textup{(\roman*)}]
\item the scalar modular curvature vanishes if and only if
$k=-h^\vee$;
\item away from the critical level, the raw genus-$g$ fiberwise
operator belongs to the CDG/coderived lane unless one uses a
curvature-killing Maurer--Cartan twist or the flat total comparison
differential.
\end{enumerate}
This statement concerns the scalar modular fiberwise curvature of the
universal affine family. It is not the periodic-CDG curvature of a
simple admissible quotient $L_k(\mathfrak g)$.
\end{proposition}

\begin{proof}
The scalar formula is the standard-landscape normalization
of Chapter~\ref{chap:landscape-census}. The trace-form double-pole
projection in the affine OPE contributes
\(\dim(\mathfrak g)k/(2h^\vee)\), while the simple-pole Lie-bracket
channel contributes \(\dim(\mathfrak g)/2\). Their sum is
\[
\frac{\dim(\mathfrak g)k}{2h^\vee}
+\frac{\dim(\mathfrak g)}{2}
=\frac{\dim(\mathfrak g)(k+h^\vee)}{2h^\vee}.
\]
This is exactly
\(\kappa^{\mathrm{KM}}(V_k(\mathfrak g))\).
The CDG identity for the raw genus-$g$ fiberwise operator is the
curved-central relation
$d_{\mathrm{fib}}^{\,2}=[\kappa(\cA)\omega_g,-]$ applied to
\(\cA=V_k(\mathfrak g)\).

The genus-$0$ PBW bar differential is square-zero by the Borcherds
identity and the Arnold relation, independently of~$k$. The curvature
above is the modular fiberwise scalar; it vanishes precisely when
$k+h^\vee=0$. The simple quotient $L_k(\mathfrak g)$ at admissible
level can acquire additional curvature from vacuum null vectors, but
that is a separate quotient-level assertion.
\end{proof}

\begin{remark}[Feigin--Frenkel center at critical level]
\label{rem:feigin-frenkel-center}
\index{Feigin--Frenkel center}
At the critical level $k = -h^\vee$,
Proposition~\ref{prop:km-bar-curvature} gives
\(\kappa^{\mathrm{KM}}(V_{-h^\vee}(\mathfrak g))=0\), so the scalar
modular curvature vanishes. The genus-$0$ PBW bar complex of the
universal affine algebra is square-zero at every level; at the critical
level a foundational result of Feigin--Frenkel~\cite{Feigin-Frenkel}
identifies the center of $\widehat{\mathfrak{g}}_{-h^\vee}$ with the
algebra of functions on $\mathfrak{g}^\vee$-opers:
\[
Z(\widehat{\mathfrak{g}}_{-h^\vee}) \cong \mathrm{Fun}(\mathrm{Op}_{\check{\mathfrak{g}}}(X)).
\]
The critical level is the fixed point of the level-shifting involution
$k \mapsto -k - 2h^\vee$, and bar-cobar inversion recovers
$\widehat{\mathfrak{g}}_{-h^\vee}$ from its bar coalgebra on the
Koszul/completed locus. See Part~\ref{part:characteristic-datum} for
the detailed treatment.
\end{remark}

\begin{corollary}[Level-shifting via Verdier duality; \ClaimStatusConditional]
\label{cor:level-shifting-part1}
\index{Feigin--Frenkel duality}
\index{level!shifted}
For any simple Lie algebra $\mathfrak{g}$ and level $k \neq -h^\vee$
on the same-family affine Koszul-duality lane,
the post-Verdier factorization-algebra identification
(Theorem~\ref{thm:bar-cobar-isomorphism-main},
Convention~\ref{conv:bar-coalgebra-identity}) gives
\begin{equation}\label{eq:level-shifting-part1}
\mathbb{D}_{\operatorname{Ran}}\bigl(
 \bar{B}^{\mathrm{ch}}(\widehat{\mathfrak{g}}_k)\bigr)
\;\simeq\;
\bigl(\widehat{\mathfrak{g}}_k\bigr)^!_\infty
\;\simeq\;
\widehat{\mathfrak{g}}_{k'},
\qquad k' = -k - 2h^\vee.
\end{equation}
Bar-cobar \emph{inversion} recovers the original algebra:
$\Omega^{\mathrm{ch}}(\bar{B}^{\mathrm{ch}}(\widehat{\mathfrak{g}}_k))
\simeq \widehat{\mathfrak{g}}_k$ on the Koszul locus
(Theorem~B). The level-shifted algebra
$\widehat{\mathfrak{g}}_{k'}$ is the
\emph{Koszul dual}, obtained by Verdier duality on
$\operatorname{Ran}(X)$, not by the cobar functor.

The dual level $k' = -k - 2h^\vee$ satisfies:
\begin{enumerate}[label=\textup{(\roman*)}]
\item $k + k' = -2h^\vee$;
\item the central charges satisfy $c(\widehat{\mathfrak{g}}_k) + c(\widehat{\mathfrak{g}}_{k'}) = 2\dim\mathfrak{g}$ \textup{(}Theorem~\ref{thm:central-charge-complementarity}\textup{)};
\item the scalar modular characteristics are opposite:
\(\kappa^{\mathrm{KM}}(V_{k'}(\mathfrak g))
=-\kappa^{\mathrm{KM}}(V_k(\mathfrak g))\).
\end{enumerate}
\end{corollary}

\begin{proof}
The Verdier intertwining on the factorization-algebra side
(Theorem~\ref{thm:bar-cobar-isomorphism-main},
Convention~\ref{conv:bar-coalgebra-identity}) identifies
$\mathbb{D}_{\operatorname{Ran}}(\bar{B}^{\mathrm{ch}}(\cA))
\simeq (\cA)^!_\infty$ for any chiral algebra $\cA$, and on the
Koszul locus $(\cA)^!_\infty \simeq \cA^!$. For
$\cA = \widehat{\mathfrak{g}}_k$, the Koszul dual
is $\cA^! = \widehat{\mathfrak{g}}_{k'}$ with
$k' = -k - 2h^\vee$: the affine same-family Verdier involution
preserves the simple-pole Lie bracket and sends the scalar modular
characteristic to its negative, since
\[
\kappa^{\mathrm{KM}}(V_{k'}(\mathfrak g))
=\frac{\dim(\mathfrak g)(k'+h^\vee)}{2h^\vee}
=-\frac{\dim(\mathfrak g)(k+h^\vee)}{2h^\vee}.
\]
The dual algebra has OPE
\[
J^{\prime a}(z) J^{\prime b}(w) \sim
\frac{k'\,(J^a, J^b)_{\mathrm{tr}}}{(z-w)^2}
+ \frac{f^{ab}{}_c\, J^{\prime c}(w)}{z-w}
\]
with $k' = -k - 2h^\vee$. Parts~(i)--(iii) follow by
direct substitution.
\end{proof}

\subsection{Computational algorithms for cobar}

\begin{remark}[Cobar complex computation]
Given a chiral coalgebra $\mathcal{C}$ with basis $\{e_i\}$, structure constants $\Delta(e_i) = \sum_{j,k} c_{jk}^i e_j \otimes e_k$, and counit $\epsilon$, the cobar complex $(\Omega^{\mathrm{ch}}(\mathcal{C}), d_{\mathrm{cobar}})$ is constructed as follows.
\begin{enumerate}
\item Initialize $\Omega^0 = \mathrm{Free}_{\mathrm{ch}}(s^{-1}\bar{\mathcal{C}})$, where $\bar{\mathcal{C}}$ is the coaugmentation coideal of Definition~\ref{def:conilpotent-cobar}.
\item For each generator $s^{-1}e_i$ with $e_i\in\bar{\mathcal C}$ and reduced coproduct $\bar\Delta(e_i)=\sum_{j,k} c_{jk}^i e_j\otimes e_k$, compute
\[
d(s^{-1}e_i)
=-s^{-1}(d_{\mathcal C}e_i)
-\sum_{j,k}(-1)^{|e_j|}c_{jk}^i\,
s^{-1}e_j\otimes s^{-1}e_k .
\]
\item Extend to products by the Leibniz rule: $d(xy) = d(x)y + (-1)^{|x|}xd(y)$.
\item For each $n$-fold product, tensor with $\Omega^*(C_{n+1}(X))$.
\item Impose the Verdier-transposed Arnold--Orlik--Solomon relations among FM boundary currents and the factorization constraints from the chiral structure.
\end{enumerate}
\end{remark}

\section{Genus 1 contributions: central extensions in the bar-cobar complex}
\label{sec:genus_1_central_extensions}

Everything constructed so far lives on a \emph{fixed} curve~$X$.
The bar differential is a residue on $\overline{C}_n(X)$; the
cobar differential inserts delta functions on $C_n(X)$; Verdier
duality exchanges the two. None of this involves the moduli of
curves.

Now let the curve vary. Replace $X$ by a family
$\pi \colon \mathcal{X} \to S$ over a base~$S$, and ask: how does
the bar complex change as $X$ moves in $\overline{\mathcal{M}}_g$?
At genus~$0$ the answer is: it does not; there are no moduli.
At genus~$1$, the propagator $K^{(1)} = \theta_1'/\theta_1$
depends on the modular parameter $\tau$, and the bar differential
acquires a correction. This correction is the central extension:
curvature $m_0 = \kappa(\cA)$ emerges because the trace of the
OPE around the $A$-cycle of the torus is nonzero. The local central
term is visible already in the genus-zero OPE; the genus-one trace
turns that local coefficient into the first scalar level of the
shadow obstruction tower
(Definition~\ref{def:shadow-postnikov-tower}):
the scalar level $\Theta_{\cA}^{\leq 2} = \kappa(\cA)$.

\begin{remark}[Central extension as tower level]
\label{rem:central-extension-tower-level}
\index{shadow obstruction tower!genus hierarchy}
The hierarchy of what becomes visible at each genus mirrors the
tower levels:
\begin{itemize}
\item Genus~$0$: $\Theta|_{\hbar=0}$, the strict bar-cobar
 adjunction (Theorem~A).
\item Genus~$1$: $\Theta_\cA^{\leq 2} = \kappa(\cA)$, the
 central extension / modular characteristic (Theorem~D).
\item Higher genus: the cubic and quartic shadows
 $\mathfrak{C}$, $\mathfrak{Q}$ contribute to the
 $g \geq 2$ bar differential via the shadow obstruction tower
 (Definition~\ref{def:shadow-postnikov-tower}).
\end{itemize}
\end{remark}

\subsection{Central trace class at genus 1}

\subsubsection{Heisenberg normalization}

Let $E_\tau=\mathbb C/(\mathbb Z+\tau\mathbb Z)$ and let
$\mathcal H_\ell(V)$ be the Heisenberg chiral algebra attached to a
finite-dimensional vector space $V$ with non-degenerate symmetric form
$(\, ,\,)$. The current modes satisfy
\[
[x_m,y_n]=\ell\,m\,(x,y)\,\delta_{m+n,0}\,\mathbf K,
\qquad x,y\in V.
\]
Equivalently, the local OPE is
\[
x(z)y(w)\sim \frac{\ell(x,y)}{(z-w)^2}.
\]
The double pole is already present in the local genus-zero OPE. What is
new at genus~$1$ is the trace functional obtained by integrating along a
homology cycle of $E_\tau$; this functional extracts the scalar
$\ell=\kappa(\mathcal H_\ell)$ as a finite shadow coefficient.

\subsection{The geometric construction: configuration spaces on the torus}

\subsubsection{Setup: the genus 1 configuration space}

Let $\mathbb{T}^2 = \mathbb{C}/\Lambda$ be a torus with period lattice $\Lambda$.
Define:
\[\mathrm{Conf}_n(\mathbb{T}^2) = \{ (z_1, \ldots, z_n) \in (\mathbb{T}^2)^n \mid z_i \neq z_j \}\]

The genus 1 bar complex is:
\[C_{\bullet}^{(1)}(\mathcal{A}) = \mathcal{C}_{\bullet}(\mathrm{Conf}_{\bullet}(\mathbb{T}^2),
\mathcal{A}^{\boxtimes \bullet})\]
chains on configuration space with coefficients in $\mathcal{A}$.

\subsubsection{The trace element}

For a primitive cycle $\gamma\in H_1(E_\tau,\mathbb Z)$ and a normalized
one-form $\omega_\gamma$ with $\int_\gamma\omega_\gamma=1$, define the
trace functional on a torus insertion by
\[
\operatorname{Tr}_\gamma(A)=\int_\gamma Y(A,z)\,\omega_\gamma .
\]
For the Heisenberg modes this normalization gives
$\operatorname{Tr}_\gamma(\mathbf 1)=1$ and
$\operatorname{Tr}_\gamma(x_m)=0$ for $m\neq 0$. The class depends on the
oriented homology class $\gamma$ and transforms under the natural
$\operatorname{SL}_2(\mathbb Z)$ action on $H_1(E_\tau)$; the extracted
scalar coefficient is the modular characteristic
$\kappa(\mathcal H_\ell)=\ell$.

\subsubsection{Explicit formula for central charge cocycle}

For the Heisenberg algebra, the genus-one trace cocycle is the bilinear
form
\[
\omega_\ell(x_m,y_n)
  =\ell\,m\,(x,y)\,\delta_{m+n,0}.
\]

\begin{theorem}[Central charge cocycle; \ClaimStatusProvedHere]\label{thm:central-charge-cocycle}
\textup{[Regime: curved-central
\textup{(}Convention~\textup{\ref{conv:regime-tags})}.]}

For the Heisenberg current Lie algebra
$V[t,t^{-1}]$, the bilinear form
\[
\omega_\ell(x\otimes t^m,y\otimes t^n)
  =\ell\,m\,(x,y)\,\delta_{m+n,0}
\]
is a Chevalley--Eilenberg $2$-cocycle. Its class is nonzero and defines
the standard central extension
\[
0\to \mathbb C\mathbf K\to
\widehat{V[t,t^{-1}]}_\ell\to V[t,t^{-1}]\to 0.
\]
On the genus-one bar side this is the finite trace shadow extracted
from the diagonal double-pole coefficient of
$\mathcal H_\ell(V)$.
\end{theorem}

\begin{proof}
The current Lie algebra $V[t,t^{-1}]$ is abelian before the central
extension, so the Chevalley--Eilenberg differential of any alternating
bilinear form is zero. The form above is alternating because
\[
\omega_\ell(y_n,x_m)
 =\ell\,n\,(y,x)\,\delta_{n+m,0}
 =-\ell\,m\,(x,y)\,\delta_{m+n,0}.
\]
It is not a coboundary: for an abelian Lie algebra, every
Chevalley--Eilenberg $1$-cochain has zero coboundary. Therefore
$[\omega_\ell]\neq 0$. In modes it gives exactly
$[x_m,y_n]=\omega_\ell(x_m,y_n)\mathbf K$, which is the standard
Heisenberg central extension.

The geometric representative is obtained by pairing the local OPE
coefficient $\ell(x,y)/(z-w)^2$ with the normalized cycle trace on
$E_\tau$. In the finite trace quotient the vacuum trace is normalized to
$1$ and nonzero one-point Heisenberg traces vanish, so the only surviving
scalar from the two-current collision is
$\ell m(x,y)\delta_{m+n,0}$. This is a finite shadow extraction, not an
unconverged graph sum and not a universal statement for arbitrary vertex
algebras.
\end{proof}

\subsection{Formal calculations: degree-by-degree analysis}

\subsubsection{Degree 0: the vacuum}

$C_0^{(1)} = \mathbb{C} \cdot \mathbbm{1}$, the vacuum state.

\subsubsection{Degree 1: trace insertions}

\[C_1^{(1)} = \operatorname{span}\{ \operatorname{Tr}_\gamma(x_n) \mid x\in V,\ n \in \mathbb{Z} \}\]

The differential $d^{(1)}: C_1^{(1)} \to C_0^{(1)}$ maps:
\begin{align}
d^{(1)}[\operatorname{Tr}_\gamma(x_n)] &= 0 \quad \text{for } n \neq 0, \\
d^{(1)}[\operatorname{Tr}_\gamma(x_0)] &= 0
\end{align}

\emph{Homology.} In the finite trace quotient, the degree-one classes are
represented by the trace insertions modulo the chosen zero-mode
relations.

\subsubsection{Degree 2: the central charge emerges}

\[C_2^{(1)} = \operatorname{span}\{
\operatorname{Tr}_\gamma(x_m \otimes y_n)
\mid x,y\in V,\ m,n\in\mathbb Z
\}\]

\emph{The computation.}
\begin{align}
d^{(1)}[\operatorname{Tr}_\gamma(x_m \otimes y_n)]
&= \operatorname{Tr}_\gamma\bigl([x_m,y_n]\bigr) \\
&= \ell\,m\,(x,y)\,\delta_{m+n,0}\cdot \mathbbm{1}.
\end{align}

The modular characteristic $\kappa(\mathcal H_\ell)=\ell$ appears only in the residue
constraint $m+n=0$. Geometrically, this is the trace of the diagonal
double-pole coefficient around the chosen genus-one cycle.

\subsubsection{Degrees 3--5: modular corrections}

In higher degrees the finite trace quotient has pairwise OPE
contractions and modular forms from the elliptic propagator. The first
new terms are:
\begin{itemize}
\item products of the central residue cocycle $\omega_\ell$;
\item modular dependence on the torus parameter $\tau$ through the
 elliptic propagator;
\item the quasi-modular weight-$2$ term $E_2(\tau)$ after the usual
 elliptic regularization.
\end{itemize}
These terms are finite-window shadow data. Passing from them to the
full tower requires the completed graph-sum formalism of
Theorem~\ref{thm:prism-higher-genus}.

\subsection{The cobar resolution: recovering central extensions}

The cobar construction recovers a central extension from the finite
trace coalgebra only after the reduced coproduct and the adjunction
counit are fixed.

\begin{theorem}[Genus 1 cobar extraction of the Heisenberg central extension; \ClaimStatusProvedHere]\label{thm:genus1-cobar-bar}
\textup{[Regime: curved-central, genus-$1$
\textup{(}Convention~\textup{\ref{conv:regime-tags})}.]}

Let $C_{\le2,\gamma}^{(1)}(\mathcal H_\ell(V))$ be the finite
genus-one trace quotient of the bar coalgebra generated by the vacuum,
one-current trace insertions, and two-current trace insertions, with
reduced coproduct dual to the cocycle
\[
\omega_\ell(x_m,y_n)=\ell\,m\,(x,y)\,\delta_{m+n,0}.
\]
Then the completed cobar algebra has
\[
H^0\!\left(\widehat{\Omega}\,
C_{\le2,\gamma}^{(1)}(\mathcal H_\ell(V))\right)
\cong
\widehat U\bigl(\widehat{V[t,t^{-1}]}_\ell\bigr)
\]
with central relation
\[
[x_m,y_n]=\ell\,m\,(x,y)\,\delta_{m+n,0}\,\mathbf K.
\]
For a general vertex algebra this conclusion requires a finite trace
quotient, a closed genus-one cocycle, and the filtered counit
quasi-isomorphism of Corollary~\ref{cor:bar-cobar-inverse}; it is not a
formal consequence of cobar inversion alone.
\end{theorem}

\begin{proof}
The quotient $C_{\le2,\gamma}^{(1)}$ is finite and conilpotent, so the
completed cobar differential is the continuous derivation determined on
desuspended generators by the reduced coproduct. The degree-two
generator dual to
$\operatorname{Tr}_\gamma(x_m\otimes y_n)$ has reduced coproduct whose
antisymmetric part is dual to
$\omega_\ell(x_m,y_n)$. Therefore the degree-zero relations in the
cobar algebra are precisely
\[
x_my_n-y_nx_m
  =\ell\,m\,(x,y)\,\delta_{m+n,0}\,\mathbf K.
\]
All other reduced coproduct components in the quotient are primitive or
central, so no additional degree-zero relations appear. This gives the
completed enveloping algebra of the standard central extension.

The argument uses only the finite Heisenberg trace quotient and the
standard cocycle of Theorem~\ref{thm:central-charge-cocycle}. For a
general vertex algebra, an arbitrary genus-one trace does not force
universality; the missing data are exactly the closed cocycle,
finite-window completion, and filtered counit quasi-isomorphism stated
in the theorem.
\end{proof}

\begin{remark}[Central extension as genus-\texorpdfstring{$1$}{1} trace component of \texorpdfstring{$\Theta$}{Theta}]
\label{rem:central-ext-theta1}
In the finite central trace lane, the cocycle
$\omega_{\kappa(\cA)} = \kappa(\cA) \cdot \mathrm{Tr}$ is the genus-$1$ trace
component $\theta_1$ of the universal MC class
(Theorem~\ref{thm:explicit-theta}):
$\theta_1 = \kappa(\cA) \cdot \mu \otimes \lambda_1$,
where $\mu = [-,-]$ is the bracket in the finite trace quotient and
$\lambda_1$ is the Hodge class. For Heisenberg and affine
Kac--Moody current algebras this recovers the standard central
extension class. For a general vertex algebra the genus-one component is
a shadow class in the relevant chiral Hochschild complex, and its
universality is an additional recognition theorem rather than a
consequence of cobar inversion.
\end{remark}

\subsection{Census constants at genus 1}

The finite genus-one trace extraction records the following constants.

\begin{itemize}
\item \emph{Kac--Moody algebras:} For the affine vertex algebra
$V_k(\mathfrak{g})$ at level~$k$, the modular characteristic is
\[
\kappa(V_k(\mathfrak{g}))
=\frac{\dim(\mathfrak{g})(k+h^\vee)}{2h^\vee}.
\]
This is the census modular characteristic, not the Sugawara central
charge. The Heisenberg algebra is a separate family with
$\kappa(\mathcal{H}_k)=k$.

\item \emph{Virasoro:} For the Virasoro vertex algebra,
$\kappa(\mathrm{Vir}_c)=c/2$.

\item \emph{$W$-algebras:} For the principal $W_N$ family,
\[
\kappa(\mathcal W_N)=c(H_N-1),
\qquad H_N=\sum_{j=1}^N\frac1j.
\]
\end{itemize}

\subsection{Genus 1 summary}

\begin{center}
\begin{tabular}{|l|l|l|}
\hline
\emph{Algebra} & \textbf{Finite Trace Datum} & \textbf{Bar-Cobar Output} \\
\hline
Central extension & $\omega_{\kappa(\mathcal A)}$ & Genus 1 cocycle \\
Modular characteristic $\kappa$ & diagonal double pole & trace coefficient \\
Level of Kac--Moody & $k+h^\vee$ curvature & finite shadow class \\
Virasoro $c$ & $c/2$ & $\operatorname{Tr}(T \otimes T)$ \\
\hline
\end{tabular}
\end{center}

\begin{remark}[Functoriality]
The finite trace extraction is functorial for morphisms
$f\colon\mathcal{A}\to\mathcal{B}$ that preserve the chosen genus-one
trace quotient and the cocycle:
\[
f^*\omega_{\kappa(\mathcal B)}=\omega_{\kappa(\mathcal A)}.
\]
Such a morphism induces
\[
C_{\le2,\gamma}^{(1)}(\mathcal{A})
\longrightarrow
C_{\le2,\gamma}^{(1)}(\mathcal{B})
\]
and hence a map of completed cobar algebras. A morphism of vertex
algebras that does not preserve the trace quotient or the modular
characteristic need not respect the genus-one central class.
\end{remark}

\subsection{Extension theory: from genus 0 to higher genus}

\subsubsection{The obstruction complex}

Not every genus 0 chiral algebra extends to higher genus. The obstructions live in specific cohomology groups:

\begin{theorem}[Extension obstruction class {\cite{BD04}}; \ClaimStatusConditional]
\label{thm:extension-obstruction-bar}
Let $\mathcal{A}$ be a chiral algebra on $\mathbb{CP}^1$ together with
a specified deformation problem for extending its factorization data
over a genus-$g$ family. Let
$\mathbb T_{\mathrm{ext}}(\mathcal A)$ be the corresponding tangent
complex of chiral operations with sewing constraints. The obstruction
class lies in
\[
\operatorname{Obs}_g(\mathcal{A})
\in
H^2\!\left(\overline{\mathcal{M}}_g,
\mathcal H^0(\mathbb T_{\mathrm{ext}}(\mathcal A))\right).
\]
On the central trace lane this projects to a class in
$H^2(\overline{\mathcal{M}}_g,\mathcal Z(\mathcal A))$.
\end{theorem}

\begin{proof}
By the deformation theory of chiral algebras on families of curves
(Theorem~\ref{thm:obstruction-quantum}), extensions over the moduli
base are controlled by the sheafified deformation complex
$\mathbb T_{\mathrm{ext}}(\mathcal A)$. The first obstruction to
lifting a first-order extension over the genus-$g$ family is the usual
degree-$2$ obstruction class:
\[
\operatorname{Obs}_g(\mathcal{A})
\in
H^2\!\left(\overline{\mathcal{M}}_g,
\mathcal H^0(\mathbb T_{\mathrm{ext}}(\mathcal A))\right).
\]
Projecting the tangent complex to its central quotient gives the
central obstruction class in
$H^2(\overline{\mathcal{M}}_g,\mathcal Z(\mathcal A))$. For the
standard comparison families this gives:
\begin{enumerate}
\item for the bosonic matter--ghost BRST system, the central obstruction
is proportional to $c_{\mathrm{matter}}-26$;
\item for the superstring ghost system in the standard normalization,
the central obstruction is proportional to $c_{\mathrm{matter}}-15$;
\item for free theories, extension requires the usual spin or charge
sector data, and the remaining scalar obstruction is read in the center
local system rather than in a traceless endomorphism sheaf.
\end{enumerate}
\end{proof}

\begin{example}[Free fermion extension]
The free fermion extends to all genera with spin structure:

For genus 1: The extension depends on the choice of spin structure (periodic/antiperiodic boundary conditions):
\[
\mathcal{F}_{E_\tau}^{\text{NS}} = \bigoplus_{n \in \mathbb{Z}} \mathcal{F}_n \quad \text{(Neveu--Schwarz)}
\]
\[
\mathcal{F}_{E_\tau}^{\text{R}} = \bigoplus_{n \in \mathbb{Z} + 1/2} \mathcal{F}_n \quad \text{(Ramond)}
\]

The partition function encodes the obstruction:
\[
Z_{\text{ferm}}(\tau) = \frac{\theta_3(0|\tau)}{\eta(\tau)} \quad \text{(NS sector)}
\]
\end{example}

\subsubsection{The tower of extensions}

\begin{theorem}[Genus-filtered completed tower; \ClaimStatusConditional]\label{thm:universal-extension-tower}
Assume that the modular bar coalgebra of $\mathcal A$ exists as a
complete genus-filtered coalgebra
\[
\barB^{\mathrm{mod}}(\mathcal A)
\simeq \widehat{\bigoplus}_{g\ge0}\barB^{(g)}(\mathcal A)
\]
and that the genus filtration is compatible with the coproduct,
curvature, and sewing maps. Set
\[
\cA_{\le g}:=
\widehat\Omega\!\left(
\barB^{\mathrm{mod}}(\mathcal A)/F_{>g}\barB^{\mathrm{mod}}(\mathcal A)
\right).
\]
Then there is an inverse tower of completed cobar algebras
\[
\cdots\longrightarrow
\cA_{\le g+1}\longrightarrow
\cA_{\le g}\longrightarrow\cdots\longrightarrow\cA_{\le0}.
\]
If the tower is Mittag--Leffler complete, its limit is
\[
\cA_\infty\simeq
\varprojlim_g \cA_{\le g}
\simeq
\widehat\Omega\bigl(\barB^{\mathrm{mod}}(\mathcal A)\bigr).
\]
\end{theorem}

\begin{proof}
The quotient
$\barB^{\mathrm{mod}}(\mathcal A)/F_{>g}\barB^{\mathrm{mod}}(\mathcal A)$
is a finite genus-window coalgebra by the compatibility hypotheses.
The quotient maps for $g+1\to g$ are morphisms of curved coalgebras, so
the completed cobar functor gives morphisms
$\cA_{\le g+1}\to\cA_{\le g}$. This constructs the inverse tower.

The equality with the completed cobar of the full modular bar coalgebra
is exactly the continuity of $\widehat\Omega$ with respect to the
chosen complete genus filtration. The Mittag--Leffler hypothesis removes
the possible $\varprojlim^1$ obstruction. Without that hypothesis the
finite-window algebras still exist, but their inverse limit need not
compute the full completed cobar object.
\end{proof}

\subsection{Spectral sequence convergence}

\begin{theorem}[Bar complex spectral sequence; \ClaimStatusProvedHere]
\label{thm:bar-complex-spectral-sequence}
\index{spectral sequence!bar-cobar|textbf}
There exists a spectral sequence:
\[E_2^{p,q} = H^p(\ConfigSpace{*}, H^q(\mathcal{A}^{\boxtimes *})) \Rightarrow H^{p+q}(\barBgeom(\mathcal{A}))\]
which converges under the following conditions:
\begin{enumerate}
\item $\mathcal{A}$ is bounded below: $\mathcal{A}_i = 0$ for $i < i_0$
\item The configuration spaces have finite cohomological dimension
\item The chiral algebra has finite homological dimension
\end{enumerate}
\end{theorem}

\begin{proof}
We filter the bar complex by \emph{configuration degree} (= bar degree):
\[F_p\barBgeom(\mathcal{A}) = \bigoplus_{n \leq p} \barBgeom^n(\mathcal{A}), \qquad p \ge 0,\]
and set $F_{-1} = 0$. The bar differential decomposes as
$d = d_{\mathrm{int}} + d_{\mathrm{fact}}$, where $d_{\mathrm{int}}$ preserves
configuration degree and $d_{\mathrm{fact}}$ raises it by one. Both respect the
filtration, so $d(F_p) \subset F_{p+1}$ and the filtration is compatible with the
differential.

\emph{Step 1 (associated graded and $E_1$ page).}
The successive quotients are
\[\mathrm{Gr}_p = F_p / F_{p-1} \cong \Omega^*(\ConfigSpace{p+1}) \otimes
\mathcal{A}^{\boxtimes(p+1)},\]
and the induced differential on $\mathrm{Gr}_p$ is $d_{\mathrm{int}}$ (the
internal differential on each factor). The $E_1$ page is therefore:
\[E_1^{p,q} = H^q(\mathrm{Gr}_p) = H^q\!\bigl(\Omega^*(\ConfigSpace{p+1})
\otimes \mathcal{A}^{\boxtimes(p+1)}\bigr).\]
The $d_1$ differential $d_1\colon E_1^{p,q} \to E_1^{p+1,q}$ is the map induced
by $d_{\mathrm{fact}}$ on the $E_1$ page. Taking cohomology yields the $E_2$ page:
\[E_2^{p,q} = H^p\!\bigl(\ConfigSpace{*},\, H^q(\mathcal{A}^{\boxtimes *})\bigr).\]

\emph{Step 2 (bounded-below and exhaustive).}
The filtration satisfies two structural properties:
\begin{enumerate}
\item \emph{Bounded below}: $F_{-1} = 0$, so configuration degree is $\ge 0$.
 In particular, $E_r^{p,q} = 0$ for $p < 0$ at every page $r \ge 0$.
\item \emph{Exhaustive}: Every element of $\barBgeom(\mathcal{A})$ has finite bar
 degree (it is a finite linear combination of terms, each lying in some
 $\barBgeom^n(\mathcal{A})$), so
 $\barBgeom(\mathcal{A}) = \bigcup_{p \ge 0} F_p\barBgeom(\mathcal{A})$.
\end{enumerate}
The bounded-below hypothesis $\mathcal{A}_i = 0$ for $i < i_0$ ensures that
$E_1^{p,q} = 0$ for $q < (p+1)\cdot i_0$, so the spectral sequence is concentrated
in a half-plane. For each fixed total degree $n = p + q$, only finitely many pairs
$(p,q)$ with $p \ge 0$ and $q \ge (p+1)\cdot i_0$ satisfy $p + q = n$; hence only
finitely many terms contribute in each total degree.

\emph{Step 3 (convergence).}
We verify the hypotheses of the classical convergence theorem for filtered complexes
(\cite{Weibel94}, Theorem~5.5.1; see also Boardman~\cite{Boardman-conditional},
Theorem~7.1):
\begin{enumerate}
\item[(a)] \emph{Bounded below}: $F_{-1} = 0$, i.e., the filtration starts at
 $p = 0$. Therefore $E_r^{p,q} = 0$ for all $p < 0$ and all $r \ge 0$.
\item[(b)] \emph{Exhaustive}: $\barBgeom(\mathcal{A}) = \bigcup_{p \ge 0}
 F_p\barBgeom(\mathcal{A})$, established in Step~2.
\end{enumerate}
The bounded-below condition implies that for each fixed bidegree $(p,q)$, the
incoming differentials $d_r\colon E_r^{p-r,\,q+r-1} \to E_r^{p,q}$ eventually have
source in the vanishing region $p - r < 0$, so $\ker(d_r)$ stabilizes at
$E_r^{p,q}$ for $r > p$. Similarly, the outgoing differentials $d_r\colon
E_r^{p,q} \to E_r^{p+r,\,q-r+1}$ stabilize because the half-plane concentration
from Step~2 forces the target to vanish for $r$ sufficiently large. Therefore
$E_\infty^{p,q} = E_{p+1}^{p,q}$ exists for every $(p,q)$, and the convergence
\[E_\infty^{p,q} \cong \mathrm{Gr}_p\, H^{p+q}(\barBgeom(\mathcal{A}))\]
holds: the $E_\infty$ page computes the associated graded of the induced filtration
on cohomology. The finite cohomological dimension of configuration spaces ensures
$E_2^{p,q} = 0$ for $p$ exceeding $\dim_{\mathrm{coh}}(\ConfigSpace{*})$ in each
total degree. This makes the filtration on $H^n(\barBgeom(\mathcal{A}))$ finite in
each degree, so the convergence is \emph{strong}: there is no $\varprojlim^1$
obstruction to recovering $H^n$ from $\mathrm{Gr}_\bullet\, H^n$, and the short
exact sequences $0 \to F_{p-1}H^n \to F_p H^n \to E_\infty^{p,n-p} \to 0$ split
(by induction on the finite filtration length).

\emph{Step 4 (Koszul degeneration).}
On the Koszul locus (Definition~\ref{def:chiral-koszul-pair}), the $E_2$ page is
concentrated: $H^q(\mathcal{A}^{\boxtimes(p+1)})$ is nonzero only for $q = p$
(diagonal concentration from bar cohomology purity). This forces all higher
differentials $d_r = 0$ for $r \ge 2$ by bidegree considerations, so the spectral
sequence degenerates at $E_2$ (Corollary~\ref{cor:spectral-degeneration}).
\end{proof}

\begin{corollary}[Degeneration {\cite{BGS96}}; \ClaimStatusProvedElsewhere]
\label{cor:spectral-degeneration}
\index{$E_2$ collapse}
If $\mathcal{A}$ is Koszul (Definition~\ref{def:chiral-koszul-pair}), the spectral sequence degenerates at $E_2$:
\[E_2^{p,q} = E_\infty^{p,q}\]
This gives:
\[H^n(\barBgeom(\mathcal{A})) = \bigoplus_{p+q=n} H^p(\ConfigSpace{*}) \otimes H^q(\mathcal{A}^!)\]
where $\mathcal{A}^!$ is the Koszul dual.
\end{corollary}

\subsection{Essential image of the bar functor}

\begin{theorem}[Essential image recognition with unit witness; \ClaimStatusProvedHere]
\label{thm:essential-image-bar}
Let $\mathcal{C}$ be a coaugmented conilpotent chiral coalgebra in the
regular holonomic finite-type ambient category. Then
$\mathcal{C}$ lies in the essential image of
\[
\barBgeom\colon
\mathrm{ChirAlg}^{\mathrm{aug}}_X\longrightarrow
\mathrm{Coalg}^{\mathrm{ch,conil}}_X
\]
up to filtered quasi-isomorphism if and only if there exist
\begin{enumerate}[label=\textup{(\roman*)}]
\item an augmented chiral algebra $\mathcal{A}$;
\item a chiral twisting morphism
 $\tau\colon\mathcal{C}\to\mathcal{A}$ satisfying
 $d\tau+\tau\smile\tau=0$;
\item logarithmic FM residue compatibility: the coderivation of
 $\mathcal{C}$ has conormal singularities only along the collision
 divisors $D_{ij}$, its residue on $D_{ij}$ is dual to the chiral
 product $\mu_{ij}$ of $\mathcal{A}$ under the finite-type pairing, and
 the codimension-two residues obey the Arnold--Orlik--Solomon boundary
 relation;
\item a filtered unit witness
 \[
 \eta_{\mathcal{C},\tau}\colon
 \mathcal{C}\longrightarrow \barBgeom(\mathcal{A})
 \]
 whose associated graded is a quasi-isomorphism and whose filtration is
 conilpotent or complete in the sense of
 Corollary~\ref{cor:bar-cobar-inverse}.
\end{enumerate}
The phrase ``essential image'' therefore includes the unit witness.
Logarithmic support and Arnold relations alone are necessary local
conditions; they are not sufficient without the twisting morphism and
the filtered quasi-isomorphism.
\end{theorem}

\begin{proof}
\emph{Necessity.}
If $\mathcal{C}\simeq\barBgeom(\mathcal{A})$, take $\tau$ to be the
universal twisting morphism
$\barBgeom(\mathcal{A})\to\mathcal{A}$ of
Proposition~\ref{prop:universal-twisting-adjunction}. The equation
$d\tau+\tau\smile\tau=0$ is precisely the chain-map condition for the
bar differential. The coderivation of $\barBgeom(\mathcal{A})$ is the
dual of the factorization part of the bar differential, hence is
supported on the FM collision divisors; the residue formula is the
definition of the chiral product, and the Arnold relation is the
codimension-two boundary identity used in the bar nilpotence proof. The
unit witness is the identity of $\barBgeom(\mathcal{A})$, transported
across the chosen filtered quasi-isomorphism
$\mathcal{C}\simeq\barBgeom(\mathcal{A})$.

\emph{Sufficiency.}
Given (i)--(iv), the map
$\eta_{\mathcal{C},\tau}$ is by hypothesis a filtered
quasi-isomorphism from $\mathcal{C}$ to the bar coalgebra of the
augmented chiral algebra $\mathcal{A}$. The logarithmic residue
compatibility says that this map is a morphism of dg coalgebras, not
only of graded coalgebras: it intertwines the FM boundary coderivation
with the bar differential determined by the product of $\mathcal{A}$. The
Maurer--Cartan equation for $\tau$ is exactly the remaining chain-map
condition under the twisting-morphism bijection
\[
\mathrm{Tw}^{\mathrm{ch}}(\mathcal{C},\mathcal{A})
\cong
\Hom_{\mathrm{dg\,CoAlg}^{\mathrm{ch,conil}}}
(\mathcal{C},\barBgeom(\mathcal{A}))
\]
of Corollary~\ref{cor:chiral-adjunction-hom-bijection}. Hence
$\mathcal{C}$ is isomorphic to an object in the bar image in the
filtered homotopy category. Conilpotent or complete convergence is
exactly the convergence input of Corollary~\ref{cor:bar-cobar-inverse},
so no infinite bar-degree sum is hidden in the recognition map.
\end{proof}

\begin{corollary}[Recognition principle with witness; \ClaimStatusConditional]
\label{cor:recognition-principle}
A chiral coalgebra $\mathcal{C}$ is recognized as a bar coalgebra once
$\Omega(\mathcal{C})$ is equipped with a strict chiral algebra structure
and the adjunction unit
\[
\mathcal{C}\longrightarrow \barBgeom(\Omega(\mathcal{C}))
\]
is a filtered quasi-isomorphism satisfying the residue compatibility of
Theorem~\ref{thm:essential-image-bar}. Strictness of
$\Omega(\mathcal{C})$ alone is not sufficient.
\end{corollary}
\begin{proof}
The bar-cobar adjunction (Theorem~\ref{thm:bar-cobar-adjunction})
provides unit and counit maps for any augmented chiral algebra and
any conilpotent coalgebra. On the Koszul locus
(Definition~\ref{def:chiral-koszul-pair}), the coalgebra-side unit
$\mathcal{C}\to\bar{B}(\Omega(\mathcal{C}))$ is a
quasi-isomorphism (Theorem~\ref{thm:bar-cobar-isomorphism-main}),
which identifies $\mathcal{C}$ as the bar of its cobar.
The theorem supplies the missing converse data: a strict cobar algebra
together with a filtered unit quasi-isomorphism and residue-compatible
twisting morphism.
\end{proof}

\subsection{BRST cohomology and string theory connection}

\begin{conjecture}[BRST/bar comparison datum; \ClaimStatusConjectured]\label{conj:brst-cohomology}
A BRST/bar comparison datum for a chosen matter--ghost system is a
chain map from its BRST complex to the geometric bar complex,
\[
\Phi_{\mathrm{BRST}} \colon
\bigl(C^\bullet_{\mathrm{BRST}}(\mathcal{A}), Q_{\mathrm{BRST}}\bigr)
\longrightarrow
\bigl(\barBgeom(\mathcal{A}), d_{\mathrm{bar}}\bigr),
\]
whose cohomological effect models the BRST/bar dictionary after
choosing the correct matter--ghost system and anomaly
prescription. On the anomaly-free genus-$0$ total
matter+ghost locus, this comparison is realized algebraically by
Theorem~\ref{thm:brst-bar-genus0}; beyond that setting it
requires the construction of the chain map, the anomaly cancellation
identity, and the sewing compatibility.

The expected correspondences are:
\begin{align}
Q_{\text{BRST}} &\leftrightarrow d_{\text{bar}} = d_{\text{int}} + d_{\text{fact}} + d_{\text{config}} \\
\text{Ghost number} &\leftrightarrow \text{Homological degree} \\
\text{BRST classes} &\leftrightarrow \text{Bar cohomology classes}
\end{align}
\end{conjecture}

\begin{remark}[String-field comparison data]
The proposed comparison has the following typed components.

\emph{String-field equation.} The string field $\Psi$ satisfies the BRST equation:
\[Q_{\text{BRST}} \Psi + \Psi \star \Psi = 0\]
where $\star$ is the string product.

\emph{Chiral-algebra expansion.} The string field decomposes as:
\[\Psi = \sum_{n=0}^\infty \Psi^{(n)} \otimes \omega^{(n)}\]
where $\Psi^{(n)} \in \mathcal{A}^{\otimes n}$ and $\omega^{(n)} \in \Omega^n(\overline{C}_n(X))$.

\emph{BRST action.} The BRST operator acts as:
\begin{align}
Q_{\text{BRST}}(\Psi^{(n)} \otimes \omega^{(n)}) &= \sum_{i=1}^n Q_i(\Psi^{(n)}) \otimes \omega^{(n)} \\
&\quad + \sum_{i<j} \mu_{ij}(\Psi^{(n)}) \otimes \text{Res}_{D_{ij}}[\omega^{(n)}] \\
&\quad + \Psi^{(n)} \otimes d_{\text{config}}\omega^{(n)}
\end{align}

The comparison map must identify this three-term BRST action with the
bar differential
$d = d_{\text{int}} + d_{\text{fact}} + d_{\text{config}}$.

\emph{Cohomology.} BRST classes are represented by
BRST-closed states modulo exact states. After the chain map and anomaly
hypotheses are supplied, the conjectural comparison predicts
\[
H^*_{\text{BRST}}
=
\text{Ker}(Q_{\text{BRST}})/\text{Im}(Q_{\text{BRST}})
\cong
H^*(\barBgeom(\mathcal{A})).
\]
\end{remark}

\begin{example}[Bosonic string theory]
For the bosonic string with central charge $c = 26$:

\emph{Ghost System.} The $(b,c)$ ghost system has OPE:
\[b(z)c(w) \sim \frac{1}{z-w}\]

\emph{BRST charge.}
\[Q_{\text{BRST}} = \oint dz \left[ c(z)T(z) + \frac{1}{2}:c(z)\partial c(z)b(z): \right]\]

\emph{Bar complex.} On the total anomaly-free matter--ghost system,
the conjectural comparison map has target the geometric bar complex
\[
\barBgeom(\mathrm{Vir}_{26}\otimes bc).
\]

\emph{Cohomology.} BRST state classes are cohomology classes of
$Q_{\text{BRST}}$. Identifying them with bar cohomology classes of weight
\((1,1)\) is the content of the comparison map, not a consequence of
bar-cobar inversion alone.
\end{example}

\begin{example}[Superstring theory]
For the superstring with central charge $c = 15$:

\emph{Superghost system.} The $(\beta,\gamma)$ system has OPE:
\[\beta(z)\gamma(w) \sim \frac{1}{z-w}\]

\emph{BRST charge.}
\[Q_{\text{BRST}} = \oint dz \left[ \gamma(z)G(z) + \frac{1}{2}:\gamma(z)\partial\gamma(z)\beta(z): \right]\]

\emph{Bar complex.} The comparison target includes the NS and R
sectors:
\[
\barBgeom(\mathcal{A}_{\text{NS}} \oplus \mathcal{A}_{\text{R}}).
\]

\emph{GSO Projection.} The GSO projection is an additional parity
projection on the BRST/string-field side. Its compatibility with the
bar complex is part of the comparison datum.
\end{example}

\begin{conjecture}[Anomaly cancellation for matter-ghost systems; \ClaimStatusConjectured]\label{conj:anomaly-cancellation}
For the \emph{coupled matter-ghost BRST system} in string theory, anomaly cancellation admits a geometric interpretation via the bar complex:

\begin{enumerate}
\item \emph{Central charge constraint.} With the standard ghost
normalizations $c_{bc}=-26$ and $c_{\beta\gamma}=-15$, the coupled
BRST differential is square-zero only on the total central-charge-zero
locus
\[
c_{\mathrm{total}}=c_{\mathrm{matter}}+c_{\mathrm{ghost}}=0,
\]
hence $c_{\mathrm{matter}}=26$ in the bosonic case and
$c_{\mathrm{matter}}=15$ in the superstring case.

\item \emph{Modular covariance.} The bar complex of the total system
transforms covariantly under $\operatorname{SL}_2(\mathbb{Z})$ only after
the anomaly class of the coupled matter--ghost system vanishes.

\item \emph{Geometric interpretation.} The anomaly is the obstruction
class for extending the total bar complex to higher genus with
compatible sewing.
\end{enumerate}

\end{conjecture}

\begin{remark}[Individual vs.\ total nilpotence]
This is \emph{not} in contradiction with the uncurved bar identity
$d_{\text{bar}}^2 = 0$ for an individual strict chiral algebra
(Theorem~\ref{thm:bar-nilpotency-complete}). In the curved central
lane the individual square is instead the inner curvature operator.
The BRST anomaly arises when one couples a matter chiral algebra to a
ghost system and requires the \emph{total} BRST charge to be
nilpotent; this imposes
$c_{\text{total}} = c_{\text{matter}} + c_{\text{ghost}} = 0$.
\end{remark}

\begin{remark}[Configuration-space anomaly class]
The anomaly arises from the failure of the \emph{total BRST
differential} to square to zero on the compactified configuration
space; it is not a statement about the uncurved individual bar
identity.

\emph{Individual nilpotence.} On an uncurved strict chiral
algebra, the bar differential satisfies $d_{\text{bar}}^2 = 0$ by
Theorem~\ref{thm:bar-nilpotency-complete}. In the curved central lane
one reads $d_{\text{bar},h}^{\,2}=[h_{\bar B},-]$ instead.

\emph{BRST coupling.} The total BRST differential
$d_{\text{BRST}} = d_{\text{matter}} + d_{\text{ghost}} +
d_{\text{coupling}}$ includes the matter--ghost interaction. Its square
is the inner action of the BRST anomaly class:
\[
d_{\text{BRST}}^2=[\mathfrak a_{\mathrm{BRST}},-],
\qquad
[\mathfrak a_{\mathrm{BRST}}]
=c_{\mathrm{total}}\,[\mathfrak a_{\mathrm{Vir}}].
\]
Here $[\mathfrak a_{\mathrm{Vir}}]$ is the Virasoro anomaly class in the
chosen normalization. The scalar coefficient is
$c_{\mathrm{total}}=c_{\mathrm{matter}}+c_{\mathrm{ghost}}$.

\emph{Cancellation.} The anomaly vanishes when
$c_{\text{matter}} + c_{\text{ghost}} = 0$. Since
$c_{\text{ghost}} = -26$ in the bosonic $bc$ system and
$c_{\text{ghost}} = -15$ in the superstring ghost system, the standard
critical values are $26$ and $15$ respectively.
\end{remark}

\begin{remark}
The geometric bar complex packages BRST cohomology, OPE residues on
configuration spaces, central-charge anomaly classes, and compatibility
with genus-one modular geometry. The comparison requires the chain map
and anomaly cancellation datum of Conjecture~\ref{conj:brst-cohomology}.
\end{remark}

\section{Relationship between bar-cobar and Koszul duality}

For a chiral Koszul pair $(\cA_1, \cA_2)$, the bar of one is the
Koszul dual coalgebra of the other, and Verdier duality on
$\operatorname{Ran}(X)$ provides the equivalence. The Koszul
property (diagonal Ext vanishing in the (bar degree, weight)
bigrading) is the condition under which this equivalence descends
from the coderived category to the ordinary derived category.

\subsection{Precise formulation of the relationship}

\begin{remark}[Criteria for chiral Koszul pairs]\label{rem:koszul-criteria}
\label{rem:koszul-chiral}
\index{Koszul duality!chiral}
\index{Koszul property}
Recall from Definition~\ref{def:chiral-koszul-pair} that a chiral Koszul pair
is specified by antecedent recognition data: a chiral twisting datum, a Koszul
twisting morphism, compatible filtrations, and Verdier-compatible dual coalgebras.
The bar/cobar quasi-isomorphisms are then proved consequences
(Theorem~\ref{thm:fundamental-twisting-morphisms} and
Theorem~\ref{thm:bar-cobar-isomorphism-main}), not part of the definition.
\end{remark}

\begin{remark}[Bar-cobar vs.\ Koszul: the fundamental distinction]\label{rem:fundamental-distinction}
\emph{Strict uncurved lane.}
\begin{itemize}
\item $\bar{B}: \mathcal{A} \to \bar{B}(\mathcal{A})$ exists and
 $d^2 = 0$ for uncurved augmented chiral algebras
 (Theorem~\ref{thm:bar-nilpotency-complete}).
\item $\Omega: \bar{B}(\mathcal{A}) \to \Omega(\bar{B}(\mathcal{A}))$ exists
 and $d^2 = 0$ for uncurved conilpotent bar coalgebras
 (Theorem~\ref{thm:cobar-d-squared-zero}).
\end{itemize}

These are constructions in the ordinary dg sense. In the curved
central lane the constructions still exist, but as CDG objects:
\[
d_{\bar B,h}^{\,2}=[h_{\bar B},-],
\qquad
d_{\Omega,h}^{\,2}=[h_\Omega,-].
\]
Strict nilpotence is restored only by a curvature-killing
Maurer--Cartan twist or by the flat comparison/coderived
totalization.

\emph{On the Koszul locus only
 \textup{(}chiral Koszul pairs,
 Definition~\textup{\ref{def:chiral-koszul-pair}):}}
\begin{itemize}
\item The counit $\Omega(\bar{B}(\mathcal{A})) \xrightarrow{\sim} \mathcal{A}$
 is a quasi-isomorphism (Theorem~\ref{thm:bar-cobar-isomorphism-main}).
\item $\bar{B}(\mathcal{A}_1) \simeq (\mathcal{A}_2^!)^\vee$ as chiral
 coalgebras (equivalently
 $\bar{B}(\mathcal{A}_1)^\vee \simeq \mathcal{A}_2^!$ as chiral algebras)
 (Theorem~\ref{thm:bar-concentration}). The Koszul dual algebra
 $\mathcal{A}_2^!$ and the Koszul dual coalgebra
 $\bar{B}(\mathcal{A}_1)$ live in different categories and are related
 by linear duality; conflating them violates the distinction of
 Remark~\ref{rem:cobar-three-functors}.
\end{itemize}

Off the Koszul locus, $\Omega(\bar{B}(\mathcal{A}))$ still exists but
the counit need not be a quasi-isomorphism; the bar-cobar object
represents $\mathcal{A}$ in the completed coderived category
$D^{\mathrm{co}}(\text{-}\mathrm{CoMod})$ of Positselski
(Theorem~\ref{thm:higher-genus-inversion}(b)).
\end{remark}

\begin{theorem}[Necessary conditions for chiral Koszul duality {\cite{FG12}}; \ClaimStatusProvedElsewhere]\label{thm:koszul-necessary}
For $(\mathcal{A}_1, \mathcal{A}_2)$ to form a chiral Koszul pair, the following must hold:
\begin{enumerate}
\item Both algebras are finitely generated over $\mathcal{D}_X$
\item The bar complexes have finite-dimensional cohomology in each degree
\item There exists a non-degenerate pairing $\langle -, - \rangle: \bar{B}(\mathcal{A}_1) \otimes \bar{B}(\mathcal{A}_2) \to \omega_X$
\end{enumerate}
\end{theorem}

\subsection{Diagram of relationships}

The relationship between bar, cobar, and Koszul duality can be summarized.
The diagram below respects the category distinction of
Remark~\ref{rem:cobar-three-functors}: the bar functor lands in
coalgebras, the Koszul dual lives in algebras, and linear duality
$(-)^\vee$ mediates between them.

\begin{center}
\begin{tikzcd}[column sep=large, row sep=huge]
\mathcal{A}_1 \arrow[r, "\bar{B}"]
 \arrow[d, shift left=2, "\text{Koszul}", "{\text{duality}}"' {description}]
 & \bar{B}(\mathcal{A}_1) \arrow[r, "(-)^\vee"]
 \arrow[dr, "\Omega"', bend left=10]
 & \bar{B}(\mathcal{A}_1)^\vee \arrow[r, "\simeq"]
 & \mathcal{A}_2^! \\
\mathcal{A}_2 \arrow[u, shift left=2]
 \arrow[r, "\bar{B}"']
 & \bar{B}(\mathcal{A}_2) \arrow[r, "(-)^\vee"']
 & \bar{B}(\mathcal{A}_2)^\vee \arrow[r, "\simeq"']
 \arrow[u, "\Omega"', bend right=10]
 & \mathcal{A}_1^!
\end{tikzcd}
\end{center}

\emph{Reading the diagram.}
\begin{itemize}
\item Horizontal arrow $\bar{B}$: always exists; sends chiral algebra to
 factorization coalgebra.
\item Horizontal arrow $(-)^\vee$: linear duality (coalgebra to algebra).
\item Horizontal equivalence $\simeq$: the Koszul condition identifies
 $\bar{B}(\mathcal{A}_i)^\vee$ with the Koszul dual algebra
 $\mathcal{A}_{3-i}^!$. The left and right columns live in different
 categories (algebras vs.\ coalgebras); conflating
 $\bar{B}(\mathcal{A}_1)$ with $\mathcal{A}_2^!$ violates the
 distinction of Remark~\ref{rem:cobar-three-functors}.
\item Vertical double arrow: Koszul duality as an operation on
 algebras (exists only for Koszul pairs).
\item Diagonal arrow $\Omega$: cobar reconstruction; the inversion
 $\Omega(\bar{B}(\mathcal{A}_i)) \simeq \mathcal{A}_i$
 (Corollary~\ref{cor:bar-cobar-inverse}). This is the round-trip,
 not a map to $\mathcal{A}_{3-i}^!$.
\end{itemize}

\subsection{Examples illustrating the distinction}

\begin{example}[Heisenberg: level shift required]\label{ex:heisenberg-koszul-vs-barcobar}
Bar-cobar inversion gives $\Omega(\bar{B}(\mathcal{H}_k)) \simeq \mathcal{H}_k$ automatically, while Koszul duality yields $\mathcal{H}_k^! \simeq \mathrm{Sym}^{\mathrm{ch}}(V^*)$ with curvature $m_0 = -k\,\omega$ (cf.\ \S\ref{sec:heisenberg-koszul}). These are distinct statements: bar-cobar inverts $\mathcal{H}_k$, but the Koszul dual is a different type of algebra (commutative vs.\ Lie). See Part~\ref{part:characteristic-datum}.
\end{example}

\section{Curved Koszul duality and quantum obstructions}
\label{sec:curved-koszul-quantum}

\begin{theorem}[Quantum deformation-obstruction complementarity; \ClaimStatusConditional]\label{thm:deformation-obstruction}
\textup{[Regime: curved-central on the Koszul locus; all genera
\textup{(}Convention~\textup{\ref{conv:regime-tags})}.]}

Let $\mathcal{A}$ be a chiral algebra on a curve $X$ lying in the
curved-central Koszul complementarity lane: the genus-$g$ completed bar
family is defined, the square-zero total comparison differential
$\Dg{g}$ or the coderived filtration is fixed, Verdier duality supplies
the center-transport map
$\mathcal Z(\mathcal A^!)\simeq\mathcal Z(\mathcal A)$, and the
obstruction-deformation pairing of Lemma~\ref{lem:obs-def-pairing} is
perfect. Then the central genus-$g$ correction spaces satisfy
\[
Q_g^Z(\mathcal{A}) \oplus Q_g^Z(\mathcal{A}^!)
\simeq
H^*(\mathcal{M}_g, \mathcal Z(\mathcal{A})).
\]
Here:
\begin{itemize}
\item $Q_g^Z(\mathcal{A})$ is the central summand of genus-$g$
 obstructions to Koszul complementarity;
\item $Q_g^Z(\mathcal{A}^!)$ is the central summand of genus-$g$
 deformations of the dual algebra;
\item $\mathcal Z(\mathcal{A})$ is the center local system of
 $\mathcal{A}$ on moduli;
\item the right side is cohomology of moduli with coefficients in that
 center local system.
\end{itemize}
\end{theorem}

\begin{proof}
Theorem~\ref{thm:quantum-complementarity-main} gives the same
decomposition from the Verdier involution~$\sigma$ and the genus
spectral sequence under these hypotheses. The spectral sequence is
taken for the square-zero total comparison differential, or equivalently
for the coderived filtration on the curved bar family; it is not the
ordinary spectral sequence of the raw fiber operator when
$\dfib^{\,2} \neq 0$.

\emph{Foundation: Curved $A_\infty$ Structures}

Following Gui--Li--Zeng~\cite{GLZ22}, a curved chiral algebra $\mathcal{A}$ has:
\begin{enumerate}
\item Multiplication: $\mu_2: \mathcal{A}^{\otimes 2} \to \mathcal{A}$
\item Higher operations: $\mu_n: \mathcal{A}^{\otimes n} \to \mathcal{A}$ for $n \geq 3$
\item Curvature: $\mu_0: \mathbb{C} \to \mathcal{A}$
\end{enumerate}
satisfying the curved $A_\infty$ relations:
\[\sum_{\substack{r+s+t=n \\ r,t \geq 0,\; s \geq 0}} (-1)^{rs+t}\, \mu_{r+1+t}(\mathrm{id}^{\otimes r} \otimes \mu_s \otimes \mathrm{id}^{\otimes t}) = 0\]

For $n=0$: $\mu_1(\mu_0) = 0$ (the curvature element is a $\mu_1$-cycle).

For $n=1$: $\mu_1^2 = [\mu_0, -]_{\mu_2}$, so $\mu_1$ is a differential only modulo curvature.

\emph{Step 1: Genus Stratification of Obstructions}

\noindent
\textbf{Notation.}
The differential $\dfib$ denotes the \emph{fiberwise} bar differential at
genus~$g$: the collision-residue differential on a fixed curve~$\Sigma_g$.
It is generically \emph{curved} ($\dfib^{\,2} \neq 0$).
The total corrected differential $\Dg{g}$, which incorporates period
corrections to restore strict nilpotence ($\Dg{g}^{\,2} = 0$), is
introduced in Convention~\ref{conv:higher-genus-differentials}.
Accordingly, any spectral-sequence or $H^*$ statement below uses either
$\Dg{g}$ or the coderived bar-degree filtration, not the raw curved
fiber complex $(\barB_g(\mathcal{A}), \dfib)$.

\medskip

The failure of the fiberwise bar differential to square to zero is measured by:
\[\dfib^{\,2}: \bar{B}^n_g(\mathcal{A}) \to \bar{B}^{n+2}_g(\mathcal{A})\]

\begin{lemma}[Obstruction operator {\cite{Positselski11}}; \ClaimStatusConditional]\label{lem:obstruction-class}
The composition $\dfib^{\,2}$ is the chain-level curvature operator
\[
\dfib^{\,2} = m_0^{(g)} \ast (-)
\]
with $m_0^{(g)} \in Z(\mathcal{A})$. It vanishes if and only if the
genus-$g$ fiberwise bar differential is square-zero on the nose.
\end{lemma}

\begin{proof}
In the CDG framework (Definition~\ref{def:chiral-CDG-coalgebra}), the genus-$g$
bar complex carries curvature $h = m_0^{(g)}$ satisfying $\dfib^{\,2} = h \ast (-)$.
Explicitly, $m_0^{(g)}$ is the genus-$g$ contribution to the bar curvature,
constructed as a fiber integral:
\[
m_0^{(g)} = \int_{\Sigma_g} \omega_g \in \bar{B}^0_g(\mathcal{A}),
\]
where $\omega_g$ is the volume form of the genus-$g$ propagator
(cf.\ Definition~\ref{def:genus-stratified-bar}).

\emph{Centrality.}
The element $m_0^{(g)}$ lies in the center $Z(\mathcal{A})$ by the
following chain-level argument. For any generator $a \in \mathcal{A}$,
the OPE $m_0^{(g)}{}_{(n)} a$ for $n \geq 0$ vanishes because $m_0^{(g)}$
has conformal weight~$0$ (it is a scalar on each conformal-weight space),
and the Borcherds identity forces the iterated residue
$[\mathrm{Res}_{D_{ij}}, \mathrm{Res}_{D_{kl}}]$ to act trivially on
generators when $D_{ij} \cap D_{kl} = \emptyset$ (factorization axiom),
while nested collisions contribute only scalar multiples (OPE associativity).

\emph{Cocycle condition.}
The curved $A_\infty$ relation at $n=0$
(Definition~\ref{def:curved-ainfty}) gives
$\mu_1(m_0^{(g)}) = 0$, so $m_0^{(g)}$ is closed. Hence
$\dfib^{\,2} = m_0^{(g)} \ast (-)$ is the chain-level obstruction to
lifting the genus-$0$ square-zero differential (where $m_0^{(0)} = 0$)
to genus~$g$. It vanishes exactly when the fiberwise genus-$g$
differential is already square-zero. Otherwise one must pass either to
the strict flat comparison differential $\Dg{g}$ of
Theorem~\ref{thm:quantum-diff-squares-zero} or to the coderived
framework; ordinary cohomology of the curved fiber complex
$(\bar{B}^{(g)}(\mathcal{A}), \dfib)$ is not the correct invariant.
\end{proof}

\emph{Step 2: Moduli space interpretation.}

The genus-$g$ corrections are parametrized by $H^*(\mathcal{M}_g, \mathcal{Z}(\mathcal{A}))$ (cohomology with coefficients in the center). The connection comes from period integrals:

\begin{lemma}[Period integral formula {\cite{Fay73}}; \ClaimStatusConditional]\label{lem:period-integral}
For $\omega \in \Omega^1(\overline{\mathcal{M}}_g)$, the genus-$g$ obstruction is:
\[\mathrm{Obs}_g(\omega) = \int_{\mathcal{C}_g} \omega \wedge
\left(\sum_{\text{cycles}} \mathrm{Res}_{D_i} \wedge \mathrm{Res}_{D_j}\right)\]
where the integral is over the universal curve $\mathcal{C}_g \to \overline{\mathcal{M}}_g$.
\end{lemma}

\begin{proof}
The fiberwise bar differential at genus~$g$ is
$\dfib = \sum_{D \subset \partial \overline{C}_n(\Sigma_g)} \mathrm{Res}_D$,
constructed from logarithmic propagators
$\eta_{ij} = d\log(z_i - z_j)$ on configuration spaces
(Definition~\ref{def:bar-differential-complete}). We compute $\dfib^{\,2}$ at the
cochain level.

\emph{Step~1: Double residue as fiber $2$-form.}
For $\phi \in \bar{B}^n_g(\mathcal{A})$, the composition $\dfib^{\,2}(\phi)$
is a sum of double residues $\mathrm{Res}_{D_i} \circ \mathrm{Res}_{D_j}$
over pairs of boundary divisors. When $D_i \cap D_j = \emptyset$
(disjoint collisions), the residues commute and cancel pairwise
(cf.\ the anticommutator $\{d_{\mathrm{fact}}, d_{\mathrm{config}}\} = 0$
of Theorem~\ref{thm:chiral-hochschild-differential}).
The surviving terms come from nested collisions $D_i \subset D_j$,
producing a $2$-form kernel $K_{ij}^{(g)}$ on~$\overline{C}_n(\Sigma_g)$.
By Lemma~\ref{lem:obstruction-class}, this kernel takes values
in~$Z(\mathcal{A})$.

\emph{Step~2: Fiber integration.}
Consider the universal family
$\pi\colon \overline{C}_n(\mathcal{C}_g) \to \overline{\mathcal{M}}_g$.
The pushforward $\pi_*(K_{ij}^{(g)})$ is a section of
$\mathcal{Z}(\mathcal{A})$ on $\overline{\mathcal{M}}_g$, the
fiber integral of the double-residue kernel. Evaluating against
$\omega \in \Omega^1(\overline{\mathcal{M}}_g)$ gives:
\[
\mathrm{Obs}_g(\omega)
= \int_{\overline{\mathcal{M}}_g} \omega \wedge \pi_*\!\left(
\sum_{i < j} K_{ij}^{(g)}\right)
= \int_{\mathcal{C}_g} \omega \wedge
\sum_{\text{cycles}} \mathrm{Res}_{D_i} \wedge \mathrm{Res}_{D_j},
\]
where the second equality is the projection formula.

\emph{Step~3: Boundary vanishing.}
The boundary contributions from $\partial\overline{C}_n(\Sigma_g)$
vanish fiberwise: on each fiber $\overline{C}_n(\Sigma_g)$, the
Arnold relations (Theorem~\ref{thm:arnold-three}) ensure
$\eta_{ij} \wedge \eta_{jk} + \eta_{jk} \wedge \eta_{ki}
+ \eta_{ki} \wedge \eta_{ij} = 0$, which forces the boundary
integral to cancel term by term.
The detailed stratum-by-stratum analysis across all boundary strata is carried out in
the proof of Theorem~\ref{thm:quantum-complementarity-main},
Steps~3--4.
\end{proof}

\emph{Step 3: Dual deformations}

Now consider the Koszul dual $\mathcal{A}^!$.  Its genus-$g$
deformation complex is not defined by the undualized cobar inverse of
$\barB(\mathcal A)$.  It is the chiral Hochschild deformation complex
of $\mathcal A^!$, computed on a chosen completed bar--cobar
resolution.  On the central summand this gives
\[
\mathrm{Def}^{Z}_g(\mathcal{A}^!)
 := H^1\bigl(\overline{\mathcal M}_g,
 \mathcal Z(\mathcal A^!)\bigr),
\]
with the full deformation complex obtained by replacing
\(\mathcal Z(\mathcal A^!)\) by
\(C^\bullet_{\mathrm{ch}}(\mathcal A^!,\mathcal A^!)\).

\begin{lemma}[Deformation space under center transport; \ClaimStatusProvedHere]\label{lem:deformation-space}
Assume the complementarity lane supplies a quasi-isomorphism of
central local systems
\[
\tau_Z\colon \mathcal Z(\mathcal A^!)\xrightarrow{\;\sim\;}
\mathcal Z(\mathcal A)
\]
compatible with the Verdier pairing on the chosen completed
bar--cobar resolutions. Then the central genus-$g$ deformation
summand of $\mathcal{A}^!$ maps naturally into cohomology with
coefficients in the transported center:
\[
\mathrm{Def}^{Z}_g(\mathcal{A}^!)
\hookrightarrow
H^*(\overline{\mathcal{M}}_g, \mathcal{Z}(\mathcal{A})).
\]
\end{lemma}

\begin{proof}
A genus-$g$ deformation of $\mathcal{A}^!$ is a flat family
$\tilde{\mathcal{A}}^! \to \overline{\mathcal{M}}_g$ with fiber at a
basepoint equal to $\mathcal{A}^!$. Infinitesimally, such deformations
are classified by
$H^1(\overline{\mathcal{M}}_g, \underline{\mathrm{End}}(\mathcal{A}^!))$,
where $\underline{\mathrm{End}}(\mathcal{A}^!)$ is the local system of
endomorphisms.

\emph{Chain-level identification.}
Choose the completed bar--cobar resolution of $\mathcal{A}^!$.
The endomorphism complex of this resolution is quasi-isomorphic to
the chiral Hochschild cochain complex, hence to the CDG Hom-complex
of Proposition~\ref{prop:cdg-hom-complex}:
\[
\underline{\mathrm{End}}(\mathcal{A}^!)
\;\simeq\;
\underline{\mathrm{Hom}}_C^{\mathrm{ch}}(\mathcal{A}^!, \mathcal{A}^!).
\]
The degree-$0$ cohomology of this complex is
$HH^0(\mathcal{A}^!) = Z(\mathcal{A}^!)$.  Bar-cobar inversion by
itself does not identify this center with $Z(\mathcal A)$; the
comparison is exactly the assumed center-transport map~$\tau_Z$.
Therefore:
\[
\mathrm{Def}^{Z}_g(\mathcal{A}^!)
= H^1(\overline{\mathcal{M}}_g, \mathcal{Z}(\mathcal{A}^!))
\hookrightarrow H^*(\overline{\mathcal{M}}_g, \mathcal{Z}(\mathcal{A})).
\]
\emph{(Note: One must use the local system $\mathcal{Z}(\mathcal{A})$,
not the trivial local system. By Harer's theorem,
$H^1(\mathcal{M}_g, \mathbb{Q}) = 0$ for $g \geq 2$, so the
deformation space would be trivially zero with trivial coefficients.)}
\end{proof}

\emph{Step 4: Perfect Pairing}

\begin{lemma}[Obstruction-deformation pairing; \ClaimStatusConditional]\label{lem:obs-def-pairing}
There is a perfect pairing:
\[\langle -, - \rangle: Q_g^Z(\mathcal{A}) \otimes Q_g^Z(\mathcal{A}^!) \to
H^{3g-3}(\overline{\mathcal{M}}_g, \mathbb{C})\]
given by the trace:
\[\langle \mathrm{Obs}, \mathrm{Def} \rangle = \mathrm{Tr}(\mathrm{Obs} \circ \mathrm{Def}).\]
\end{lemma}

\begin{proof}
\emph{Step~1: Cochain-level pairing.}
At the cochain level, the Verdier duality involution $\sigma$
on the center local system $\mathcal{Z}(\mathcal{A})$
(Theorem~\ref{thm:verdier-bar-cobar}) induces an involution
$\sigma^*$ on $C^*(\overline{\mathcal{M}}_g, \mathcal{Z}(\mathcal{A}))$.
Define the pairing on cochains by:
\[
\langle \alpha, \beta \rangle_{\sigma}
:= \int_{\overline{\mathcal{M}}_g}
\mathrm{Tr}_{\mathcal{Z}(\mathcal{A})}(\alpha \wedge \sigma(\beta)),
\]
where $\mathrm{Tr}_{\mathcal{Z}(\mathcal{A})}$ is the trace on the
center. This is well-defined because $\sigma$ is an isometry for the
Killing form on~$\mathcal{Z}(\mathcal{A})$ and the integral converges
on the compact orbifold $\overline{\mathcal{M}}_g$.

\emph{Step~2: Eigenspace decomposition.}
By Theorem~\ref{thm:quantum-complementarity-main}(ii), the Verdier
involution satisfies $\sigma^2 = \mathrm{id}$, so
$H^*(\overline{\mathcal{M}}_g, \mathcal{Z}(\mathcal{A}))$ decomposes
into $\pm 1$-eigenspaces:
$Q_g^Z(\mathcal{A}) = \ker(\sigma^* - \mathrm{id})$ and
$Q_g^Z(\mathcal{A}^!) = \ker(\sigma^* + \mathrm{id})$.
The pairing $\langle -, - \rangle_{\sigma}$ restricts to zero on
each eigenspace (since $\sigma(\beta) = \pm\beta$ and the
antisymmetry $\langle \alpha, \sigma\beta \rangle
= -\langle \sigma\alpha, \beta \rangle$ forces vanishing on
equal eigenvalues).

\emph{Step~3: Non-degeneracy.}
On the cross-terms $Q_g^Z(\mathcal{A}) \otimes Q_g^Z(\mathcal{A}^!)$,
the pairing reduces to Serre duality on $\overline{\mathcal{M}}_g$
with coefficients in $\mathcal{Z}(\mathcal{A})$. Serre duality is
non-degenerate on the compact orbifold $\overline{\mathcal{M}}_g$,
so the pairing is perfect.
\end{proof}

\emph{Step 5: Center Cohomology}

\begin{lemma}[Center as obstruction-deformation space; \ClaimStatusConditional]\label{lem:center-cohomology}
Assume the center-transport hypothesis of
Lemma~\ref{lem:deformation-space} and assume the transported
obstruction and deformation subspaces exhaust the $\pm1$ eigenspaces
of the Verdier involution on
\(H^*(\overline{\mathcal{M}}_g,\mathcal Z(\mathcal A))\). Then the
direct sum $Q_g^Z(\mathcal{A}) \oplus Q_g^Z(\mathcal{A}^!)$ identifies with:
\[H^*(\overline{\mathcal{M}}_g, \mathcal{Z}(\mathcal{A}))\]
\end{lemma}

\begin{proof}
The Verdier involution $\sigma$ on the center local system
$\mathcal{Z}(\mathcal{A})$ is an involution ($\sigma^2 = \mathrm{id}$)
that acts on cohomology. Since we work over~$\mathbb{C}$, the
eigenvalues of $\sigma^*$ on
$H^*(\overline{\mathcal{M}}_g, \mathcal{Z}(\mathcal{A}))$ are $\pm 1$,
giving a direct sum decomposition:
\[
H^*(\overline{\mathcal{M}}_g, \mathcal{Z}(\mathcal{A}))
= \ker(\sigma^* - \mathrm{id})
\oplus \ker(\sigma^* + \mathrm{id}).
\]
It remains to identify the two summands.

\emph{The $+1$-eigenspace.}
By Lemma~\ref{lem:obstruction-class}, the chain-level obstruction is
governed by the central coefficient $m_0^{(g)}$. Its flat-side
comparison class lies in the $\sigma$-invariant part of
$H^*(\overline{\mathcal{M}}_g, \mathcal{Z}(\mathcal{A}))$, because
$\sigma(m_0^{(g)}) = m_0^{(g)}$ (the curvature is Verdier-self-dual:
it is constructed from the volume form $\omega_g$, which is a canonical
section of $K_{\Sigma_g}$, and Verdier duality acts as the identity on
canonical sections). By the exhaustion hypothesis in the statement,
these obstruction classes fill the invariant summand, giving
$Q_g^Z(\mathcal{A}) = \ker(\sigma^* - \mathrm{id})$.

\emph{The $-1$-eigenspace.}
By Lemma~\ref{lem:deformation-space},
$\mathrm{Def}_g(\mathcal{A}^!)$ embeds in
$H^*(\overline{\mathcal{M}}_g, \mathcal{Z}(\mathcal{A}))$.
The anti-invariance of the transported deformation classes is not a
formal consequence of bar-cobar inversion; it is part of the
Verdier-compatible center-transport hypothesis. Under that hypothesis,
\[
Q_g^Z(\mathcal{A}^!) = \ker(\sigma^* + \mathrm{id}).
\]

\emph{Exhaustion.}
Since $\ker(\sigma^* - \mathrm{id}) \oplus \ker(\sigma^* + \mathrm{id})$
is the full space (over~$\mathbb{C}$, any involution is diagonalizable),
the decomposition is exhaustive.
\end{proof}

\emph{Step 6: Conclusion}

Combining Lemmas \ref{lem:obstruction-class}, \ref{lem:period-integral},
\ref{lem:deformation-space}, \ref{lem:obs-def-pairing}, and \ref{lem:center-cohomology},
we conclude:
\[
Q_g^Z(\mathcal{A}) \oplus Q_g^Z(\mathcal{A}^!)
\simeq
H^*(\mathcal{M}_g, \mathcal Z(\mathcal{A}))
\]
as claimed.

\end{proof}

\begin{remark}[Explicit formulas for low genus]\label{rem:explicit-low-genus-curved}

\emph{Genus 0.} $\mathcal{M}_0 = \text{pt}$, so $H^*(\mathcal{M}_0, Z(\mathcal{A}))
= Z(\mathcal{A})$. There are no quantum corrections.

\emph{Genus 1.} The moduli space $\mathcal{M}_1$ has $H^1(\mathcal{M}_1, \mathbb{Q}) = 0$ (as for all $\mathcal{M}_g$), but $\dim \mathcal{M}_1 = 1$. Quantum corrections enter through the \emph{partition function}, which depends on the modulus $\tau \in \mathfrak{h}/\mathrm{SL}_2(\mathbb{Z})$:
\[Q_1(\mathcal{H}_k) = \mathbb{C} \cdot k\]
where $k$ is the level of the Heisenberg algebra. Here $Q_1$ is computed from the center $Z(\mathcal{A})$ evaluated over $\mathcal{M}_1$, not from $H^1$.

The dual deformation:
\[Q_1(\text{Sym}(V^*)) = \mathbb{C} \cdot [V^* \wedge V^*]\]
measures how the symmetric algebra deforms from genus 0 to genus 1.

\emph{Genus 2.} $\dim \mathcal{M}_2 = 3$ and $H^2(\mathcal{M}_2, \mathbb{Q}) = \mathbb{Q} \cdot \lambda_1$ (generated by the Hodge class). For $\widehat{\mathfrak{sl}}_2$ at
critical level, the obstructions are:
\[Q_2(\widehat{\mathfrak{sl}}_2) = \mathbb{C} \cdot \lambda_1 \oplus \mathbb{C} \cdot
[\alpha, \alpha]\]
where $\lambda_1 \in H^2(\mathcal{M}_2)$ is the first Hodge class and $[\alpha, \alpha]$
is the self-commutator of the simple root.
\end{remark}

\begin{corollary}[Curved differential formula; \ClaimStatusConditional]\label{cor:curved-differential}
For a curved chiral algebra $\mathcal{A}$ with curvature $\mu_0$, the fiberwise genus-$g$ bar
differential is schematically of the form
\[\dfib = \dzero + \mu_0 \otimes \omega_g\]
where $\omega_g \in H^*(\mathcal{M}_g, \mathcal{Z}(\mathcal{A}))$ is the quantum correction class (living in the cohomology of $\mathcal{M}_g$ with coefficients in the center $\mathcal{Z}(\mathcal{A})$).

Its square is the chain-level obstruction operator
\[\dfib^{\,2} = [\mu_0, -] \otimes \omega_g^2,\]
whose cohomological shadow appears only after passage to the flat
comparison differential $\Dg{g}$ or the corresponding Chern--Weil
class.
\end{corollary}
\begin{proof}
The genus-$g$ bar differential receives quantum corrections from the
period matrix of the curve
(Theorem~\ref{thm:quantum-diff-squares-zero}). When $\mathcal{A}$ has
nonzero curvature $m_0 = \mu_0(1) \in \mathcal{A}^2$, the correction
class $\omega_g$ lives in
$H^*(\mathcal{M}_g, \mathcal{Z}(\mathcal{A}))$ by
Theorem~\ref{thm:quantum-complementarity-main}. Write
$\dfib = \dzero + m_0 \otimes \omega_g$, where the second term acts
via the inner derivation $[m_0, -]_{m_2}$ tensored with the modular
class. Squaring:
\[
\dfib^{\,2}
= \underbrace{\dzero^2}_{=\,0}
+ \dzero(m_0 \otimes \omega_g)
+ (m_0 \otimes \omega_g)\dzero
+ (m_0 \otimes \omega_g)^2.
\]
The cross terms vanish: $\dzero(m_0) = m_1(m_0) = 0$ by the curved
$A_\infty$ relation at $n = 0$
(Theorem~\ref{thm:curvature-central}(1)), and
$(m_0 \otimes \omega_g) \circ \dzero$ vanishes on cocycles by the
same centrality. The surviving term
$(m_0 \otimes \omega_g)^2 = [m_0, -]_{m_2} \otimes \omega_g^2$
uses the $n = 1$ curved $A_\infty$ relation $m_1^2 = [m_0, -]_{m_2}$
(Theorem~\ref{thm:curvature-central}(2)).
\end{proof}

\section{Synthesis}
\label{sec:cobar-supplement}

\begin{definition}[Finite Hall--Borcherds recognition gate;
\ClaimStatusConditional]
\label{def:finite-hall-borcherds-recognition-gate}
Fix a Mukai-weight cutoff $W$. A \emph{finite Hall--Borcherds
recognition gate} for $\mathbf H_{\Delta_5}$ up to weight $W$ is the
following finite datum.
\begin{enumerate}[label=\textup{(\roman*)}]
\item a positive-half CoHA model
$\mathcal H^+_{K3\times E,\le W}$ whose Hall generators are identified
with the positive BKM root spaces of $\mathfrak g_{\Delta_5}$ up to
weight $W$;
\item a completed Hall--Drinfeld double
$\widehat{\mathcal D}_\hbar(\mathcal H^+_{K3\times E})_{\le W}$ with
the chosen completion topology and the inclusion of the Cartan and
negative halves fixed;
\item a chain map
\[
\gamma_W\colon
\Omega_X\bigl(B_{\mathrm{ch}}(\mathbf H_{\Delta_5})\bigr)_{\le W}
\longrightarrow
\mathbf H_{\Delta_5,\le W}
\]
compatible with the Hall product, the bar/cobar differentials, and the
weight filtration;
\item source provenance fixing all normalisations: Schiffmann--Vasserot
and Davison for the positive-half CoHA, Borcherds for the automorphic
product, Gritsenko--Nikulin for the $\Delta_5$ root multiplicities, and
Drinfeld and Etingof--Kazhdan for the associator convention;
\item a parity/root fixture specifying the three real roots
$\alpha_{-1},\alpha_0,\alpha_1$, the imaginary-root parity, and the
positive cone used to pass from Fourier modes to BKM generators;
\item an independent coefficient match: for every root in the cutoff,
the Hall multiplicity, the Borcherds product exponent, and the
$\phi_{0,1}^{K3}$ Fourier coefficient give the same signed
multiplicity.
\end{enumerate}
A recognition system is a compatible family of such gates for all
cutoffs $W$, with transition maps commuting with the $\gamma_W$.
\end{definition}

\begin{remark}[Cobar construction on $B_{\mathrm{ch}}(\mathbf H_{\Delta_5})$;
\ClaimStatusConditional]
\label{rem:cobar-1}
Assume a recognition system in the sense of
Definition~\ref{def:finite-hall-borcherds-recognition-gate}. The cobar
construction $\Omega_X$, applied to the bar coalgebra of the K3 chiral
bialgebra $\mathbf H_{\Delta_5}$, then carries a comparison
\[
\Omega_X\bigl(B_{\mathrm{ch}}(\mathbf H_{\Delta_5})\bigr)
\;=\;
T_{\mathrm{ch}}\!\bigl(s^{-1}\overline{B_{\mathrm{ch}}(\mathbf H_{\Delta_5})}\bigr)
\;\xrightarrow{\;\gamma\;}\;
\mathbf H_{\Delta_5},
\]
where the tensor is taken in the Hall-theoretic ambient category
$\mathrm{CoHA}_{K3\times E}$ and the differential is the cobar-dual of
the bar differential $d_B$ decomposed in \S\ref{V2-sec:bar-cobar-review-supplement}.
If every finite map $\gamma_W$ is a quasi-isomorphism and the
quasi-isomorphisms commute with the transition maps, then $\gamma$
gives the K3 incarnation of bar--cobar inversion
\[
\Omega_X\bigl(B_{\mathrm{ch}}(\mathbf H_{\Delta_5})\bigr)
\;\simeq\;
\mathbf H_{\Delta_5}.
\]
At $\hbar^2 = -1/8$ the stronger Hall-subcategory equivalence is the
additional statement that the completed Hall--Drinfeld double in the
recognition gate identifies the source and target as Hall objects; away
from this Lusztig specialisation the output is only a filtered
quasi-isomorphism, with filtration controlled by the Siegel-modular
weight. Cross-ref Vol.~III,
Chapter~\ref{V3-ch:k3-chiral-bialgebra-platonic}, Theorem~A.
\end{remark}

\begin{remark}[Cobar differential four-piece decomposition;
\ClaimStatusConditional]
\label{rem:cobar-2}
Under the same finite recognition system, the cobar differential
$d_\Omega$ on
$\Omega_X(B_{\mathrm{ch}}(\mathbf H_{\Delta_5}))$ decomposes, by
duality to the bar differential, as
\[
d_\Omega \;=\; d^\vee_{\mathrm{Hall}} \;+\; d^\vee_{\mathrm{Sieg}} \;+\; \hbar\, d^\vee_{\mathrm{assoc}} \;+\; \hbar^2\, d^\vee_{\Phi_{10}/\eta^{24}},
\]
where the four pieces correspond, respectively, to Schiffmann--Vasserot
2012 CoHA comultiplication (dual of Hall multiplication), Borcherds 1995
singular-theta correspondence comultiplication (dual of Siegel
multiplication), the Drinfeld associator coboundary at order $\hbar$,
and the Gritsenko--Nikulin $\Phi_{10}/\eta^{24}$ twist at order $\hbar^2$.
On the raw curved object the square is the inner commutator with the
MC defect:
\[
d_{\Omega,\Theta}^{\,2}
=
\left[
D\cdot \Theta_{\mathbf H_{\Delta_5}}
+ \tfrac{1}{2}[\Theta_{\mathbf H_{\Delta_5}},\Theta_{\mathbf H_{\Delta_5}}],
-
\right].
\]
The closedness $d_{\Omega,\Theta}^{\,2} = 0$ is the assertion that this K3
Maurer--Cartan defect vanishes on the MC-twisted or coderived
totalization. Primary-lit: Schiffmann--Vasserot 2012,
Borcherds 1995, Drinfeld 1990, Gritsenko--Nikulin 1997--98.
\end{remark}

\begin{definition}[Cobar-side ordered convolution dGLA for
$\mathbf{H}_{\Delta_5}$; \ClaimStatusConditional]
\label{def:cobar-convolution-dgla-k3}
Conditionally on a recognition system as in
Definition~\ref{def:finite-hall-borcherds-recognition-gate}, and
dually to Definition~\ref{def:convolution-dgla-k3} of
Chapter~\ref{chap:bar-construction}, define the \emph{cobar-side}
ordered convolution dGLA
\[
 \mathfrak{C}^\Omega_{\mathbf{H}_{\Delta_5}}
 \;:=\;
 \mathrm{Hom}_{\Bbbk}\!\bigl(B_{\mathrm{ch}}(\mathbf{H}_{\Delta_5}),\,
 \mathbf{H}_{\Delta_5}\bigr)
 \;\cong_{\mathrm{Rec}}\;
 \mathrm{Hom}_{\mathrm{chAlg}}\!\bigl(
 \Omega_X\bigl(B_{\mathrm{ch}}(\mathbf{H}_{\Delta_5})\bigr),\,
 \mathbf{H}_{\Delta_5}\bigr),
\]
where $\cong_{\mathrm{Rec}}$ denotes the bar--cobar adjunction
(Theorem~\ref{thm:bar-cobar-inversion-qi}) applied to
$\mathbf{H}_{\Delta_5}$ after the finite chain maps $\gamma_W$ have
recognized the Hall model. Before this gate is supplied, the
unconditional object is only the first Hom complex. Explicitly: an
element
$\phi\in\mathfrak{C}^{\Omega,p}$ of cohomological degree $p$ is a
$\Bbbk$-linear map
$\phi\colon B_{\mathrm{ch}}(\mathbf{H}_{\Delta_5})\to
\mathbf{H}_{\Delta_5}$
shifting degree by $p$, working in the weight-completed pro-category
of Remark~\ref{rem:c1-weight-completion-k3-convolution} of the bar
chapter.
\begin{enumerate}[label=\textup{(\alph*)}]
\item \emph{Generators.} After the recognition gate has fixed the
positive-half CoHA model, the Drinfeld double completion, the
parity/root fixture, and the independent coefficient match, a basis for
$B_{\mathrm{ch}}(\mathbf{H}_{\Delta_5})$ in each finite weight window
is given by ordered tensor words
$w = s^{-1}x_{i_1}\otimes\cdots\otimes s^{-1}x_{i_n}$
with $x_{i_k}\in\{e_\alpha, f_\alpha, h_\alpha\}$ ranging over the
BKM generators of $\mathbf{H}_{\Delta_5}$: three real roots $\alpha_{-1},
\alpha_0, \alpha_1$ (the $\mathrm{II}_{1,1}$-hyperbolic triple of
Nikulin~1995), and imaginary roots indexed by the positive-cone
Fourier modes of the K3 elliptic genus $\phi_{0,1}^{K3}$ with
Gritsenko--Nikulin multiplicities $c_{\Delta_5}(N,\ell)$. A basis
element of the recognized finite piece
$\mathfrak{C}^{\Omega}_{\mathbf{H}_{\Delta_5},\le W}$ is specified by a
pair $(w, y)$ with $w\in B_{\mathrm{ch},\le W}$ and
$y\in\mathbf{H}_{\Delta_5,\le W}$: it maps $w\mapsto y$ and kills all
other recognized basis words. The completed dGLA is the inverse limit
over the finite recognized pieces.
\item \emph{Differential.} The cobar-side differential is
 \[
 (D\phi)(w) \;=\; d_A\bigl(\phi(w)\bigr)\;-\;(-1)^p \phi(d_B w),
 \]
 identical in form to the bar-side formula but evaluated on
 $B_{\mathrm{ch}} \supset B^{\mathrm{ord}}$. On the ordered subcomplex
 $B^{\mathrm{ord}}\hookrightarrow B_{\mathrm{ch}}$ the restriction
 recovers $\mathfrak{C}_{\mathbf{H}_{\Delta_5}}$; on the symmetrised
 quotient $B_{\mathrm{ch}}^{\Sigma}$ it factors through
 coinvariants.
\item \emph{Bracket.} The convolution cup
 $[\phi,\psi] = \phi\smile\psi - (-1)^{pq}\psi\smile\phi$,
 with $\smile = \mu\circ(\cdot\otimes\cdot)\circ\bar\Delta_{\mathrm{ch}}$
 using the deconcatenation coproduct of $B_{\mathrm{ch}}$.
\end{enumerate}
\end{definition}

\begin{proposition}[Bar-cobar naturality of MC;
\ClaimStatusConditional]
\label{prop:bar-cobar-MC-naturality-k3}
Under the adjunction $\Omega^{\mathrm{ch}}\dashv B^{\mathrm{ch}}$
\textup{(}Theorem~\ref{thm:bar-cobar-adjunction}\textup{)}, the
universal Maurer--Cartan element
$\Theta_{\mathbf{H}_{\Delta_5}}\in\mathfrak{C}^1_{\mathbf{H}_{\Delta_5}}[[\hbar]]$
of Remark~\ref{rem:theta-universal-MC-k3}
of the bar chapter corresponds under the universal
twisting bifunctor $\mathrm{Tw}^{\mathrm{ch}}(-,-)$
\textup{(}Proposition~\ref{prop:universal-twisting-adjunction}\textup{)} to
a unique twisting morphism
$\pi^{\Theta}\colon B_{\mathrm{ch}}(\mathbf{H}_{\Delta_5})\to
\mathbf{H}_{\Delta_5}$. Equivalently: the universal obstruction at
the bar level transports to a universal curving datum at the cobar level
via the Positselski / Getzler--Jones co/bar--MC dictionary
\textup{(}Positselski 2011 Chapter 6; Getzler--Jones 1990
\textup{Illinois J.}\ 34\textup{)}.
\end{proposition}

\begin{proof}
The adjunction is established in
Theorem~\ref{thm:bar-cobar-adjunction}; the universal twisting morphism
is constructed in Proposition~\ref{prop:universal-twisting-adjunction}.
The only point to verify is that the MC condition
$D\Theta + \tfrac{1}{2}[\Theta,\Theta]=0$ translates to the twisting
condition $d\pi + \pi\smile\pi = 0$; this is a standard rephrasing
(Loday--Vallette 2012 Chapter 11, Proposition 11.1.1), with the
sign $\tfrac{1}{2}$ vs.\ $\smile$ absorbed into the super-antisymmetry
of $[\cdot,\cdot]$ versus the non-symmetric $\smile$.
\end{proof}

\begin{remark}[Cobar side MC verification at
$\hbar^3, \hbar^4$; \ClaimStatusConditional]
\label{rem:cobar-MC-verification-k3}
The MC-twisted cobar-side closedness
$d_{\Omega,\Theta}^{\,2}=0$ on
$\Omega_X(B_{\mathrm{ch}}(\mathbf{H}_{\Delta_5}))$ translates, via
Proposition~\ref{prop:bar-cobar-MC-naturality-k3}, to the bar-side MC
equation verified in
Theorem~\ref{thm:MC-hbar3-hbar4-k3} of Chapter~\ref{chap:bar-construction}.
The $\hbar^3$ term is $\phi^{(3)} = \zeta(3)c_{\mathrm{symm}} +
(25/3)c_{\mathrm{timelike}} + (\Phi_{10}/\eta^{24})c_{\Phi_{10}}$
with $D\phi^{(3)}=0$ the pentagon equation
(Vol.~III, Theorem~\ref{V3-thm:pentagon-sieg-bor});
the $\hbar^4$ term is $\phi^{(4)}=0$ with
$[\phi^{(3)},\phi^{(3)}]_{\hbar^4}=0$ by $\hbar$-degree counting
($\hbar^3\cdot\hbar^3 = \hbar^6$); the next nontrivial MC
equation is the $\hbar^6$ super-Jacobi identity
$D\phi^{(6)} = -\tfrac{1}{2}[\phi^{(3)},\phi^{(3)}]$,
with $\phi^{(6)}$ reading off the six independent
MZV--Borcherds-leg cohomology representatives
$\zeta(3)^2 c_6^{(1)} + (25/3)\zeta(3) c_6^{(2)} + (25/3)^2 c_6^{(3)}
+ \zeta(3)(\Phi_{10}/\eta^{24})c_6^{(4)} + (\Phi_{10}/\eta^{24})^2 c_6^{(5)}
+ \zeta(5) c_6^{(6)}$, verified via MZV shuffle relations
(Brown 2012 \textup{Ann.~Math.}~175 \S 3) and Borcherds product
shuffle (Borcherds 1995 \textup{Invent.~Math.}~120 \S 8).
\end{remark}
